 
% new file names from date 20171119: 
% block1_171119.tex 
\renewcommand{\thisfile}{block1\_190328} 


\cleardoublepage
\mysection{Basic Facts on Categories} \label{sec:basic}
\AYcopyright

It is assumed that the reader has a working knowledge of categories and 
functors. References for this material are \cite{Mac1}, \cite{Mac2}, 
\cite{Rot}, \cite{HiSt}, \cite{We}, \cite{KaSc1} and \cite{KaSc2}. In this 
section we review the relevant material, and establish notation. 

\mysubsection{Set Theory} \label{subsec:set-theor}
In this book we will not try to be precise about issues of set theory. The
blanket assumption is that we are given a {\em Grothendieck universe}
$\cat{U}$. This is a ``large'' infinite set.
A {\em small set}, or a $\cat{U}$-small set, is a set $S$ that is an element of 
$\cat{U}$. We want all the products $\prod_{i \in I} S_i$ and disjoint unions 
$\coprod_{i \in I} S_i$, with $I$ and $S_i$ small sets, 
to be small sets too. 
(This requirement is not crucial for us, and it is more a matter of convenience.
When dealing with higher categories, one usually needs a hierarchy of 
universes anyhow.) We assume that the axiom of choice holds in $\cat{U}$. 

A {\em $\cat{U}$-category} is a category $\cat{C}$ whose
set of objects $\opn{Ob}(\cat{C})$ is a subset of $\cat{U}$, 
and for every $C, D \in \opn{Ob}(\cat{C})$ the set of morphisms 
$\opn{Hom}_{\cat{C}}(C, D)$ is small. 
If $\opn{Ob}(\cat{C})$ is also small, then $\cat{C}$ is called a {\em small 
category}. See \cite{SGA-4}, \cite{Mac2} or \cite[Section 1.1]{KaSc2}. Another 
approach, involving ``sets'' vs ``classes'',  can be found in \cite{Ne1}. 

We denote by $\cat{Set}$ the category of all small sets. So 
$\opn{Ob}(\cat{Set}) = \cat{U}$, and $\cat{Set}$ is a $\cat{U}$-category.
A group (or a ring, etc.) is called small if its underlying set is
small. We denote by $\cat{Grp}$, $\cat{Ab}$, $\cat{Rng}$ and 
$\cat{Rng_{\mrm{c}}}$ the categories of small groups, small abelian groups, 
small rings  and small commutative rings respectively.  
For a small ring $A$ we denote by $\cat{Mod} A$ the category of 
small left $A$-modules. 

By default we work with $\cat{U}$-categories, and from now on $\cat{U}$ will
remain implicit. There are several places in which we shall encounter 
set theoretical issues (regarding functor categories and localization of 
categories); but these problems can be solved by introducing a bigger universe
$\cat{V}$ such that $\cat{U} \in \cat{V}$.


\mysubsection{Notation and Conventions} \label{subsec:notation}

Let $\cat{C}$ be a category. 
We often write $C \in \cat{C}$ as an abbreviation for 
$C \in \opn{Ob}(\cat{C})$. For an object $C$, its identity automorphism is 
denoted by $\opn{id}_C$. The identity functor of $\cat{C}$ is denoted by 
$\opn{Id}_{\cat{C}}$. 

The opposite category of $\cat{C}$ is $\cat{C}^{\mrm{op}}$.
It has the same objects as $\cat{C}$, but the morphism sets are 
$\opn{Hom}_{\cat{C}^{\mrm{op}}}(C_0, C_1) :=
\opn{Hom}_{\cat{C}}(C_1, C_0)$,
and composition is reversed. Of course 
$(\cat{C}^{\mrm{op}})^{\mrm{op}} = \cat{C}$. 
The identity functor of $\cat{C}$ can be viewed as a contravariant functor 
$\opn{Op} : \cat{C} \to \cat{C}^{\mrm{op}}$. 
To be explicit, on objects we take 
$\opn{Op}(C) := C$. As for morphisms, given a morphism 
$\phi : C_0 \to C_1$ in $\cat{C}$, we let 
$\opn{Op}(\phi) : \opn{Op}(C_1) \to \opn{Op}(C_0)$
be the morphism $\opn{Op}(\phi) := \phi$ in $\cat{C}^{\mrm{op}}$.
The inverse functor $\cat{C}^{\mrm{op}} \to \cat{C}$ is also denoted by 
$\opn{Op}$.
Thus $\opn{Op} \circ \opn{Op} = \opn{Id}$. 

A contravariant functor $F : \cat{C} \to \cat{D}$
is the same as a covariant functor 
$F \circ \opn{Op} : \cat{C}^{\mrm{op}} \to \cat{D}$.
By default all functors will be covariant, unless explicitly mentioned 
otherwise. Contravariant functors will almost always we dealt with by replacing 
the source 
category with its opposite.

\begin{dfn} \label{dfn:4125}
Let $\K$ be a commutative ring. By a 
{\em central $\K$-ring}%
\index{Central $\K$-ring} 
we mean a ring 
$A$, together with a ring homomorphism $\K \to A$, called the structural 
homomorphism, such that the image of 
$\K$ is inside the center of $A$. 

The category of central $\K$-rings, whose morphisms are the ring homomorphisms
$f : A \to B$ that respect the structural homomorphisms from $\K$, is denoted 
by $\catt{Rng} \centover \K$.  
\end{dfn}

Traditionally, a central $\K$-ring was called a ``unital associative 
$\K$-algebra''. Of course all rings and ring homomorphisms are unital.
When $\K = \Z$, a central $\K$-ring is just a ring, and then we 
sometimes use the notation $\catt{Rng}$. 

\begin{exa}  \label{exa:1035}
Let $\K$ be a nonzero commutative ring, and let $n$ be a positive integer. Then 
$\opn{Mat}_n(\K)$, the ring of $n \times n$ matrices with entries in $\K$, is a 
central $\K$-ring. 
\end{exa}

\begin{dfn} \label{dfn:4126}
Let $A$ be a ring. 
We denote by $\cat{Mod} A$, or by the abbreviated notation $\dcat{M}(A)$, the 
category of left $A$-modules.
\end{dfn}

Rings and modules are very important for us, so let us also put forth the next 
convention.

\begin{conv} \label{conv:2490} %\mbox{}
Below are the default implicit assumptions for linear structures and 
operations. 
\begin{enumerate}
\item There is a nonzero commutative base ring $\K$ (e.g.\ the ring of integers 
$\Z$ of a field). 

\item The unadorned tensor symbol $\ot$ means $\ot_{\K}$. 

\item All rings are central $\K$-rings (see Definition 
\ref{dfn:4125}), all ring homomorphisms are over $\K$, and all bimodules are 
$\K$-central. 

\item Generalizing (3), all linear categories are $\K$-linear (see 
Definition \ref{dfn:1083}), and all linear functors are $\K$-linear (see 
Definition \ref{dfn:1031}).

\item For a ring $A$, by all $A$-modules are left $A$-modules, unless 
explicitly stated otherwise. 
\end{enumerate}
\end{conv}

Right $A$-modules are left modules over the opposite ring $A^{\mrm{op}}$, and 
this is the way we shall most often deal with them. 
Morphisms in the categories of groups, rings, $A$-modules, etc.\ will 
usually be called group homomorphism, ring homomorphisms, $A$-module 
homomorphisms, etc., respectively. 

\begin{conv} \label{conv:3170}
We will try to keep the following font and letter conventions: 
\begin{itemize}
\item $f : C \to D$ is a morphism between objects in a category. 
\item $F : \cat{C} \to \cat{D}$ is a functor between categories. 
\item $\eta : F \to G$ is morphism of functors (i.e.\ a natural transformation) 
between functors $F, G : \cat{C} \to \cat{D}$.
\item $\phi, \psi, \phi_i : M \to N$ are morphisms between objects in an 
abelian category $\cat{M}$. 
\item $F : \cat{M} \to \cat{N}$ is a linear functor between abelian 
categories.
\item The category of complexes in an abelian category $\cat{M}$ is 
$\dcat{C}(\cat{M})$.
\item If $\cat{M}$ is a module category, and $M \in \opn{Ob}(\cat{M})$,
then elements of $M$ will be denoted by $m, n, m_i, \ldots$.
\end{itemize}
\end{conv}




\mysubsection{Epimorphisms and Monomorphisms} \label{subsec:zero}
Let $\cat{C}$ be a category. Recall that a morphism $f : C \to D$ in $\cat{C}$ 
is called an {\em isomorphism} if there is a morphism 
$g : D \to C$ such that $f \circ g = \opn{id}_D$ and 
$g \circ f = \opn{id}_C$. The morphism $g$ is called the {\em inverse} of $f$, 
it is unique (if it exists), and it is denoted by $f^{-1}$. 
An isomorphism is often denoted by this shape of arrow: 
$f : C \iso D$. 

A morphism $f : C \to D$ in $\cat{C}$ is called an
{\em epimorphism} if it has the right cancellation property: for every
$g, g' : D \to E$, $g \circ f = g' \circ f$ implies $g = g'$. 
An epimorphism is often denoted by this shape of arrow: 
$f : C \surj D$. 

A morphism $f : C \to D$ is called a {\em monomorphism} if it has the left
cancellation property: for every
$g, g' : E \to C$, $f \circ g = f \circ g'$ implies $g = g'$. 
A monomorphism is often denoted by this shape of arrow: 
$f : C \inj D$. 

\begin{exa}
In $\cat{Set}$ the monomorphisms are the injections, and the epimorphisms are
the surjections. A morphism $f : C \to D$ in $\cat{Set}$ that is both a
monomorphism and an epimorphism is an isomorphism. The same holds in the 
category $\cat{Mod} A$ of left modules over a ring $A$. 
\end{exa}

This example could be misleading, because the property of being an epimorphism 
is often not preserved by forgetful functors, as the next exercise shows. 

\begin{exer} \label{exer:1310}
Consider the category of rings $\cat{Rng}$.
Show that the forgetful functor 
$\cat{Rng} \to \cat{Set}$ respects monomorphisms, but it does not respect
epimorphisms. 
(Hint: show that the inclusion $\Z \to \mbb{Q}$ is an epimorphism in 
$\cat{Rng}$.)  
\end{exer}

By a {\em subobject} of an object $C \in \cat{C}$ we mean a monomorphism 
$f : C' \inj C$ in $\cat{C}$. We sometimes write $C' \subseteq C$ in this 
situation, but this is only notational (and does not mean inclusion of sets). 
We say that two subobjects $f_0 : C'_0 \inj C$ and $f_1 : C'_1 \inj C$ of $C$
are {\em isomorphic} if  there is an isomorphism $g : C'_0 \iso C'_1$ such that 
$f_1 \circ g = f_0$.

Likewise, by a {\em quotient} of $C$ we mean an 
epimorphism $g : C \surj C''$ in $\cat{C}$. There is an analogous notion of 
isomorphic quotients. 

\begin{exer} \label{exer:1676}
Let $\cat{C}$ be a category, and let $C$ be an object of $\cat{C}$. 
\begin{enumerate}
\item Suppose $f_0 : C'_0 \inj C$ and $f_1 : C'_1 \inj C$ are subobjects of
$C$. Show that there is at most one morphism $g : C'_0 \to C'_1$ such that 
$f_1 \circ g = f_0$; and if $g$ exists, then it is a monomorphism. 

\item Show that isomorphism is an equivalence relation on the set of subobjects 
of $C$. Show that the set of equivalence classes of subobjects of $C$ is 
partially ordered by ``inclusion''. (Ignore set-theoretical issues.) 

\item Formulate and prove the analogous statements for quotient objects. 
\end{enumerate}
\end{exer}

An {\em initial object} in a category $\cat{C}$ is an object 
$C_0 \in \cat{C}$,
such that for every object $C \in \cat{C}$ there is exactly one morphism 
$C_0 \to C$. Thus the set $\opn{Hom}_{\cat{C}}(C_0, C)$ is a singleton.
A {\em terminal object} in $\cat{C}$ is an object $C_{\infty} \in \cat{C}$,
such that for every object $C \in \cat{C}$ there is exactly one morphism 
$C \to C_{\infty}$.

\begin{dfn}
A {\em zero object} in a category $\cat{C}$ is an object which is both initial
and terminal. 
\end{dfn}

Initial, terminal and zero objects are unique up to unique isomorphisms (but
they need not exist).

\begin{exa}
In $\cat{Set}$, $\varnothing$ is an initial object, and every singleton is a
terminal object. There is no zero object.
\end{exa}

\begin{exa}
In $\cat{Mod} A$, every trivial module (with only the zero element) is a 
zero object, and we denote this module by $0$. This is allowed, since any other
zero module is uniquely isomorphic to it.  
\end{exa}


\mysubsection{Products and Coproducts}
Let $\cat{C}$ be a category. By a {\em collection of objects of $\cat{C}$} 
indexed by a (small) set $I$, we mean a function $I \to \opn{Ob}(\cat{C})$,
$i \mapsto C_i$. We usually denote this collection using the curly brackets 
notation: $\{ C_i \}_{i \in I}$.

Given a collection $\{ C_i \}_{i \in I}$ of objects of $\cat{C}$, 
its {\em product} is a pair $(C, \{ p_i \}_{i \in I})$
consisting of an object $C \in \cat{C}$, and a collection
$\{ p_i \}_{i \in I}$ of morphisms $p_i : C \to C_i$, called 
{\em projections}. The pair $(C, \{ p_i \}_{i \in I})$ 
must have this universal property: given an object $D$ and morphisms 
$f_i : D \to C_i$, there is a unique morphism $f : D \to C$  s.t.\ 
$f_i = p_i \circ f$. Of course if a product
$(C, \{ p_i \}_{i \in I})$ exists, then it is unique up to a
unique isomorphism; and we usually write
$\prod_{i \in I} C_i := C$, leaving the projection morphisms implicit. 

\begin{exa}
In $\cat{Set}$ and $\cat{Mod} A$ all products exist,
and they are the usual cartesian products. 
\end{exa}

For a collection $\{ C_i \}_{i \in I}$ of objects
of $\cat{C}$, its {\em coproduct} is a pair $(C, \{ e_i \}_{i \in I})$,
consisting of an object $C$ and a collection $\{ e_i \}_{i \in I}$
of morphisms $e_i : C_i \to C$, called {\em embeddings}. 
The pair $(C, \{ e_i \}_{i \in I})$ 
must have this universal property: given an object $D$ and morphisms 
$f_i : C_i \to D$, there is a unique morphism $f : C \to D$  s.t.\ 
$f_i = f \circ e_i$. If a coproduct
$(C, \{ e_i \}_{i \in I})$ exists, then it is unique up to a
unique isomorphism; and we write
$\coprod_{i \in I} C_i := C$,  leaving the embeddings implicit. 

\begin{exa} \label{exa:103}
In $\cat{Set}$ the coproduct is the disjoint union. 
In $\cat{Mod} A$ the coproduct is the direct sum. 
\end{exa}

Product and coproducts are very degenerate cases of {\em limits} and {\em 
colimits} respectively. In this book we will not need to use limits and 
colimits in their most general form. All we shall need is inverse limits and 
direct limits indexed by $\N$; and these will be recalled in Subsection 
\ref{subsec:inv-dir-lim} below. 

We do need to talk about {\em fibered products}. Let $\cat{C}$ be some 
category. Recall that a commutative diagram 
\begin{equation} \label{eqn:3625}
\UseTips \xymatrix @C=6ex @R=6ex {
E
\ar[d]_{g_1}
\ar[r]^{g_2}
&
D_2
\ar[d]^{f_2}
\\
D_1
\ar[r]^{f_1}
&
C
}
\end{equation}
in $\cat{C}$ is called {\em cartesian} if for every object $E' \in \cat{C}$, 
with morphisms 
$g_1' : E' \to D_1$ and $g_2' : E' \to D_2$ that satisfy 
$f_1 \circ g'_1 = f_2 \circ g'_2$, there exists a unique morphism 
$h : E' \to E$ such that 
$g'_i = g_i \circ h$.
See commutative diagram below. 
\[ \UseTips \xymatrix @C=6ex @R=6ex {
E'
\ar@{-->}@(d,ul)[ddr]_{g'_1}
\ar@{-->}@(r,ul)[drr]^{g'_2}
\ar@{-->}[dr]_{h}
\\
&
E
\ar[d]_{g_1}
\ar[r]^{g_2}
&
D_2
\ar[d]^{f_2}
\\
&
D_1
\ar[r]^{f_1}
&
C
} \]

A cartesian diagram is also called a {\em pullback diagram}, and the object $E$ 
is called the {\em fibered product} of $D_1$ and $D_2$ over $C$, with notation
$D_1 \times_C D_2 := E$.
This notation leaves the morphisms implicit. Of course if a fibered product 
exists, than it is unique up to a unique isomorphism that commutes with the 
given arrows. 

There is a dual notion: {\em fibered coproduct}. The input is morphisms 
$C \to D_1$ and $C \to D_2$ in $\cat{C}$, and the fibered coproduct 
$D_1 \sqcup_{C} D_1$ in $\cat{C}$ is just the fibered product in 
the opposite category $\cat{C}^{\mrm{op}}$.



\mysubsection{Equivalence of Categories}
Recall that a functor $F : \cat{C} \to \cat{D}$ is an {\em equivalence} if
there exist a functor $G : \cat{D} \to \cat{C}$, and isomorphisms of 
functors (i.e.\ natural isomorphisms)
$G \circ F \iso \opn{Id}_{\cat{C}}$ and 
$F \circ G \iso  \opn{Id}_{\cat{D}}$. Such a functor $G$ is called a {\em
quasi-inverse} of $F$, and it is unique up to isomorphism (if it exists).  

The functor $F : \cat{C} \to \cat{D}$ is {\em full} (resp.\  {\em faithful})
if for every $C_0, C_1 \in \cat{C}$ the function 
\[ F : \opn{Hom}_{\cat{C}}(C_0, C_1) \to 
\opn{Hom}_{\cat{D}} \bigl( F(C_0), F(C_1) \bigr) \]
is surjective (resp.\ injective). 

We know that $F : \cat{C} \to \cat{D}$ is an equivalence if and only if these 
two conditions hold:
\begin{enumerate}
\rmitem{i} $F$ is essentially surjective on objects. This means that for every
$D \in \cat{D}$  there is an isomorphism $F(C) \iso D$ for some $C \in \cat{C}$.

\rmitem{ii} $F$ is fully faithful (i.e.\ full and faithful). 
\end{enumerate}

\begin{exer}
If you are not sure about the last claim (characterization of equivalences), 
then prove it. (Hint: use the axiom of choice to construct a quasi-inverse of 
$F$.)  
\end{exer}

A functor $F : \cat{C} \to \cat{D}$ is called an {\em isomorphism of 
categories} if it is bijective on sets of objects and on sets of morphisms. It 
is clear that an isomorphism of categories is an equivalence. If $F$ is an 
isomorphism of categories, then it has an inverse isomorphism 
$F^{-1} : \cat{D} \to \cat{C}$, which is unique. 



\mysubsection{Bifunctors} \label{subsec:bifunc}
Let $\cat{C}$ and $\cat{D}$ be categories. Their product is the category 
$\cat{C} \times \cat{D}$ defined as follows: the set of objects is 
\[ \opn{Ob}(\cat{C} \times \cat{D}) := 
\opn{Ob}(\cat{C}) \times \opn{Ob}(\cat{D}) . \] 
The sets of morphisms are 
\[ \opn{Hom}_{\cat{C} \times \cat{D}} \big( (C_0, D_0), (C_1, D_1) \big) :=
\opn{Hom}_{\cat{C} } (C_0, C_1) \times 
\opn{Hom}_{\cat{D} } (D_0, D_1)  . \]
The composition is 
\[ (f_1, g_1) \circ (f_0, g_0) := 
(f_1 \circ f_0, g_1 \circ g_0) , \]
and the identity morphism on an object $(C, D)$ is
$( \opn{id}_C,  \opn{id}_D)$. 

A {\em bifunctor} from $(\cat{C}, \cat{D})$ to $\cat{E}$ is a functor 
$F : \cat{C} \times \cat{D} \to \cat{E}$.
The extra information that is implicit when we call $F$ a bifunctor 
is that the source category $\cat{C} \times \cat{D}$ is a product. 


\mysubsection{Representable Functors} \label{subsec:repfunc}

Let $\cat{C}$ be a category. An object $C \in \cat{C}$ gives rise 
to a functor 
\begin{equation} \label{eqn:3025}
\opn{Y}_{\cat{C}}(C) : \cat{C}^{\mrm{op}} \to \cat{Set} , \quad
\opn{Y}_{\cat{C}}(C) := \opn{Hom}_{\cat{C}}(-, C) .  
\end{equation}
Explicitly, the functor $\opn{Y}_{\cat{C}}(C)$ sends an object 
$D \in \cat{C}$ to the set 
$\opn{Y}_{\cat{C}}(C)(D) :=  \opn{Hom}_{\cat{C}}(D, C)$,
and a morphism
$\psi : D_0 \to D_1$ in $\cat{C}$ goes to the function 
\[ \opn{Y}_{\cat{C}}(C)(\psi) := \opn{Hom}_{\cat{C}}(\psi,  \opn{id}_C) : 
\opn{Hom}_{\cat{C}}(D_1, C) \to \opn{Hom}_{\cat{C}}(D_0, C) . \]

Now suppose we are given a morphism 
$\phi : C_0 \to C_1$ in $\cat{C}$. There is a morphism of functors (a natural 
transformation) 
\begin{equation} \label{eqn:3026}
\opn{Y}_{\cat{C}}(\phi) := \opn{Hom}_{\cat{C}}(-, \phi) : 
\opn{Y}_{\cat{C}}(C_0) \to \opn{Y}_{\cat{C}}(C_1) .
\end{equation}

Consider the category 
$\cat{Fun}(\cat{C}^{\mrm{op}}, \cat{Set})$,
whose objects are the functors $F : \cat{C}^{\mrm{op}} \to \cat{Set}$, and
whose morphisms are the morphisms of functors
$\eta : F_0 \to F_1$.  
There is a set-theoretic difficulty here: the sets of objects and morphisms 
of $\cat{Fun}(\cat{C}^{\mrm{op}}, \cat{Set})$
are too big (unless $\cat{C}$ is a small category), and this is not a 
$\cat{U}$-category. Hence we must enlarge the universe, as mentioned in 
Subsection \ref{subsec:set-theor}. 

\begin{dfn} \label{dfn:3025}
The {\em Yoneda functor} of the category $\cat{C}$ is the functor
\[ \opn{Y}_{\cat{C}} : \cat{C} \to \cat{Fun}(\cat{C}^{\mrm{op}}, \cat{Set}) \]
described by formulas (\ref{eqn:3025}) and (\ref{eqn:3026}).
\end{dfn}

\begin{thm}[Yoneda Lemma] \label{thm:3025}
The Yoneda functor $\opn{Y}_{\cat{C}}$ is fully faithful. 
\end{thm}

See \cite[Section III.2]{Mac2} or \cite[Section 1.4]{KaSc2} for a proof. The 
proof is not hard, but it is very confusing. 

A functor $F : \cat{C}^{\mrm{op}} \to \cat{Set}$ is called 
{\em representable} if there is an isomorphism of functors 
$\eta : F \iso \opn{Y}_{\cat{C}}(C)$ for some object $C \in \cat{C}$.
Such an object $C$ is said to {\em represent} the functor $F$. 
The Yoneda Lemma says that $\opn{Y}_{\cat{C}}$ is an equivalence
from $\cat{C}$ to the category of representable functors. 
Thus the pair $(C, \eta)$ is unique up to a unique isomorphism (if it exists).
Note that the isomorphism of sets
$\eta_C : F(C) \iso \opn{Y}_{\cat{C}}(C)(C)$ 
gives a special element $\til{\eta} \in F(C)$ such that 
$\eta_C(\til{\eta}) = \opn{id}_C$. 

Dually, any object $C \in \cat{C}$ gives rise to a functor
\begin{equation} \label{eqn:3027}
\opn{Y}_{\cat{C}}^{\vee}(C) : \cat{C} \to \cat{Set} , \quad
\opn{Y}_{\cat{C}}^{\vee}(C) := \opn{Hom}_{\cat{C}}(C, -) . 
\end{equation}
A morphism $\phi : C_0 \to C_1$ in $\cat{C}$ induces a morphism of 
functors
\begin{equation} \label{eqn:3028}
\opn{Y}_{\cat{C}}^{\vee}(\phi) := \opn{Hom}_{\cat{C}}(\phi, -) : 
\opn{Y}_{\cat{C}}^{\vee}(C_1) \to \opn{Y}_{\cat{C}}^{\vee}(C_0) .
\end{equation}

\begin{dfn} \label{dfn:3026}
The {\em dual Yoneda functor} of the category $\cat{C}$ is the functor
\[ \opn{Y}_{\cat{C}}^{\vee} : \cat{C}^{\mrm{op}} \to 
\cat{Fun}(\cat{C}, \cat{Set}) \]
described by formulas (\ref{eqn:3027}) and (\ref{eqn:3028}).
\end{dfn}

\begin{thm}[Dual Yoneda Lemma] \label{thm:3026}
The dual Yoneda functor $\opn{Y}_{\cat{C}}^{\vee}$ is fully faithful. 
\end{thm}

This is also proved in \cite[Section III.2]{Mac2} and 
\cite[Section 1.4]{KaSc2}. 

A functor $F : \cat{C} \to \cat{Set}$ is called {\em corepresentable} if there 
is an isomorphism of functors
$\eta : F \iso \opn{Y}_{\cat{C}}^{\vee}(C)$ for some object $C \in \cat{C}$.
The object $C$ is said to {\em corepresent} the functor $F$.
The dual Yoneda Lemma says that the functor $\opn{Y}_{\cat{C}}^{\vee}$ is 
an equivalence from $\cat{C}^{\mrm{op}}$ to the category of corepresentable 
functors. The identity automorphism $\opn{id}_C$ corresponds to 
a special element $\til{\eta} \in F(C)$.
 


\mysubsection{Inverse and Direct Limits} \label{subsec:inv-dir-lim}
 
We are only interested in direct and inverse limits indexed by the ordered set 
$\N$. For a general discussion see \cite{Mac1} or \cite{KaSc1}. 

Let $\cat{C}$ be a category. Recall that an 
{\em $\N$-indexed direct system} in $\cat{C}$ is data  
\[ \bigl( \{ C_k \}_{k \in \N}, \{ \mu_{k} \}_{k \in \N} \bigr) , \]
where $C_k$ are objects of $\cat{C}$, and 
$\mu_k : C_k \to C_{k + 1}$
are morphisms, that we call {\em transitions}.
A {\em direct limit} of this system is  data  
$\bigl( C , \{ \ep_{k} \}_{k \in \N} \bigr)$,
where $C \in \cat{C}$, and $\ep_k : C_k \to C$ are morphisms, that we call 
{\em abutments}, such that $\ep_{k + 1} \circ \mu_k = \ep_k$ for all $k$. The 
universal property required is this: if 
$\bigl( C' , \{ \ep'_{k} \}_{k \in \N} \bigr)$
is another pair such that $\ep'_{k + 1} \circ \mu_k = \ep'_k$, then
there is a unique morphism 
$\ep : C \to C'$ such that 
$\ep'_{k} = \ep \circ \ep_k$. 
If a direct limit $C$ exists, then of course it is unique, up to a unique 
isomorphism. We then write 
$\lim_{k \to} C_k := C$,
and call this the direct limit of the system $\{ C_k \}_{k \in \N}$, keeping 
the transitions and the abutments implicit. 
Sometimes we look at the morphisms 
\[ \mu_{k_0, k_1} := \mu_{k_1 - 1} \circ \cdots \circ \mu_{k_0} :
C_{k_0} \to C_{k_1} \]
for $k_0 < k_1$, and $\mu_{k, k} := \opn{id}_{C_k}$. 

By an {\em $\N$-indexed inverse system} in the category $\cat{C}$ we mean 
data 
\[ \bigl( \{ C_k \}_{k \in \N}, \{ \mu_k \}_{k \in \N} \bigr) , \]
where $\{ C_k \}_{k \in \N}$ is a collection of objects, and 
$\mu_k : C_{k + 1} \to C_{k}$ 
are morphisms, also called transitions. 
An {\em inverse limit} of this system is data  
$\bigl( C , \{ \ep_{k} \}_{k \in \N} \bigr)$,
where $C \in \cat{C}$, and $\ep_k : C \to C_k$ are morphisms, that we also call 
abutments, such that $\mu_k \circ \ep_{k + 1} = \ep_{k}$. These satisfy an 
analogous universal property. 
If an inverse limit $C$ exists, then it is unique, up to a unique 
isomorphism. We then write 
$\lim_{\lar k} C_k := C$,
and we call this the inverse limit of the system $\{ C_k \}_{k \in \N}$.
We define the morphisms 
\[ \mu_{k_0, k_1} := \mu_{k_0} \circ \cdots \circ \mu_{k_1 - 1} :
C_{k_1} \to C_{k_0} \]
for $k_0 < k_1$, and $\mu_{k, k} := \opn{id}_{C_k}$. 

\begin{exer} \label{exer:3010}
We can view the ordered set $\N$ as a category, with a single 
morphism $k \to l$ when $k \leq l$, and no morphisms otherwise. 
\begin{enumerate}
\item Interpret $\N$-indexed direct and inverse systems in $\cat{C}$ as 
functors $F : \N \to \cat{C}$ and $G : \N^{\mrm{op}} \to \cat{C}$ respectively. 

\item Let $\bar{\N} := \N \cup \{ \infty \}$. 
Interpret the direct and inverse limits of $F$ and $G$ respectively as 
functors $\bar{F} : \bar{\N} \to \cat{C}$ and 
$\bar{G} : \bar{\N}^{\mrm{op}} \to \cat{C}$,
extending $F$ and $G$, with suitable universal properties. 
\end{enumerate}
\end{exer}

\begin{exer} \label{exer:3011}
Prove that $\N$-indexed direct and inverse limits exist in the categories 
$\cat{Set}$ and $\cat{Mod} A$, for any ring $A$. Give explicit formulas.
\end{exer}

\begin{exa} \label{ea:1740}
Let $\cat{M}$ be the category of finite abelian groups. 
The inverse system $\{ M_k \}_{k \in \N}$, where
$M_k := \Z / (2^k)$, and the transition 
$\mu_k :  M_{k + 1} \to M_k$ is the canonical surjection, 
does not have an inverse limit in $\cat{M}$. We can also make 
$\{ M_k \}_{k \in \N}$ into a direct system, in which the transition
$\nu_k : M_k \to  M_{k + 1}$ 
is multiplication by $2$. The direct limit does not exist in $\cat{M}$.
\end{exa}

If  $\{ C_k \}_{k \in \N}$ is a direct system in $\cat{C}$, and 
$D \in \cat{C}$ is any object, then there is an induced inverse system 
$\bigl\{ \opn{Hom}_{\cat{C}}(C_k, D) \}_{k \in \N}$
in $\cat{Set}$, and it has a limit. If 
$C := \lim_{k \to} C_k$ exists, then the abutments $\ep_k : C_k \to C$ induce 
a morphism 
\begin{equation} \label{eqn:3010}
\opn{Hom}_{\cat{C}}(C, D) \to 
\lim_{\leftarrow k} \, \opn{Hom}_{\cat{C}}(C_k, D) 
\end{equation}
in $\cat{Set}$. 

Similarly, if $\{ C_k \}_{k \in \N}$ is an inverse system in $\cat{C}$, 
and $D \in \cat{C}$ is any object, then there is an induced inverse system 
$\bigl\{ \opn{Hom}_{\cat{C}}(D, C_k) \}_{k \in \N}$
in $\cat{Set}$, and it has a limit. If $C := \lim_{\lar k} C_k$ exists,
then the abutments $\ep_k : C \to C_k$ induce a morphism 
\begin{equation} \label{eqn:3011}
\opn{Hom}_{\cat{C}}(D, C) \to 
\lim_{\leftarrow k} \, \opn{Hom}_{\cat{C}}(D, C_k) , 
\end{equation}
in $\cat{Set}$. 


\begin{prop} \label{prop:1535}
Let $\cat{C}$ be a category. 
\begin{enumerate}
\item Let  $\{ C_k \}_{k \in \N}$ be a direct system in $\cat{C}$,
and let $\bigl( C , \{ \ep_{k} \}_{k \in \N} \bigr)$ be data as in the 
definition of a direct limit. Then  
$C = \lim_{k \to} C_k$ if and only if for every object 
$D \in \cat{C}$, the function \tup{(\ref{eqn:3010})} is bijective. 

\item Let  $\{ C_k \}_{k \in \N}$ be an inverse system in $\cat{C}$,
and let $\bigl( C , \{ \ep_{k} \}_{k \in \N} \bigr)$ be data as in the 
definition of an inverse limit. Then  
$C = \lim_{\lar k} C_k$ if and only if for every object 
$D \in \cat{C}$, the function \tup{(\ref{eqn:3011})} is bijective. 
\end{enumerate}
\end{prop}

\begin{exer}  \label{exer:1535}
Prove Proposition \ref{prop:1535}.
\end{exer}

\begin{rem} \label{rem:4090}
In many cases inverse and direct limits do not exist in $\cat{C}$ because ``its 
objects are too small''. This happens in the category 
$\cat{Set}_{\mrm{fin}}$ of finite sets, and also in the category 
$\catt{Ab}_{\mrm{fin}}$ of finite abelian groups (see Example \ref{ea:1740} 
above).  

There is a very effective method to enlarge $\cat{C}$ just enough so that the 
bigger category will have the desired limits. This is done by means of the 
categories $\cat{Ind}_{}(\cat{C})$ and $\cat{Pro}_{}(\cat{C})$ of {\em 
ind-objects} and {\em pro-objects} of $\cat{C}$, respectively. 
See [KaSc1, Sec 1.11] and [KaSc2, Sec 6.1] for detailed discussions.

Here are two examples. For $\cat{C} = \cat{Set}_{\mrm{fin}}$, the 
category $\cat{Ind}_{}(\cat{Set}_{\mrm{fin}})$ is (canonically equivalent to) 
$\cat{Set}$. 
The category $\cat{Pro}_{}(\cat{Set}_{\mrm{fin}})$ is the category of 
totally disconnected compact Hausdorff topological spaces, 
whose morphisms are the continuous functions. 

For $\cat{C} = \catt{Ab}_{\mrm{fin}}$, 
the category $\cat{Ind}_{}(\cat{Ab}_{\mrm{fin}})$ of ind-objects is the 
category $\cat{Ab}$ of abelian groups. The category 
$\cat{Pro}_{}(\cat{Ab}_{\mrm{fin}})$ of pro-objects is the category of profinite
abelian groups, whose morphisms are the continuous group homomorphisms. 
\end{rem}






\cleardoublepage
\mysection{Abelian Categories and Additive Functors} \label{sec:abelian}

\AYcopyright

The concept of {\em abelian category} is an extremely useful abstraction of a
category of modules. It was introduced by A. Grothendieck in his foundational 
paper \cite{Gro} from 1957. 
References for this material are \cite{Mac1}, \cite{Mac2}, 
\cite{Rot}, \cite{HiSt}, \cite{We}, \cite{KaSc1} and \cite{KaSc2}.


\mysubsection{Linear Categories} \label{subsec:1}

\begin{dfn} \label{dfn:1083}
Let $\K$ be a commutative ring. 
A  {\em $\K$-linear category}%
\index{Linear category}%
\index{Linear category!K@$\K$-linear category}
is a category $\cat{M}$, endowed with 
a $\K$-module structure on each of the sets of morphisms  
$\opn{Hom}_{\cat{M}}(M_0, M_1)$. The condition is this:
\begin{itemize}
\item For all $M_0, M_1, M_2 \in \cat{M}$ the composition function 
\[ \begin{aligned}
& \opn{Hom}_{\cat{M}}(M_1, M_2) \times \opn{Hom}_{\cat{M}}(M_0, M_1)
\to \opn{Hom}_{\cat{M}}(M_0, M_2) \\
& (\phi_1, \phi_0) \mapsto \phi_1 \circ \phi_0 
\end{aligned} \]
is $\K$-bilinear.
\end{itemize}
If $\K = \Z$, we say that $\cat{M}$ is a {\em linear category}.
\end{dfn}

As already mentioned in Convention \ref{conv:2490}, in our book all linear 
categories are $\K$-linear. 

\begin{prop} \label{prop:1035}
Let $\cat{M}$ be a $\K$-linear category.
\begin{enumerate}
\item For every object $M \in \cat{M}$, the set 
$\opn{End}_{\cat{M}}(M) := \opn{Hom}_{\cat{M}}(M, M)$,
with its given addition operation, and with the operation of composition,
is a central $\K$-ring. 

\item For every two objects $M_0, M_1 \in \cat{M}$, the set 
$\opn{Hom}_{\cat{M}}(M_0, M_1)$,
with its given addition operation, and with the operations of composition,
is a left module over the ring $\opn{End}_{\cat{M}}(M_1)$, and 
a right module over the ring $\opn{End}_{\cat{M}}(M_0)$.
Furthermore, these left and right actions commute with each other. 
\end{enumerate}
\end{prop}

\begin{exer} \label{exer:3255}
Prove Proposition \ref{prop:1035}.
\end{exer}

This result can be reversed:

\begin{exa} \label{exa:104}
Let $A$ be a central $\K$-ring.  
Define a category $\cat{M}$ like this: there is a single object $M$, and its
set of morphisms is
$\opn{Hom}_{\cat{M}}(M, M) := A$.
Composition in $\cat{M}$ is the multiplication of $A$. Then $\cat{M}$ is a
$\K$-linear category.
\end{exa}

For a central $\K$-ring $A$, the opposite ring $A^{\mrm{op}}$ has the same 
$\K$-module structure as $A$, but the multiplication is reversed. 

\begin{exer} \label{exer:1030}
Let $A$ be a nonzero ring. Let $P, Q \in \cat{Mod} A$ 
be distinct free $A$-modules of rank $1$.
\begin{enumerate}
\item Prove that there is a ring isomorphism 
$\opn{End}_{\cat{Mod} A}(P) \cong A^{\mrm{op}}$.
Is this ring isomorphism canonical?

\item Let $\cat{M}$ be the full subcategory of $\cat{Mod} A$ on the set of 
objects $\{ P, Q \}$. Compare the linear category $\cat{M}$ to the ring of 
matrices $\opn{Mat}_2(A^{\mrm{op}})$.
\end{enumerate}
\end{exer}
 


\mysubsection{Additive Categories}

\begin{dfn} 
An  {\em additive category}%
\index{Additive category}
is a linear category $\cat{M}$
satisfying these two conditions:
\begin{enumerate}
\rmitem{i} $\cat{M}$ has a zero object $0$.

\rmitem{ii} $\cat{M}$ has finite coproducts.
\end{enumerate}
\end{dfn}

Observe that $\opn{Hom}_{\cat{M}}(M, N) \neq \varnothing$
for every $M, N \in \cat{M}$, since this is an abelian group. For the zero 
object $0 \in \cat{M}$ we have 
$\opn{Hom}_{\cat{M}}(M, 0) = \opn{Hom}_{\cat{M}}(0, M) = 0$, 
the zero abelian group.  We denote the unique arrows $0 \to M$ and 
$M \to 0$ also by $0$. So the numeral $0$ has a lot of meanings; but they are 
always (hopefully) clear from the context. 
The coproduct in a linear category $\cat{M}$ is usually denoted by $\bigoplus$, 
and is called the {\em direct sum}; cf.\ Example \ref{exa:103}.

\begin{exa}
Let $A$ be a central $\K$-ring. The category $\cat{Mod} A$ is a $\K$-linear 
additive category. 
The full subcategory $\cat{F} \subseteq \cat{Mod} A$ on the free modules is
also additive. 
\end{exa}

\begin{prop} \label{prop:1}
Let $\cat{M}$ be a linear category. Let $\{ M_i \}_{i \in I}$ be a  
finite collection of objects of $\cat{M}$, and assume the coproduct
$M = \bigoplus_{i \in I} M_i$ exists, with embeddings 
$e_i : M_i \to M$. 

\begin{enumerate}
\item For each $i$ let $p_i : M \to M_i$ be the unique morphism 
such that $p_i \circ e_i = \opn{id}_{M_i}$, and 
 $p_i \circ e_j = 0$ for $j \neq i$. 
Then $(M, \{ p_i \}_{i \in I})$ is a product of the collection 
$\{ M_i \}_{i \in I}$.

\item $\sum_{i \in I} e_i \circ p_i = \opn{id}_M$.
\end{enumerate}
\end{prop}

\begin{exer} \label{exer:2490}
Prove this proposition.
\end{exer}

Part (1) of Proposition \ref{prop:1} directly implies:

\begin{cor} \label{cor:1030}
An additive category has finite products. 
\end{cor}

\begin{dfn} \label{dfn:1025}
Let $\cat{M}$ be an additive category, and let $\cat{N}$ be a full subcategory
of $\cat{M}$. We say that $\cat{N}$ is a {\em full additive subcategory} of
$\cat{M}$ if $\cat{N}$ contains the zero object, and is closed under finite 
direct sums.
\end{dfn}

\begin{exer}
In the situation of Definition \ref{dfn:1025}, show that the category $\cat{N}$ 
is itself additive.
\end{exer}

\begin{exa} \label{exa:1030}
Consider the linear category $\cat{M}$ from Example \ref{exa:104},
built from a ring $A$. It does not have a zero object (unless the ring $A$ is 
the zero ring), so it is not additive. 

A more puzzling question is this: Does $\cat{M}$ have finite direct 
sums? This turns out to be equivalent to whether or not $A \cong A \oplus A$ as 
right $A$-modules. 
One can show that when $A$ is nonzero and commutative, or nonzero and 
noetherian, then $A \not\cong A \oplus A$ in $\cat{Mod} A^{\mrm{op}}$.
On the other hand, if we take a field $\K$, and a countable rank $\K$-module
$N$, then $A := \opn{End}_{\K}(N)$ will satisfy $A \cong A \oplus A$.
\end{exa}


\begin{prop} \label{prop:1030}
Let $\cat{M}$ be a linear category, and $N \in \cat{M}$. The following 
conditions are equivalent\tup{:}
\begin{enumerate}
\rmitem{i} The ring  $\opn{End}_{\cat{M}}(N)$ is trivial. 
\rmitem{ii} $N$ is a zero object of $\cat{M}$.
\end{enumerate}
\end{prop}

\begin{proof} \mbox{}

\smallskip \noindent
(ii) $\Rightarrow$ (i): Since the set $\opn{End}_{\cat{M}}(N)$ is a singleton, 
it must be the trivial ring. 

\medskip \noindent 
(i) $\Rightarrow$ (ii): If the ring $\opn{End}_{\cat{M}}(N)$ is 
trivial, then all left and right modules over it must be trivial. 
Now use Proposition \ref{prop:1035}(2). 
\end{proof}



\mysubsection{Abelian Categories} \label{subsec:ab-cats}

\begin{dfn}
Let $\cat{M}$ be an additive category, and let $f : M \to N$ be a morphism in
$\cat{M}$. A {\em kernel} of $f$ is a pair $(K, k)$, consisting of an object
$K \in \cat{M}$ and a morphism $k : K \to M$, with these two properties:
\begin{enumerate}
\rmitem{i}  $f \circ k = 0$. 
\rmitem{ii} If $k' : K' \to M$ is a morphism in $\cat{M}$ such that 
$f \circ k' = 0$, then there is a unique morphism $g : K' \to K$ 
such that $k' = k \circ g$.
\end{enumerate}
\end{dfn}

In other words, the object $K$ represents the functor 
$\cat{M}^{\mrm{op}} \to \cat{Ab}$, 
\[ K' \mapsto \{ k' \in \opn{Hom}_{\cat{M}}(K', M) \mid f \circ k' = 0 \} . \]
The kernel of $f$ is of course unique up to a unique isomorphism (if it
exists), and we denote if by $\opn{Ker}(f)$. Sometimes $\opn{Ker}(f)$ refers
only to the object $K$, and other times it refers only to the morphism
$k$; as usual, this should be clear from the context.

\begin{dfn}
Let $\cat{M}$ be an additive category, and let $f : M \to N$ be a morphism in
$\cat{M}$. A {\em cokernel} of $f$ is a pair $(C, c)$, consisting of an object
$C \in \cat{M}$ and a morphism $c : N \to C$, with these two properties:
\begin{enumerate}
\rmitem{i}  $c \circ f = 0$. 
\rmitem{ii} If $c' : N \to C'$ is a morphism in $\cat{M}$ such that 
$c' \circ f = 0$, then there is a unique morphism $g : C \to C'$ 
such that $c' = g \circ c$.
\end{enumerate}
\end{dfn}

In other words, the object $C$ corepresents the functor 
$\cat{M} \to \cat{Ab}$, 
\[ C' \mapsto \{ c' \in \opn{Hom}_{\cat{M}}(N, C') \mid c' \circ f = 0 \} . \]
The cokernel of $f$ is of course unique up to a unique isomorphism (if it
exists), and we denote if by $\opn{Coker}(f)$. Sometimes $\opn{Coker}(f)$ 
refers only to the object $C$, and other times it refers only to the morphism
$c$; as usual, this should be clear from the context.

\begin{exa}
In $\cat{Mod} A$ all kernels and cokernels exist. 
Given $f : M \to N$, the kernel is 
$k : K \to M$, where 
$K := \{ m \in M \mid f(m) = 0 \}$,
and the $k$ is the inclusion. The cokernel is $c : N \to C$, where 
$C := N / f(M)$, and $c$ is the canonical projection.
\end{exa}

\begin{prop}
Let  $\cat{M}$ be an additive category, and let $f : M \to N$ be a morphism 
in $\cat{M}$.
\begin{enumerate}
\item If $k : K \to M$ is a kernel of $f$, then $k$  is a monomorphism.
\item If $c : N \to C$ is a cokernel of $f$, then $c$ is an epimorphism.
\end{enumerate}
\end{prop}

\begin{exer} 
Prove the proposition.
\end{exer}

\begin{dfn}
Assume the additive category $\cat{M}$ has kernels and cokernels. 
Let $f : M \to N$ be a morphism in $\cat{M}$.
\begin{enumerate}
\item Define the {\em image} of $f$ to be 
$\opn{Im}(f) := \opn{Ker}(\opn{Coker}(f))$.

\item Define the {\em coimage} of $f$ to be 
$\opn{Coim}(f) := \opn{Coker}(\opn{Ker}(f))$.
\end{enumerate}
\end{dfn}

The image is familiar, but the coimage is probably not. The next diagram 
should help. 
We start with a morphism $f : M \to N$ in $\cat{M}$. 
The kernel and cokernel of $f$ fit into this diagram:
$K \xar{k} M \xar{f} N \xar{c} C$. 
Inserting $\al := \opn{Coker}(k) = \opn{Coim}(f)$ and 
$\be := \opn{Ker}(c) = \opn{Im}(f)$
we get the following commutative diagram (solid arrows):
\begin{equation} \label{eqn:1015}
\UseTips  \xymatrix @C=8ex @R=6ex {
K
\ar[r]^{k}
\ar[dr]_{0}
& 
M
\ar[r]^{f}
\ar[d]_{\al}
\ar@{-->}[dr]^{\ga}
&
N
\ar[r]^{c}
&
C
\\
&
M'
\ar@{-->}[r]_{f'}
&
N'
\ar[ur]_{0}
\ar[u]_{\be}
}
\end{equation}
Since $c \circ f = 0$ there is a unique morphism $\ga$ (the dashed arrow)
making the diagram commutative. Now 
$\be \circ \ga \circ k = f \circ k = 0$; and $\be$ is a monomorphism; so 
$\ga \circ k = 0$. Hence there is a unique morphism $f' : M' \to N'$ 
making the diagram commutative.
We conclude that $f : M \to N$ induces a morphism 
\begin{equation}  \label{eqn:2}
f' : \opn{Coim}(f) \to \opn{Im}(f) . 
\end{equation}

\begin{dfn} \label{dfn:1}
An {\em abelian category}%
\index{Abelian category}
is an additive category $\cat{M}$ with these two extra properties:
\begin{enumerate}
\rmitem{i} All morphisms in $\cat{M}$ admit kernels and cokernels.
\rmitem{ii} For every morphism $f : M \to N$ in $\cat{M}$, the induced morphism 
$f'$ in equation (\ref{eqn:2}) is an isomorphism.
\end{enumerate}
\end{dfn}

Here is a less precise but (maybe) easier to remember way to state property 
(ii). Because $M' = \opn{Coker}(\opn{Ker}(f))$ and 
$N' = \opn{Ker}(\opn{Coker}(f))$, 
we see that 
\begin{equation} \label{eqn:1019}
\opn{Coker}(\opn{Ker}(f)) = \opn{Ker}(\opn{Coker}(f))  .
\end{equation}

{}From now on we forget all about the coimage. 

\begin{exer}
For any ring $A$, prove that the category $\cat{Mod} A$ is abelian. 
\end{exer}

This includes the category $\cat{Ab} = \cat{Mod} \Z$, from which the name 
derives. 

\begin{dfn} \label{dfn:1015}
Let $\cat{M}$ be an abelian category, and let $\cat{N}$ be a full subcategory
of $\cat{M}$. We say that $\cat{N}$ is a 
{\em full abelian subcategory}%
\index{Abelian category! full  abelian subcategory of}
of $\cat{M}$ if the zero object belongs to $\cat{N}$, and $\cat{N}$ is closed  
in $\cat{M}$ under taking finite direct sums, kernels and cokernels.
\end{dfn}

\begin{exer}
In the situation of Definition \ref{dfn:1015}, show that the category $\cat{N}$ 
is itself abelian.
\end{exer}

\begin{exa}
Let $\cat{M}_1$ be the category of finitely generated abelian groups,
and let $\cat{M}_0$ be the category of finite abelian groups. 
Then $\cat{M}_1$ is a full abelian subcategory of $\cat{Ab}$, and 
$\cat{M}_0$ is a full abelian subcategory of $\cat{M}_1$. 
\end{exa}

\begin{exer}
Let $\cat{N}$ be the full subcategory of $\cat{Ab}$ whose objects are the
finitely generated free abelian groups. It is an additive subcategory of
$\cat{Ab}$  (since it is closed under direct sums).
\begin{enumerate}
\item Show that $\cat{N}$ is closed under kernels in $\cat{Ab}$.

\item Show that $\cat{N}$ is not closed under cokernels in $\cat{Ab}$, so it 
is not a full abelian subcategory of $\cat{Ab}$.

\item Show that $\cat{N}$ has cokernels (not the same as those of $\cat{Ab}$).
Still, it fails to be an abelian category.
\end{enumerate}
\end{exer}

\begin{exer}
The category $\cat{Grp}$ is not linear of course. Still, it does have a 
zero object (the trivial group). Show that $\cat{Grp}$ has kernels and 
cokernels, but condition (ii) of Definition \ref{dfn:1} fails. 
\end{exer}

\begin{exer}
Let $\cat{Hilb}$ be the category of Hilbert spaces over $\mbb{C}$. 
The morphisms are the continuous $\mbb{C}$-linear homomorphisms. Show that 
$\cat{Hilb}$ is a $\mbb{C}$-linear additive category with kernels and 
cokernels, 
but it is not an abelian category.
\end{exer}

\begin{exer}
Let $A$ be a ring. Show that $A$ is {\em left noetherian} iff the category 
$\cat{Mod}_{\mrm{f}} A$ of  finitely generated left modules is a full abelian 
subcategory of $\cat{Mod} A$. 
\end{exer}

\begin{exa} \label{exa:1}
Let $(X, \mcal{A})$ be a {\em ringed space}; namely $X$ is a topological space 
and $\mcal{A}$ is a sheaf of rings on $X$. 
There is a very detailed discussion of sheaves of modules in 
\cite[Section II.1]{Har} and \cite[Sections 2.1-2.2]{KaSc1}. 
Let us mention only one aspect of it. 

We denote by $\cat{Mod}^{\mrm{pr}} \mcal{A}$ the category of {\em presheaves} 
of left $\mcal{A}$-modules on $X$. This is an abelian category. 
Given a morphism $\phi : \mcal{M} \to \mcal{N}$ in 
$\cat{Mod}^{\mrm{pr}} \mcal{A}$, its kernel is the presheaf 
$\opn{Ker}^{\mrm{pr}}(\phi)$ defined by 
\[ \Gamma \bigl( U, \opn{Ker}^{\mrm{pr}}(\phi) \bigr) := 
\opn{Ker} \big( \phi : \Gamma(U, \mcal{M}) \to \Gamma(U, \mcal{N}) \big)  \]
on every open set $U \subseteq X$. 
The cokernel of $\phi$ is the  presheaf $\opn{Coker}^{\mrm{pr}}(\phi)$
defined by 
\[ \Gamma \bigl( U,  \opn{Coker}^{\mrm{pr}}(\phi) \bigr) : 
\opn{Coker} \big( \phi : \Gamma(U, \mcal{M}) \to \Gamma(U, \mcal{N}) \big) . \]

Now let $\cat{Mod} \mcal{A}$ be the full subcategory of 
$\cat{Mod}^{\mrm{pr}} \mcal{A}$ consisting of {\em sheaves}. 
It is a full additive subcategory of $\cat{Mod}^{\mrm{pr}} \mcal{A}$, closed 
under kernels. However, $\cat{Mod} \mcal{A}$ is not closed under 
cokernels inside $\cat{Mod}^{\mrm{pr}} \mcal{A}$,
and hence it is not a full abelian subcategory.

Nonetheless, $\cat{Mod} \mcal{A}$ is itself an abelian category, but
with different cokernels. Given a morphism 
$\phi : \mcal{M} \to \mcal{N}$ in $\cat{Mod} \mcal{A}$,
its cokernel in $\cat{Mod} \mcal{A}$ is the sheaf $\opn{Coker}(\phi)$ that's  
associated to the presheaf $\opn{Coker}^{\mrm{pr}}(\phi)$.
\end{exa}

\begin{prop} \label{prop:2490}
Let $\cat{M}$ be a linear category.
\begin{enumerate}
\item The opposite category $\cat{M}^{\mrm{op}}$ has a canonical structure 
of a linear category. 

\item If $\cat{M}$ is additive, then $\cat{M}^{\mrm{op}}$ is also additive.

\item If $\cat{M}$ is abelian, then $\cat{M}^{\mrm{op}}$ is also abelian.
\end{enumerate}
\end{prop}

\begin{proof}
(1) Since
$\opn{Hom}_{\cat{M}^{\mrm{op}}}(M, N) = \opn{Hom}_{\cat{M}}(N, M)$,
this is an abelian group. The bilinearity of the composition in 
$\cat{M}^{\mrm{op}}$ is clear.

\medskip \noindent 
(2) The zero objects in $\cat{M}$ and $\cat{M}^{\mrm{op}}$ 
are the same. Existence of finite coproducts in 
$\cat{M}^{\mrm{op}}$ is because of existence of finite products in $\cat{M}$;
see Proposition \ref{prop:1}(1). 

\medskip \noindent 
(3) $\cat{M}^{\mrm{op}}$ has kernels and cokernels, since 
$\opn{Ker}_{\cat{M}^{\mrm{op}}}(\opn{Op}(\phi)) = \opn{Coker}_{\cat{M}}(\phi)$ 
and vice versa. Also the symmetric condition (ii) of Definition \ref{dfn:1} 
holds. 
\end{proof}

\begin{prop}
Let $\phi : M \to N$ be a morphism in an abelian category $\cat{M}$. 
\begin{enumerate}
\item $\phi$ is a monomorphism if and only if $\opn{Ker}(\phi) = 0$. 
\item $\phi$ is an epimorphism  if and only if $\opn{Coker}(\phi) = 0$. 
\item $\phi$ is an isomorphism  if and only if it is both a monomorphism and an 
epimorphism.
\end{enumerate}
\end{prop}

\begin{exer}
Prove this proposition.
\end{exer}

Consider a diagram 
\begin{equation} \label{eqn:3665}
\bsym{S} = \bigl( \cdots  M_{-1} \xar{\phi_{-1}}
M_{0} \xar{\phi_0} M_{1} \xar{\phi_{1}} M_{2} \cdots \bigr)
\end{equation}
in an abelian category $\cat{M}$, extending finitely or infinitely to either 
side. Such a diagram is called a {\em sequence in $\cat{M}$}. 
An object $M_i$ appearing in $\bsym{S}$ is called {\em interior in $\bsym{S}$} 
if there is an object $M_{i - 1}$ appearing to the left of it, and an object 
$M_{i + 1}$ appearing to the right of it. 

\begin{dfn} \label{dfn:3635}
\index{Exact sequence}
Let $\bsym{S}$ be a sequence in the abelian category $\cat{M}$, with notation 
as in (\ref{eqn:3665}). 
\begin{enumerate}
\item Suppose $M_i$ is an interior object in $\bsym{S}$.
We say that the sequence $\bsym{S}$ is {\em exact at $M_i$} if 
$\opn{Im}(\phi_{i - 1}) =  \opn{Ker}(\phi_{i})$,
as subobjects of $M_i$. 

\item The sequence $\bsym{S}$ is said to be {\em exact} if it is exact at all 
of its interior objects.
\end{enumerate}
\end{dfn}

\begin{exa} \label{exa:3635}
A morphism $\phi : M \to N$ in an abelian category $\cat{M}$ is a monomorphism 
iff $0 \to M \xar{\phi} N$
is an exact sequence. The morphism $\phi$ is  an epimorphism iff the sequence
$M \xar{\phi} N \to 0$ is exact. 
\end{exa}

\begin{dfn} \label{dfn:4095}
A {\em short exact sequence} in an abelian category $\cat{M}$ is as exact 
sequence of the form 
\[ \bsym{S} = \bigl( 0 \to M_0 \xar{\phi_0} M_1 \xar{\phi_1} M_2 \to 0 \bigr) . 
\] 
\end{dfn}

\begin{prop} \label{prop:3637}
Let $\cat{M}$ be a $\K$-linear abelian category.
\begin{enumerate}
\item Let 
$0 \to M' \xar{\phi} M \xar{\psi} M''$
be an exact sequence in $\cat{M}$. Then for every $L \in \cat{M}$ the sequence 
\[ 0 \to \opn{Hom}_{\cat{M}}(L, M') \xar{\opn{Hom}(\opn{id}, \phi)} 
\opn{Hom}_{\cat{M}}(L, M) \xar{\opn{Hom}(\opn{id}, \psi)}
\opn{Hom}_{\cat{M}}(L, M'') \]
in $\cat{Mod} \K$ is exact. 

\item Let 
$M' \xar{\phi} M \xar{\psi} M'' \to 0$
be an exact sequence in $\cat{M}$. Then for every $N \in \cat{M}$ the sequence 
\[ 0 \to \opn{Hom}_{\cat{M}}(M'', N) \xar{\opn{Hom}(\psi, \opn{id})} 
\opn{Hom}_{\cat{M}}(M, N) \xar{\opn{Hom}(\phi, \opn{id})}
\opn{Hom}_{\cat{M}}(M', N) \]
in $\cat{Mod} \K$ is exact. 
\end{enumerate}
\end{prop}

\begin{exer} \label{exer:3635}
Prove Proposition \ref{prop:3637}. (Hint: use the definitions of kernel, 
cokernel and image.)
\end{exer}

 
\mysubsection{A Method for Producing Proofs in Abelian Categories} 
\label{subsec:sheaf-tricks}

A well-known difficulty in the theory of abelian categories is this: formulas 
that are easy to prove for a module category 
$\cat{M} = \cat{Mod} A$, using {\em elements}, are often 
very hard to prove in an abstract abelian category $\cat{M}$ (directly from the 
axioms). 

A neat solution to this difficulty was found by P. Freyd and B. Mitchell:

\begin{thm}[Freyd-Mitchell]
Let $\cat{M}$ be a small abelian category. Then $\cat{M}$ is equivalent to a
full abelian subcategory of $\cat{Mod} A$, for a suitable ring $A$. 
\end{thm}

\begin{rem} \label{rem:4091}
This is a deep and difficult result. See \cite{Fre}, a book that's basically 
devoted to proving this theorem. 

A modern proof can be found in
in \cite[Theorem 9.6.10]{KaSc2} -- but it too is very involved. Roughly 
speaking, they show that the abelian category $\cat{Ind}(\cat{M}^{\mrm{op}})$ 
of ind-objects of $\cat{M}$ has an injective cogenerator. This implies that the 
abelian category $\cat{Pro}(\cat{M})$ of pro-objects, that's equivalent 
to $\cat{Ind}(\cat{M}^{\mrm{op}})^{\mrm{op}}$,
has a projective generator, say $P$. Defining the ring 
$A := \opn{End}_{\cat{Pro}(\cat{M})}(P)^{\mrm{op}}$,
there is an equivalence of abelian categories
$\cat{Pro}(\cat{M}) \approx \cat{Mod} A$.
On the other hand, the Yoneda functor is a fully faithful embedding 
$\cat{M} \to \cat{Pro}(\cat{M})$.
\end{rem}

The Freyd-Mitchell Theorem implies that for 
purposes of finitary calculations in the abelian category $\cat{M}$ 
(e.g.\ checking whether a sequence is exact, 
see Definition \ref{dfn:3635}) we can assume that objects of $\cat{M}$ 
have elements. This often simplifies the work.

Since we do not give a proof of the Freyd-Mitchell Theorem in our book, 
we feel it is improper to use it. As a substitute, we provide the two 
``sheaf tricks'' below, namely Propositions \ref{prop:3635} and 
\ref{prop:3636}, with full proofs. Later in the book these sheaf tricks are 
used to give relatively easy proofs of several results on abstract abelian 
categories, most notably Theorem \ref{thm:3640} 
on the existence of the long exact 
cohomology sequence. The method of proof using these tricks is explained in  
Remark \ref{rem:3655}. It is not as slick as the method that the Freyd-Mitchell 
Theorem offers; but at least we have self-contained proofs. See Remark 
\ref{rem:3630} for some background (influence and history) on the sheaf tricks. 

First an important technical lemma. In an abelian category $\cat{M}$ we have 
finite products, and they are also coproducts (see Proposition 
\ref{prop:1}). 
 
\begin{lem} \label{lem:3630}
Let $\cat{M}$ be an abelian category. Consider a diagram 
\[ \tag{D}
\UseTips \xymatrix @C=6ex @R=6ex {
M_1 \times M_2
\ar[r]^(0.6){p_2}
\ar[d]_{p_1}
&
M_2
\ar[d]^{\phi_2}
\\
M_1
\ar[r]^{\phi_1}
&
N
} \]
in $\cat{M}$, where $p_i$ are the projections. 
Define the object 
\[ L := \opn{Ker}(\phi_1 \circ p_1 - \phi_2 \circ p_2) \sub M_1 \times M_2 , \]
with inclusion morphism 
$e : L \to M_1 \times M_2$.
Next define the morphisms 
$\psi_i := p_i \circ e : L \to M_i$. 
\begin{enumerate}
\item The diagram 
\[ \UseTips \xymatrix @C=6ex @R=6ex {
L
\ar[r]^{\psi_2}
\ar[d]_{\psi_1}
&
M_2
\ar[d]^{\phi_2}
\\
M_1
\ar[r]^{\phi_1}
&
N
} \]
is cartesian, and $L = M_1 \times_N M_2$. 

\item If $\phi_1$ is an epimorphism, then $\psi_2$ is an epimorphism.
\end{enumerate}
\end{lem}

Note that the diagram (D) is not assumed to be commutative. In case (D) does 
happen to be commutative, then 
$M_1 \times_N M_2 = M_1 \times M_2$. 

\begin{proof}
(1) The fact that $L$ (with the morphisms $\psi_i$) is the fibered product is 
immediate from the definitions of product and kernel. 

\medskip \noindent
(2) Here we follow [Mac2, Sec VIII.4]. 
Let $\rho$ be a morphism such that 
\[ \rho \circ (\phi_1 \circ p_1 - \phi_2 \circ p_2) = 0 . \]
Consider the embedding 
\[ e_1 : M_1 \to M_1 \oplus M_2 = M_1 \times M_2 . \]
Then 
\[ 0 = \rho \circ (\phi_1 \circ p_1 - \phi_2 \circ p_2) \circ e_1  
= \rho \circ \phi_1 . \]
Because $\phi_1$ is an epimorphism, it follows that $\rho = 0$. 
We conclude that 
\begin{equation} \label{eqn:3630}
\phi_1 \circ p_1 - \phi_2 \circ p_2 : M_1 \times M_2 \to N
\end{equation}
is an epimorphism. Thus (\ref{eqn:3630}) is the cokernel of $e$. 

Next let $\si : M_2 \to P$ be a morphism such that $\si \circ \psi_2 = 0$. 
Since $\psi_2 = p_2 \circ e$, we get 
$\si \circ p_2 \circ e = 0$. 
Hence $\si \circ p_2$ factors through $\opn{Coker}(e)$.
Namely there is a morphism $\si' : N \to P$ such that 
\begin{equation} \label{eqn:3626}
\si \circ p_2 = \si' \circ (\phi_1 \circ p_1 - \phi_2 \circ p_2) : 
M_1 \times M_2 \to P . 
\end{equation}
But $p_2 \circ e_1 = 0$, and therefore
\[ 0 = \si \circ p_2 \circ e_1 = 
\si' \circ (\phi_1 \circ p_1 - \phi_2 \circ p_2) \circ e_1 = 
\si' \circ \phi_1 . \]
As $\phi_1$ is an epimorphism, it follows that $\si' = 0$. 
Finally, using (\ref{eqn:3626}), we get 
\[ \si = \si \circ p_2 \circ e_2 = 
\si' \circ (\phi_1 \circ p_1 - \phi_2 \circ p_2) \circ e_2 = 
- \si' \circ \phi_2 = 0 . \]
We conclude that $\psi_2$ is an epimorphism. 
\end{proof}

Before giving the precise statements, here is a heuristic.% 
\index{Sheaf Tricks (in abelian categories)|(}

\begin{rem} \label{rem:3655}
The sheaf tricks (Propositions \ref{prop:3635} and \ref{prop:3636}) work like 
this: we pretend that the objects of our $\K$-linear abelian 
category $\cat{M}$ are ``sheaves on an imaginary topological space $X$''; 
objects playing this role are denoted by letters $M, N, \ldots$. 
The objects of $\cat{M}$ are also ``open sets in the topological space $X$'';
and objects playing this role are denoted by letters $U, V, \ldots$.
Given a sheaf $M$ and an open set $U$, there is a $\K$-module 
$\Ga(U, M)$ of ``sections of $M$ over $U$''. 
These ``sections'' can be added and subtracted, being elements of a 
$\K$-module. Sometimes we require ``refinement'', i.e.\ replacing an 
``open set $U$'' by a ``covering $V \surj U$'', giving rise to an embedding 
$\K$-modules $\Ga(U, M) \inj \Ga(V, M)$. This allegory is made precise in 
Definition \ref{dfn:3621} below. 
\end{rem}

\begin{dfn}[Sheaf Metaphor] \label{dfn:3621}
Let $\cat{M}$ be a $\K$-linear abelian category. 
\begin{enumerate}
\item Objects of $\cat{M}$ are called sheaves or open sets, depending on the 
role they play in each context. 

\item For an open set $U \in \cat{M}$ and a sheaf $M \in \cat{M}$ we write 
$\Ga(U, M) := \lb \opn{Hom}_{\cat{M}}(U, M)$.
This is a $\K$-module, and we call it the {\em module of sections of 
the sheaf $M$ over the open set $U$}.

\item Given a morphism $\rho: V \to U$ of open sets in $\cat{M}$, and a sheaf 
$M \in \cat{M}$, we use this notation for the resulting $\K$-module 
homomorphism:
\[ \rho^* := \opn{Hom}_{\cat{M}}(\rho, \opn{id}_M) : 
\Ga(U, M) \to \Ga(V, M) . \]
We call $\rho^*$ the {\em pullback along $\rho$}.  


\item Given a morphism $\phi : M \to N$ of sheaves in $\cat{M}$, and an open 
set $U \in \cat{M}$, we use this notation for the resulting $\K$-module 
homomorphism:
\[ \Ga(U, \phi) := \opn{Hom}_{\cat{M}}(\opn{id}_U, \phi) : 
\Ga(U, M) \to \Ga(U, N) . \]

\item If $\rho : V \to U$ is a morphism of open sets in $\cat{M}$ that is an  
epimorphism, then we call $\rho$ a {\em covering of $U$}. The homomorphism 
$\rho^*$ in this case is called the {\em restriction of $M$ from $U$ to $V$}.   
\end{enumerate}   
\end{dfn}

Note that given a covering $\rho : V \surj U$, the restriction 
homomorphism $\rho^*$ is injective (see Proposition \ref{prop:3637}(2)). By 
slight abuse of notation, when the covering $\rho$ is clear from the context, 
we will often identify the $\K$-module $\Ga(U, M)$ with its image in 
$\Ga(V, M)$ under $\rho^*$. 
Another convenient abuse of notation is writing
$\phi$ instead of $\Ga(U, \phi)$, for a morphism $\phi : M \to N$ in $\cat{M}$ 
and open set $U$. Combined, these notations give sense to such 
a commutative diagram
\begin{equation} \label{eqn:4935}
\UseTips \xymatrix @C=8ex @R=6ex {
\Ga(U, M)
\ar[r]^{\phi}
\ar[d]_{\sub}
&
\Ga(U, N)
\ar[d]^{\sub}
\\
\Ga(V, M)
\ar[r]^{\phi}
&
\Ga(V, N)
} 
\end{equation}
in $\cat{Mod} \K$. This will be used in item (3) of the next proposition. 

\begin{prop}[First Sheaf Trick] \label{prop:3635}
\index{Sheaf Tricks (in abelian categories)|)}
Let $\phi : M \to N$ be a morphism in an abelian category $\cat{M}$.
\begin{enumerate}
\item The morphism $\phi$ is zero if and only if
for every $U \in \cat{M}$ the homomorphism $\Ga(U, \phi)$ 
in $\cat{Mod} \K$ is zero.

\item  The morphism $\phi$ is a monomorphism if and only if for every $U \in 
\cat{M}$ the homomorphism $\Ga(U, \phi)$ 
in $\cat{Mod} \K$ is injective. 

\item  The morphism  $\phi$ is an epimorphism if and only if for every $U \in 
\cat{M}$ and every section $n \in \Ga(U, N)$, there exists a covering $V \surj 
U$ and a section $m \in \Ga(V, M)$, such that 
$\phi(m) = n$ in $\Ga(V, N)$.
\end{enumerate}
\end{prop}

The heuristic interpretation of item (3) of the proposition is this: 
$\phi : M \to N$ is an epimorphism in $\cat{M}$ if and only if it is ``locally 
surjective''.

\begin{proof}
(1) If $\phi = 0$ then for every $m \in \Ga(U, M)$ 
we have $\Ga(U, \phi)(m) = \phi \circ m = 0$ in $\Ga(U, N)$; thus 
$\Ga(U, \phi) = 0$. 
Conversely, if $\Ga(U, \phi) = 0$ for all $U$, then take $U := M$ and 
$m := \opn{id}_M \in \Ga(U, M)$. We obtain 
$\phi = \Ga(U, \phi)(m) = 0$. 

\medskip \noindent
(2) First assume $\phi$ is a monomorphism. Take an arbitrary object $U$ and 
a section $m \in \Ga(U, M)$. So $m : U \to M$ is a morphism in $\cat{M}$. 
 If $\phi \circ m = 0$, then, by the definition of a monomorphism, we must have 
$m = 0$. Thus $\Ga(U, \phi)$ is injective.

Conversely, assume that $\Ga(U, \phi)$ is injective for every $U$.
Take $U := \opn{Ker}(\phi)$ and let
$m : U \to M$ be the inclusion. Then 
$\Ga(U, \phi)(m) = \phi \circ m = 0$.
But then $m = 0$ and $U = 0$. Thus $\phi$ is a monomorphism. 

\medskip \noindent
(3) First assume $\phi$ is an epimorphism. Consider a section 
$n \in \Ga(U, N)$. So we have morphisms 
$\phi : M \surj N$ and $n : U \to N$. Let
$V := M \times_N U$, the fibered product. By Lemma \ref{lem:3630}(2)
the projection $\rho : V \to U$ is an epimorphism, i.e.\ a covering in our 
terminology. Now the other projection
$m : V \to M$ satisfies $\phi \circ m = n \circ \rho$. 
This means that $\phi(m) = n$ in $\Ga(V, N)$.   

Conversely, let us take $U := N$ and 
$n := \opn{id}_N \in \Ga(U, N)$. There exists an epimorphism 
$\rho : V \surj U$ and a morphism $m : V \to M$ such that 
$\phi \circ m = n \circ \rho = \rho$.
We see that $\phi \circ m$ is an epimorphism, and hence $\phi$ is an 
epimorphism.
\end{proof}

\begin{exa} \label{exa:4090}
Suppose 
\[ \bsym{S} = \bigl(  0 \to M' \xar{\phi} M \xar{\psi} M'' \to 0 \bigr) \]
is a sequence in $\cat{M}$. Here is how exactness of $\bsym{S}$ is tested using 
the first sheaf trick. 

By item (1) of the first sheaf trick, 
the condition $\psi \circ \phi = 0$ is same as 
$\Ga(U, \psi \circ \phi) = 0$ for every $U \in \cat{M}$. 

Exactness at $M'$ is the same as this condition: for every $U \in \cat{M}$ and 
every section $m' \in \Ga(U, M')$, if $\phi(m') = 0$ then $m' = 0$. 
This is by item (2). 

Exactness at $M''$ is the same as this condition: 
for every $U \in \cat{M}$ and every section 
$m'' \in \Ga(U, M'')$, there exists a covering $V \surj U$ and a section 
$m \in \Ga(V, M)$, such that 
$\psi(m) = m''$ in $\Ga(V, M'')$.

Finally, exactness at $M$ is the same as this condition (given that 
$\psi \circ \phi = 0$): for every $U \in \cat{M}$ and every section 
$m \in \Ga(U, M)$ such that $\psi(m) = 0$, there exists a covering 
$V \surj U$ and a section $m' \in \Ga(V, M')$, such that 
$\phi(m') = m$ in $\Ga(V, M'')$.
\end{exa}

\begin{prop}[Second Sheaf Trick] \label{prop:3636}
\index{Sheaf Tricks (in abelian categories)|)}
Let $M$ be an object in an abelian category $\cat{M}$. Suppose 
$\rho_1 : V_1 \to U$ and $\rho_2 : V_2 \to U$ are morphisms in $\cat{M}$, 
such that 
$(\rho_1, \rho_2) : V_1 \oplus V_2 \to U$
is an epimorphism. Let $W := V_1 \times_U V_2$, with projection morphisms 
$\si_1 : W \to V_1$ and $\si_2 : W \to V_2$. 
Then the sequence 
\[ 0 \to \Ga(U, M) \xar{(\rho_1^*, \rho_2^*)} \Ga(V_1, M) \times \Ga(V_2, M)
\xar{(\si_1^*, - \si_2^*)} \Ga(W, M) \]
in $\cat{Mod} \K$ is exact. 
\end{prop}

\begin{proof}
Note that $V_1 \times V_2 = V_1 \oplus V_2$. 
By Lemma \ref{lem:3630}(1) there is an exact sequence 
\[ 0 \to W \xar{(\si_1, -\si_2)} V_1 \oplus V_2 \xar{(\rho_1, \rho_2)} U 
\to 0 \]
in $\cat{M}$. By Proposition \ref{prop:3637}(2) we get an exact sequence
\[ 0 \to \Ga(U, M) \xar{(\rho_1^*, \rho_2^*)} \Ga(V_1 \oplus V_2, M)
\xar{(\si_1, -\si_2)^*} \Ga(W, M) \]
in $\cat{Ab}$. But 
\[ \Ga(V_1 \oplus V_2, M) \cong \Ga(V_1, M) \times \Ga(V_2, M) . \qedhere \]
\end{proof}

\begin{rem} \label{rem:3630}
As can be seen in the practical application of our sheaf tricks (e.g.\ 
in the proofs of Proposition \ref{prop:3620} and Theorem \ref{thm:3640})
these tricks reduce proofs about an abstract abelian category $\cat{M}$, to 
proofs that are just like in the concrete case of the  abelian category
$\cat{M} = \cat{Ab} X$ of sheaves of abelian groups on a topological 
space $X$. This is not as easy as the case 
$\cat{M} = \cat{Mod} A$ of modules over a ring $A$, that the Freyd-Mitchell 
Theorem permits, but -- at least for someone with experience in algebraic 
geometry -- our method is quite straightforward. 

Our sheaf tricks are inspired by \cite[tag 05PL]{SP}. There it is shown 
that the geometric allegory is genuine, except that instead of a topological 
space $X$, there is a {\em site} $X$ (in the sense of Grothendieck). The 
underlying category of the site $X$ is the abelian category $\cat{M}$ itself, 
and the coverings of $X$ are the epimorphisms in $\cat{M}$. It is proved in 
\cite{SP} that the category $\cat{M}$ embeds as a full abelian subcategory of 
$\cat{Ab} X$, the category of sheaves of abelian groups on $X$. 

Proposition \ref{prop:3635} is very similar to
\cite[Theorem VIII.3]{Mac2}, where what 
we call ``sections'' are called ``members''. But S. MacLane's method loses the 
abelian group structure (the members are not elements of abelian groups). 
An improvement of MacLane's method can be found in an unpublished note of 
G. Bergman \cite{Ber}. 
\end{rem}




\mysubsection{Additive Functors} \label{subsec:add-funcs}

\begin{dfn} \label{dfn:1031}
Let $\cat{M}$ and $\cat{N}$ be $\K$-linear categories. A functor 
$F : \cat{M} \to \cat{N}$ is called a {\em $\K$-linear functor}%
\index{Linear functor}%
\index{Linear functor!K@$\K$-linear functor}
if for every $M_0, M_1 \in \cat{M}$ the function 
\[ F : \opn{Hom}_{\cat{M}}(M_0, M_1) \to
\opn{Hom}_{\cat{N}} \bigl( F(M_0), F(M_1) \bigr) \]
is a $\K$-linear homomorphism. 

When the base ring $\K$ is implicit (cf.\ Convention \ref{conv:2490}), 
we sometimes say that $F$ is an {\em additive functor}, with the same meaning 
as $\K$-linear functor. 
\end{dfn}

Additive functors commute with finite direct sums. More precisely:

\begin{prop} \label{prop:2}
Let $F : \cat{M} \to \cat{N}$ be an additive functor between linear categories,
let $\{ M_i \}_{i \in I}$ be a finite collection of objects of $\cat{M}$, and
assume that  the direct sum $(M, \{ e_i \}_{i \in I})$ of the collection 
$\{ M_i \}_{i \in I}$ exists in $\cat{M}$. Then 
$\big( F(M), \{ F(e_i) \}_{i \in I} \big)$ is a direct sum of the collection 
$\{ F(M_i) \}_{i \in I}$ in $\cat{N}$.  
\end{prop}

\begin{exer} \label{exer:1825}
Prove Proposition \ref{prop:2}. (Hint: use Proposition \ref{prop:1}.)
\end{exer}

Note that the proposition above also talks about finite products, because of 
Proposition \ref{prop:1}.

\begin{prop} \label{prop:3600}
Suppose 
$F, G : \cat{K} \to \cat{L}$
are additive functors between linear categories, and 
$\eta : F \to G$ is a morphism of functors. 
Let $M, M' , N$ be objects of $\cat{K}$, and assume that 
$N \cong M \oplus M'$. Then the following two conditions are equivalent. 
\begin{enumerate}
\rmitem{i} $\eta_N : F(N) \to G(N)$ is an isomorphism. 

\rmitem{ii} $\eta_M : F(M) \to G(M)$ and 
$\eta_{M'} : F(M') \to G(M')$ are isomorphisms. 
\end{enumerate}
\end{prop}

\begin{exer} \label{exer:3601}
Prove Proposition \ref{prop:3600}.
\end{exer}

\begin{exa} \label{exa:4935}
Let $f : A \to B$ be a ring homomorphism. There are two additive functors 
associated to $f$~: the 
{\em restriction functor}%
\index{Restriction functor}
$\opn{Rest}_f : \cat{Mod} B \to  \cat{Mod} A$
and the {\em induction functor}\index{Induction functor}
$\opn{Ind}_f : \cat{Mod} A \to  \cat{Mod} B$.
Given a $B$-module $N$, the $A$-module $\opn{Rest}_f(N)$ has the 
same underlying $\K$-module, and $A$ acts on it through $f$.
For an $A$-module $M$, the induced $B$-module is 
$\opn{Ind}_f(M) := B \ot_A M$.
\end{exa}

\begin{prop} \label{prop:109}
Let $F : \cat{M} \to \cat{N}$ be an additive functor between linear
categories. Then\tup{:}
\begin{enumerate}
\item For every $M \in \cat{M}$ the function 
$F : \opn{End}_{\cat{M}}(M) \to \opn{End}_{\cat{N}}(F(M))$
is a ring homomorphism.

\item For every $M_0, M_1\in \cat{M}$ the function 
\[ F : \opn{Hom}_{\cat{M}}(M_0, M_1) \to 
\opn{Hom}_{\cat{N}} \bigl( F(M_0), F(M_1) \bigr) \]
is a homomorphism of left $\opn{End}_{\cat{M}}(M_1)$-modules, and of 
right $\opn{End}_{\cat{M}}(M_0)$-modules.

\item If $M$ is a zero object of $\cat{M}$, then $F(M)$ is a 
zero object of $\cat{N}$.
\end{enumerate}
\end{prop}

\begin{proof} \mbox{}

\smallskip \noindent
(1) By Definition \ref{dfn:1031} the function $F$ respects addition. By the 
definition of a functor, it respects multiplication and units. 

\medskip \noindent 
(2) Immediate from the definitions, like (1). 

\medskip \noindent 
(3) Combine part (1) with Proposition \ref{prop:1030}.
\end{proof}

\begin{dfn}
\index{Exact functor}
Let $F : \cat{M} \to \cat{N}$ be an additive functor between abelian
categories. 
\begin{enumerate}
\item $F$ is called {\em left exact} if it commutes with kernels. 
Namely for every morphism $\phi : M_0 \to M_1$ in $\cat{M}$, with 
kernel $k : K \to M_0$, the morphism $F(k) : F(K) \to F(M_0)$ is 
a kernel of $F(\phi) : F(M_0) \to F(M_1)$.

\item $F$ is called {\em right exact} if it commutes with cokernels. 
Namely for every morphism $\phi : M_0 \to M_1$ in $\cat{M}$, with 
cokernel $c : M_1 \to C$, the morphism $F(c) : F(M_1) \to F(C)$ is 
a cokernel of $F(\phi) : F(M_0) \to F(M_1)$.

\item $F$ is called {\em exact} if it is both left exact and right exact.
\end{enumerate}
\end{dfn}

Let us illustrate this. Suppose $\phi : M_0 \to M_1$
is a morphism in $\cat{M}$, with  kernel $K$ and cokernel $C$. Applying $F$ to
the sequence  
$K \xar{k} M_0 \xar{\phi} M_1 \xar{c} C$
in $\cat{M}$ 
we get the solid arrows in this diagram 
\[ \UseTips  \xymatrix @C=8ex @R=6ex {
F(K)
\ar[r]^{F(k)}
\ar@{-->}[dr]_{\psi}
& 
F(M_0)
\ar[r]^{F(\phi)}
&
F(M_1)
\ar@{-->}[d]
\ar[r]^{F(c)}
&
F(C)
\\
&
\opn{Ker}_{\cat{N}}(F(\phi))
\ar@{-->}[u]
&
\opn{Coker}_{\cat{N}}(F(\phi))
\ar@{-->}[ur]_{\chi}
} \]
in $\cat{N}$. 
Because $\cat{N}$ is abelian, we get the vertical dashed arrows: the kernel 
and cokernel of $F(\phi)$. 
The slanted dashed arrows exist and are unique because 
$F(\phi) \circ F(k) = 0$ and $F(c) \circ F(\phi) = 0$. 
Left exactness of $F$ requires $\psi$ to be an isomorphism, and 
right exactness requires $\chi$ to be an isomorphism.

Recall that a short exact sequence in $\cat{M}$ is an exact sequence of the 
form 
\begin{equation} \label{eqn:4096}
\bsym{S} = \bigl( 0 \to M_0 \xar{\phi_0} M_1 \xar{\phi_1} M_2 \to 0 \bigr) .
\end{equation}

\begin{prop}  \label{prop:1826}
Let $F : \cat{M} \to \cat{N}$ be an additive functor between abelian
categories. 
\begin{enumerate}
\item The functor $F$ is left exact if and only if 
for every short exact sequence $\bsym{S}$
in $\cat{M}$, with notation \tup{(\ref{eqn:4096})}, the sequence 
\[ 0 \to F(M_0) \xar{F(\phi_0)} F(M_1) \xar{F(\phi_0)} F(M_2)  \]
is exact in $\cat{N}$.

\item The functor $F$ is right exact if and only if 
for every short exact sequence $\bsym{S}$ in $\cat{M}$, with the notation 
with notation \tup{(\ref{eqn:4096})}, the sequence 
\[ F(M_0) \xar{F(\phi_0)} F(M_1) \xar{F(\phi_1)} F(M_2) \to 0 \]
is exact in $\cat{N}$.
\end{enumerate}
\end{prop}

\begin{exer} \label{exer:1826}
Prove Proposition \ref{prop:1826}. (Hint: $M_0 \cong \opn{Ker}(M_1 \to M_2)$ 
etc.)
\end{exer}

\begin{exa} \label{exa:4680}
Let $A$ be a commutative ring, and let $M$ be a fixed $A$-module. 
Define functors 
$F, G : \cat{Mod} A \to \cat{Mod} A$ like this:
$F(N) := M \ot_A N$ and $G(N) := \opn{Hom}_A(M, N)$.
Then $F$ is right exact and $G$ is left exact.
\end{exa}

\begin{exa} \label{exa:4681}
Let $f : A \to B$ be a ring homomorphism. 
In the notation of Example \ref{exa:4935}, the restriction functor 
$\opn{Rest}_f$ is exact, and the induction functor
$\opn{Ind}_f$ is right exact. 
\end{exa}

\begin{prop}
Let $F : \cat{M} \to \cat{N}$ be an additive functor between abelian
categories. If $F$ is an equivalence then it is exact.
\end{prop}

\begin{proof}
We will prove that $F$ respects kernels; the proof for cokernels is similar.
Take a morphism $\phi : M_0 \to M_1$ in $\cat{M}$, with kernel $K$. 
We have this diagram (solid arrows):
\[ \UseTips  \xymatrix @C=8ex @R=6ex {
M 
\ar@{-->}[d]_{\psi}
\ar@{-->}[dr]^{\theta}
\\
K
\ar[r]^{k}
& 
M_0
\ar[r]^{\phi}
&
M_1
} \]
Applying $F$ we obtain 
 this diagram (solid arrows):
\[ \UseTips  \xymatrix @C=8ex @R=6ex {
N = F(M)
\ar@{-->}[d]_{F(\psi)}
\ar@{-->}[dr]^{\bar{\theta}}
\\
F(K)
\ar[r]^{F(k)}
& 
F(M_0)
\ar[r]^{F(\phi)}
&
F(M_1)
} \]
in $\cat{N}$. 
Suppose $\bar{\theta} : N \to F(M_0)$ is a morphism in $\cat{N}$ s.t.\ 
$F(\phi) \circ \bar{\theta} = 0$.
Since $F$ is essentially surjective on objects, there is some $M \in \cat{M}$
with an isomorphism $\al : F(M) \iso N$. After replacing $N$ with $F(M)$ and 
$\bar{\theta}$ with $\bar{\theta} \circ \al$, we can assume that  
$N = F(M)$. 

Now since $F$ is fully faithful, there is a unique $\theta : M \to M_0$ s.t.\ 
$F(\theta) = \bar{\theta}$; and $\phi \circ \theta = 0$. 
So there is a unique $\psi : M \to K$ s.t.\  
$\theta = k \circ \psi$. It follows that 
$F(\psi) : F(M) \to F(K)$ is the unique morphism s.t.\ 
$\bar{\theta} = F(k) \circ F(\psi)$.
\end{proof}

Here is a result that could afford another proof of the previous proposition. 

\begin{prop} \label{prop:1827}
Let $F : \cat{M} \to \cat{N}$ be an additive functor between linear
categories. Assume $F$ is an equivalence, with quasi-inverse $G$. Then 
$G : \cat{N} \to \cat{M}$ is an additive functor.
\end{prop}

\begin{exer} \label{exer:1827}
Prove Proposition \ref{prop:1827}.
\end{exer}

\begin{dfn} \label{dfn:3620}
Consider abelian categories $\cat{M}$ and $\cat{N}$. 
Suppose we are given a sequence 
\[ \cdots F_{-1} \xar{\phi_{-1}} F_0 \xar{\phi_0} F_1 \xar{\phi_1} F_2 \cdots \]
(finite or infinite on either side), 
where each $F_i : \cat{M} \to \cat{N}$ is an additive functor, and each
$\phi_i : F_i \to F_{i + 1}$
is a morphism of functors. We say that this sequence is an {\em exact sequence 
of additive functors} if for every object $M \in \cat{M}$ the sequence 
\[ \cdots F_{-1}(M) \xar{\phi_{-1, M}} F_0(M) \xar{\phi_{0, M}} F_1(M) 
\xar{\phi_{1, M}} F_2(M) \cdots \]
in $\cat{N}$\ is exact.
\end{dfn}

\begin{prop} \label{prop:3620}
Let $\cat{M}$ and $\cat{N}$ be abelian categories, and let 
$F_0 \xar{\phi_0} F_1 \xar{\phi_1} F_2$
be a sequence of additive functors $\cat{M} \to \cat{N}$.
\begin{enumerate}
\item If  
$F_0 \xar{\phi_0} F_1 \xar{\phi_1} F_2 \to 0$
is an exact sequence of additive functors, and if the functors $F_0$ and $F_1$ 
are both right exact, then the functor $F_2$ is right exact. 

\item If  
$0 \to F_0 \xar{\phi_0} F_1 \xar{\phi_1} F_2$
is an exact sequence of additive functors, and if the functors $F_1$ and $F_2$ 
are both left exact, then the functor $F_0$ is left exact.
\end{enumerate}
\end{prop}

\begin{proof} \mbox{}

\smallskip \noindent
(1) Let 
$M' \xar{\si} M \xar{\tau} M'' \to 0$ 
be an exact sequence in $\cat{M}$. We must show that 
$F_2(M') \to F_2(M) \to F_2(M'') \to 0$ 
is an exact sequence in $\cat{N}$. Let us examine the commutative diagram 
\begin{equation} \label{eqn:3628}
\UseTips \xymatrix @C=8ex @R=6ex {
F_0(M')
\ar[r]^{\phi_{0, M'}}
\ar[d]^{F_0(\si)}
&
F_1(M')
\ar[d]^{F_1(\si)}
\ar[r]^{\phi_{1, M'}}
&
F_2(M')
\ar[r]
\ar[d]^{F_2(\si)}
&
0
\\
F_0(M)
\ar[r]^{\phi_{0, M}}
\ar[d]^{F_0(\tau)}
&
F_1(M)
\ar[r]^{\phi_{1, M}}
\ar[d]^{F_1(\tau)}
&
F_2(M)
\ar[r]
\ar[d]^{F_2(\tau)}
&
0
\\
F_0(M'')
\ar[r]^{\phi_{0, M''}}
\ar[d]
&
F_1(M'')
\ar[r]^{\phi_{1, M''}}
\ar[d]
&
F_2(M'')
\ar[r]
\ar[d]
&
0
\\
0
&
0
&
0
}
\end{equation}
in $\cat{N}$. It is known that the rows and the first two columns are exact. We 
must prove that the third column is exact. 

First let us prove that $F_2(\tau)$ is an epimorphism. This is easy: 
we know that $F_1(\tau)$ and $\phi_{1, M''}$ are epimorphisms; hence 
$\phi_{1, M''} \circ F_1(\tau) = F_2(\tau) \circ \phi_{1, M}$
is an epimorphism; and thus $F_2(\tau)$ is an epimorphism.

More challenging is the proof that the third column is exact at $F_2(M)$, 
namely that 
$\opn{Im}(F_2(\si)) \to \opn{Ker}(F_2(\tau))$
is an epimorphism. For this we shall use the first sheaf trick (Proposition 
\ref{prop:3635}) and a diagram chase. Consider a section 
$m_2 \in \Ga(U, \opn{Ker}(F_2(\tau)))$
on some ``open set'' $U$, i.e.\ for some object $U \in \cat{M}$. 
We shall prove that there exist a covering $V \surj U$ and a section 
$m'_2 \in \Ga(V, F_2(M'))$ such that 
$F_2(\si)(m'_2) = m_2$ 
in $\Ga(V, F_2(M))$. 

Because $\phi_{1, M}$ is an epimorphism, there is a covering 
$U_1 \surj U$ and a section 
$m_1 \in \Ga(U_1, F_1(M))$
such that 
$\phi_{1, M}(m_1) = m_2$ in $\Ga(U_1, F_2(M))$. 
Let
$m''_1 := F_1(\tau)(m_1)  \in \Ga(U_1, F_1(M''))$.
We have 
$\phi_{1,M''}(m''_1) = F_2(\tau)(m_2) = 0$. 
This means that 
$m''_1 \in \Ga(U_1, \opn{Ker}(\phi_{1,M''}))$.

The exactness of the third row says that for some covering 
$U_2 \surj U_1$ there is a section  
$m''_0 \in \Ga(U_2, F_0(M''))$
such that 
$\phi_{0, M''}(m''_0 ) = m''_1$ in $\Ga(U_2, F_1(M''))$.

The exactness of the first column implies that for some covering 
$U_3 \surj U_2$ there is a section  
$m_0 \in \Ga(U_3, F_0(M))$
such that 
$m''_0 = F_0(\tau)(m_0)$ in $\Ga(U_3, F_0(M''))$.
Define 
$\til{m}_1 := m_1 - \phi_{0, M}(m_0)$ in $\Ga(U_3, F_1(M))$.
Note that 
$\phi_{1, M}(\til{m}_1) = \phi_{1, M}(m_1) = m_2$ in $\Ga(U_3, F_2(M))$.
Also $F_1(\tau)(\til{m}_1) = 0$, i.e.\ 
$\til{m}_1 \in \Ga(U_3, \opn{Ker}(F_1(\tau)))$.

Due to the exactness of the second column, for some covering 
$V \surj U_3$ there is a section 
$m'_1 \in \Ga(V, F_1(M'))$
such that 
$\til{m}_1 = F_1(\si)(m'_1)$ in $\Ga(V, F_1(M))$.
Define
$m'_2 := \phi_{1, M'}(m'_1) \in \Ga(V, F_2(M'))$.
Then 
$F_2(\si)(m'_2) = \phi_{1, M}(\til{m}_1) = m_2$ in $\Ga(V, F_2(M))$.

\medskip \noindent 
(2) See next exercise.
\end{proof}

\begin{exer} \label{exer:3625}
Prove part (2) of Proposition \ref{dfn:3620}. (Hint: imitate the proof of part 
(1); but this is easier.)
\end{exer}

We end this subsection with a discussion of additive contravariant functors. 
Suppose $\cat{M}$ and $\cat{N}$ are linear categories. A contravariant functor 
$F : \cat{M} \to \cat{N}$ is said to be additive if it satisfies the condition 
in Definition \ref{dfn:1031}, with the obvious changes. 

\begin{prop} \label{prop:2491}
Let $\cat{M}$ and $\cat{N}$ be linear categories. Put on $\cat{M}^{\mrm{op}}$ 
the canonical linear structure \tup{(}see Proposition \tup{\ref{prop:2490})}.

\begin{enumerate}
\item The functor $\opn{Op} : \cat{M} \to \cat{M}^{\mrm{op}}$
is an additive contravariant functor. 

\item If $F : \cat{M} \to \cat{N}$ is an additive contravariant functor, then 
$F \circ \opn{Op} : \cat{M}^{\mrm{op}} \to \lb \cat{N}$
is an additive functor; and vice versa. 
\end{enumerate}
\end{prop}

\begin{exer} 
Prove Proposition \ref{prop:2491}.
\end{exer}

In view of Proposition \ref{prop:1826}, we can give an unambiguous definition 
of left and right exact contravariant functors. 

\begin{dfn} \label{dfn:4096}
Let $F : \cat{M} \to \cat{N}$ be an additive contravariant functor between 
abelian categories. 
\begin{enumerate}
\item $F$ a {\em left exact contravariant functor} if for 
every short exact sequence $\bsym{S}$ in $\cat{M}$, in the notation of 
(\ref{eqn:4096}), the sequence 
\[ 0 \to F(M_2) \xar{F(\phi_1)} F(M_1) \xar{F(\phi_0)} F(M_0) \]
in $\cat{N}$ is exact. 

\item $F$ is a {\em right exact contravariant functor} 
if for every short exact sequence $\bsym{S}$ in $\cat{M}$, the sequence 
\[ F(M_2) \xar{F(\phi_1)} F(M_1) \xar{F(\phi_0)} F(M_0) \to 0 \]
in $\cat{N}$ is exact.

\item $F$ is an {\em exact contravariant functor} if it sends every short 
exact sequence $\bsym{S}$ in $\cat{M}$ to a short exact sequence in $\cat{N}$. 
\end{enumerate}
\end{dfn}

\begin{prop} \label{prop:2492} 
Let $\cat{M}$ and $\cat{N}$ be abelian categories. 
Recall that $\cat{M}^{\mrm{op}}$ is also an abelian category.

\begin{enumerate}
\item The functor $\opn{Op} : \cat{M} \to \cat{M}^{\mrm{op}}$
is an exact contravariant functor. 

\item If $F : \cat{M} \to \cat{N}$ is an exact contravariant functor, then 
$F \circ \opn{Op} : \cat{M}^{\mrm{op}} \to  \cat{N}$
is an exact functor; and vice versa. Likewise for left exactness and right 
exactness. 
\end{enumerate}
\end{prop}

\begin{exer} 
Prove Proposition \ref{prop:2492}.
\end{exer}

\begin{exa} \label{exa:4682}
Let $A$ be a commutative ring, and let $M$ be a fixed $A$-module. 
Define the contravariant functor $F : \cat{Mod} A \to \cat{Mod} A$ to be 
$F(N) := \opn{Hom}_A(N, M)$.
Then $F$ is a left exact contravariant function.
\end{exa}


Sometimes $\cat{M}$ and $\cat{M}^{\mrm{op}}$ are equivalent as abelian 
categories, as the next exercise shows. For a counterexample see Remark 
\ref{rem:2510} below. 

\begin{exer} \label{exer:2515}
Let $\K$ be a field, and consider the category 
$\cat{M} := \cat{Mod}_{\mrm{f}} \K$
of finitely generated $\K$-modules (traditionally known as ``finite dimensional 
vector spaces over $\K$''). This is a $\K$-linear abelian category. 
Find a $\K$-linear equivalence
$F : \cat{M}^{\mrm{op}} \to \cat{M}$.
\end{exer}



\mysubsection{Projective Objects}
In this subsection $\cat{M}$ is an abelian category.

A {\em splitting} of an epimorphism 
$\psi : M \to M''$ in $\cat{M}$ is a morphism
$\al : M'' \to M$ s.t.\ $\psi \circ \al = \opn{id}_{M''}$.
A splitting of a monomorphism 
$\phi : M' \to M$ is  a morphism $\be : M \to M'$ s.t.\ 
$\be \circ \phi = \opn{id}_{M'}$.
A splitting of a short exact sequence 
\begin{equation} \label{eqn:1035}
0 \to M' \xar{\phi} M \xar{\psi} M'' \to 0
\end{equation}
is a splitting of the epimorphism $\psi$, or equivalently a splitting of the
mono\-morphism $\phi$. The short exact sequence is said to be {\em split} if it
has some splitting. 

\begin{exer} \label{exe:1035}
Show how to get from a splitting of $\phi$  to a splitting of $\psi$, and vice 
versa. Show how any of those gives rise to an isomorphism 
$M \cong M' \oplus M''$. 
\end{exer}

\begin{dfn}
An object $P \in \cat{M}$ is called a 
{\em projective object}%
\index{Abelian category! projective object in}
if for every morphism $\ga : P \to N$ and every {\em epimorphism}
$\psi : M \surj N$, there exists a morphism 
$\til{\ga} : P \to N$ such that 
$\psi \circ \til{\ga} = \ga$. 
\end{dfn}

This is described in the following commutative diagram in $\cat{M}$~: 
\[ \UseTips  \xymatrix @C=8ex @R=6ex {
&
P
\ar[d]^{\ga}
\ar@{-->}[dl]_{\til{\ga}}
\\
M 
\ar@{->>}[r]_{\psi}
&
N
} \]


\begin{prop} \label{prop:1670}
The following conditions are equivalent for $P \in \cat{M}$~\tup{:}
\begin{enumerate}
\rmitem{i} $P$ is projective.
\rmitem{ii} The additive functor 
$\opn{Hom}_{\cat{M}}(P, -) : \cat{M} \to \cat{Ab}$
is exact.
\rmitem{iii} Every short exact sequence \tup{(\ref{eqn:1035})} with 
$M'' = P$ is split.
\end{enumerate}
\end{prop}

\begin{proof}
Exercise.
\end{proof}

\begin{dfn}
We say $\cat{M}$ {\em has enough projectives} if every $M \in \cat{M}$ admits
an epimorphism $P \to M$ from a projective object $P$. 
\end{dfn}

\begin{exer} \label{exer:1670}
Let $A$ be a ring. 
\begin{enumerate}
\item Prove that an $A$-module $P$ is projective iff it is a 
direct summand of a free module; i.e.\ $P \oplus P' \cong Q$ for some module 
$P'$ and free module $Q$.

\item Prove that the category $\cat{Mod} A$ has enough projectives.
\end{enumerate} 
\end{exer}

\begin{exer} \label{exer:1675}
Let $\cat{M}$ be the category of finite abelian groups.
Prove that the only projective object in $\cat{M}$ is $0$. So $\cat{M}$ does 
not have enough projectives.
(Hint: use Proposition \ref{prop:1670}.)
\end{exer}

\begin{exa} \label{exa:2870}
Consider the scheme $X := \mbf{P}^1_{\K}$, the projective line over a field
$\K$, with structure sheaf $\mcal{O}_X$. 
The category $\cat{Coh} \mcal{O}_X$ of coherent $\mcal{O}_X$-modules is
abelian (it is a full abelian subcategory of $\cat{Mod} \mcal{O}_X$, cf.\ 
Example \ref{exa:1}). One can show that the only projective object of 
$\cat{Coh} \mcal{O}_X$ is $0$, but this is quite involved. 

Let us only indicate why the sheaf $\mcal{O}_X$ is not a projective
object of $\cat{Coh} \mcal{O}_X$. 
Denote by $t_0, t_1$ the homogeneous coordinates of $X$. These belong to 
$\Gamma(X, \mcal{O}_X(1))$, so each determines a homomorphism of sheaves
$t_j : \mcal{O}_X(i) \to \mcal{O}_X(i+1)$. We get a sequence 
\begin{equation} \label{eqn:5020}
0 \to \mcal{O}_X(-2) \xar{\sbmat{t_0 & t_1}} \mcal{O}_X(-1)^{\oplus 2} 
\xar{\sbmat{-t_1 \\ t_0}} \mcal{O}_X \to 0
\end{equation}
in $\cat{Coh} \mcal{O}_X$, which is known to be exact.
It is also known that 
$\Gamma(X, \mcal{O}_X) = \K$ and  
$\Gamma(X, \mcal{O}_X(-1)) = 0$. Therefore the sequence (\ref{eqn:5020}) is not 
split. 
\end{exa}


\mysubsection{Injective Objects}
In this subsection $\cat{M}$ is an abelian category.

\begin{dfn}
An object $I \in \cat{M}$ is called an 
{\em injective object}%
\index{Abelian category! injective object in}
if for every morphism $\ga : M \to I$ and every {\em monomorphism} 
$\psi : M \inj N$, there exists a morphism 
$\til{\ga} : N \to I$ such that 
$\til{\ga} \circ \psi = \ga$. 
\end{dfn}

This is depicted in the following commutative diagram in 
$\cat{M}$~:
\[ \UseTips  \xymatrix @C=8ex @R=6ex {
I
\\
M 
\ar[u]^{\ga}
\ar@{>->}[r]_{\psi}
&
N
\ar@{-->}[ul]_{\til{\ga}}
} \]


\begin{prop} \label{prop:2493}
The following conditions are equivalent for $I \in \cat{M}$\tup{:}
\begin{enumerate}
\rmitem{i} $I$ is injective.
\rmitem{ii} The additive functor 
$\opn{Hom}_{\cat{M}}(-, I) : \cat{M}^{\mrm{op}} \to \cat{Ab}$
is exact.

\rmitem{iii} Every short exact sequence \tup{(\ref{eqn:1035})} with 
$M' = I$ is split.
\end{enumerate}
\end{prop}

\begin{exer}
Prove  Proposition \ref{prop:2493}.
\end{exer}

Recall that $\opn{Op} : \cat{M} \to \cat{M}^{\mrm{op}}$ is an exact functor. 

\begin{prop} \label{prop:2494}
An object $J \in \cat{M}$ is injective if and only if the object 
$\opn{Op}(J) \in \cat{M}^{\mrm{op}}$ is projective. 
\end{prop}

\begin{exer}
Prove Proposition \ref{prop:2494}. 
\end{exer}

\begin{exa}
Let $A$ be a ring. Unlike projectives, the structure of injective objects in
$\cat{Mod} A$ is very complicated, and not much is known (except that they
exist). However if $A$ is a commutative noetherian ring then we know this:
every injective module $I$ is a direct sum of indecomposable injective 
modules; and the indecomposables are parameterized by $\opn{Spec}(A)$, the 
set of prime ideals of $A$. 
These facts are due to Matlis; see Subsection \ref{subsec:mr-inj-res} in the 
book. 
\end{exa}

\begin{dfn}
We say $\cat{M}$ {\em has enough injectives} if every $M \in \cat{M}$ admits
a monomorphism $M \to I$ to an injective object $I$. 
\end{dfn}

Here are a few results about injective objects. Recall that modules over a ring 
are always left modules by default. 

\begin{prop} \label{prop:1045}
Let $f : A \to B$ be a ring homomorphism, and let $I$ be an injective 
$A$-module. Then $J := \opn{Hom}_A(B, I)$ is an  injective $B$-module.
\end{prop}

\begin{proof}
Note that $B$ is a left $A$-module via $f$, and a right $B$-module. This makes
$J$ into a left $B$-module. In a formula: for $\phi \in J$ and $b, b' \in B$ we
have
$(b \cd \phi)(b') = \phi(b' \cd b)$. 

Now given any $N \in \cat{Mod} B$ there is an isomorphism 
\begin{equation} \label{eqn:4}
\opn{Hom}_B(N, J) =  \opn{Hom}_B(N, \opn{Hom}_A(B, I)) \cong 
\opn{Hom}_A(N, I) .
\end{equation}
 This is a natural isomorphism (of functors in $N$). So the functor 
$\opn{Hom}_B(-, J)$ is exact, and hence $J$ is injective. 
\end{proof}

\begin{thm}[Baer Criterion] \label{thm:151}
Let $A$ be a ring and $I$ an $A$-module. Assume that 
every $A$-module homomorphism $\a \to I$ from a left ideal $\a \subseteq A$
extends to a homomorphism $A \to I$. Then the module $I$ is injective.
\end{thm}

\begin{proof}
Consider an $A$-module $M$, a submodule $N \subseteq M$, and a homomorphism 
$\ga : N \to I$. We have to prove that $\ga$ extends to a homomorphism 
$M \to I$. Look at the pairs $(N', \ga')$ consisting of a submodule 
$N' \subseteq M$ that contains $N$, and a homomorphism 
$\ga' : N' \to I$ that extends $\ga$. The set of all such pairs is ordered by 
inclusion, and it satisfies the conditions of Zorn's Lemma. Therefore there 
exists a maximal pair $(N', \ga')$. We claim that $N' = M$. 

Otherwise, there is an element $m \in M$ that does not belong to $N'$. 
Define $N'' := N' + A \cd m$, so $N' \subsetneq N'' \sub M$. Let 
$\a := \{ a \in A \mid a \cd m \in N' \}$,
which is a left ideal of $A$. There is a short exact sequence 
\[ 0 \to \a \xar{\al} N' \oplus A \to N'' \to 0 \]
of $A$-modules, where 
$\al(a) := (a \cd m, -a)$. Let $\phi : \a \to I$ be the homomorphism
$\phi(a) := \ga'(a \cd m)$. 
By assumption, it extends to a homomorphism 
$\til{\phi} : A \to I$. We get a homomorphism 
$\ga' + \til{\phi}: N' \oplus A \to I$
that vanishes on the image of $\al$. 
Thus there is an induced homomorphism 
$\ga'' : N'' \to I$. This contradicts the maximality of $(N', \ga')$.
\end{proof}


\begin{lem} \label{lem:1045}
The $\Z$-module $\mbb{Q} / \Z$ is injective. 
\end{lem}

\begin{proof}
By the Baer criterion, it is enough to consider a homomorphism
$\ga : \a \to \mbb{Q} / \Z$ for an ideal $\a = n \cd \Z \subseteq \Z$. We may 
assume that $n \neq 0$. Say $\ga(n) = r + \Z$ with $r \in \mbb{Q}$. 
Then we can extend $\ga$ to $\til{\ga} : \Z \to \mbb{Q} / \Z$ with 
$\til{\ga}(1) := r / n + \Z$. 
\end{proof}

\begin{lem}  \label{lem:1046}
Let $\{ I_x \}_{x \in X}$ be a collection of injective objects of $\cat{M}$. If
the product $\prod_{x \in X} I_x$ exists in $\cat{M}$, then it is an 
injective object.
\end{lem}

\begin{proof}
Exercise.
\end{proof}

\begin{thm} \label{thm:1115}
Let $A$ be any ring. The category $\cat{Mod} A$ has enough injectives.
\end{thm}

\begin{proof} The proof is done in a few steps. 

\smallskip \noindent
Step 1. Here $A = \Z$. Take any nonzero $\Z$-module $M$ and any nonzero 
$m \in M$. Consider the cyclic submodule $M' := \Z \cd m \subseteq M$. 
There is a homomorphism $\ga' : M' \to \mbb{Q} / \Z$ s.t.\ 
$\ga'(m) \neq 0$. Indeed, if $M' \cong \Z$, then we take any 
$r \in \mbb{Q} - \Z$; and if $M' \cong \Z / (n)$ for some $n > 0$, then we 
take $r := 1 / n$. In either case, we define 
$\ga'(m) := r + \Z \in \mbb{Q} / \Z$.
Since $\mbb{Q} / \Z$ is an injective $\Z$-module, $\ga'$ extends 
to a homomorphism $\ga : M \to \mbb{Q} / \Z$.
By construction we have $\ga(m) \neq 0$. 

\medskip \noindent 
Step 2. Now $A$ is any ring, $M$ is any nonzero $A$-module, and $m \in M$ a 
nonzero element. Define the $A$-module 
$I := \opn{Hom}_{\Z}(A, \mbb{Q} / \Z)$, which, according to Lemma 
\ref{lem:1045} and Proposition \ref{prop:1045}, is an injective $A$-module. 
Let $\ga : M \to \mbb{Q} / \Z$ be a $\Z$-linear homomorphism such that 
$\ga(m) \neq 0$. Such $\ga$ exists by step 1. Let $\th : I \to \mbb{Q} / \Z$ 
be the $\Z$-linear homomorphism that sends an element $\chi \in I$ to 
$\chi(1) \in \mbb{Q} / \Z$.
The adjunction formula (\ref{eqn:4}) gives an $A$-module homomorphism 
$\psi : M \to I$ s.t.\ $\th \circ \psi = \ga$. We note that 
$(\th \circ \psi)(m) = \ga(m) \neq 0$, and hence $\psi(m) \neq 0$.  

\medskip \noindent 
Step 3. Here $A$ and $M$ are arbitrary. Let $I$ be as in step 2. For 
every nonzero  $m \in M$ there is an $A$-linear homomorphism $\psi_m : M \to I$ 
such that $\psi_m(m) \neq 0$.  For $m = 0$ let $\psi_0 : M \to I$ be an 
arbitrary homomorphism (e.g.\ $\psi_0 = 0$). 
Define the $A$-module $J := \prod_{m \in M} I$. There is a homomorphism 
$\psi := \prod_{m \in M} \psi_m : M \to J$,
and it is easy to check that $\psi$ 
is a monomorphism. By Lemma \ref{lem:1046}, $J$ is an injective $A$-module. 
\end{proof}

\begin{exer} \label{exer:1040}
At the price of getting a bigger injective module, we can make the 
construction of injective resolutions functorial. Let
$I := \opn{Hom}_{\Z}(A, \mbb{Q} / \Z)$ as above. Given an $A$-module $M$, 
consider the set 
$X(M) := \opn{Hom}_A(M, I) \cong \opn{Hom}_{\Z}(M, \mbb{Q} / \Z)$.
Let 
$J(M) := \prod_{\psi \in X(M)} I$. There is a ``tautological'' homomorphism 
$\phi_M : M \to J(M)$. Show that $\phi_M$ is a monomorphism, 
$J : M \mapsto J(M)$ is a functor, and $\phi : \opn{Id} \to J$ is a natural 
transformation. 

Is the functor $J : \cat{Mod} A \to \cat{Mod} A$ additive?
\end{exer}

\begin{exa}
Let $\cat{N}$ be the category of torsion abelian groups, and $\cat{M}$ the
category  of finite abelian groups. Then 
$\cat{N} \subseteq \cat{Ab}$
and $\cat{M} \subseteq \cat{N}$ are
full abelian subcategories. $\cat{M}$ has no projectives nor injectives except
$0$ (see Exercise \ref{exer:1675} regarding projectives). The only projective 
in $\cat{N}$ is $0$. However, it can be shown that $\cat{N}$ has enough 
injectives; see \cite[Lemma III.3.2]{Har} or \cite[Proposition 4.6]{Ye1}.
\end{exa}

\begin{prop} \label{prop:1046}
If $A$ is a left noetherian ring, then every direct sum of injective $A$-modules
is an injective module. 
\end{prop}

\begin{exer} \label{exer:1829}
Prove Proposition \ref{prop:1046}.
(Hint: use the Baer criterion.)
\end{exer}

\begin{exer} \label{exer:1045}
Here we study injectives in the category $\cat{Ab} = \cat{Mod} \Z$. 
By Lemma \ref{lem:1045}, the module $I := \Q / \Z$ is injective. For a 
(positive) prime number $p$, we denote by $\what{\Z}_p$ the ring of $p$-adic 
integers, and by $\what{\Q}_p$ its field of fractions (namely the $p$-adic 
completions of $\Z$ and $\Q$ respectively). Define the abelian group
$I_p := \what{\Q}_p / \what{\Z}_p$. 
\begin{enumerate}
\item Show that $I_p$ is an injective object of $\cat{Ab}$. 
\item Show that $I_p$ is indecomposable (i.e.\ it is not the direct sum of 
two nonzero objects). 
\item Show that $I \cong \bigoplus_p I_p$.
\item The theory (see Subsection \ref{subsec:mr-inj-res}) tells us that there 
is another indecomposable injective 
object in $\cat{Ab}$, besides the $I_p$. Try to identify it.
\end{enumerate}
\end{exer}

The abelian category $\cat{Mod} \mcal{A}$ associated to a ringed space $(X, 
\mcal{A})$ was introduced in Example \ref{exa:1}.

\begin{prop} \label{prop:2875}
Let $(X, \mcal{A})$ be a ringed space. The category $\cat{Mod} \mcal{A}$
has enough injectives. 
\end{prop}

\begin{proof}
Let $\mcal{M}$ be an $\mcal{A}$-module. Take a point $x \in X$. The stalk 
$\mcal{M}_x$ is a module over the ring $\mcal{A}_x$, and by Theorem 
\ref{thm:1115} we can
find an embedding $\phi_x : \mcal{M}_x \to I_x$ into an injective
$\mcal{A}_x$-module. 
Let $g_x : \{ x \} \to X$ be the inclusion, which we may view as a map of
ringed spaces from $( \{ x \}, \mcal{A}_x)$ to $(X, \mcal{A})$.
Define $\mcal{I}_x := {g_x}_* (I_x)$, which is an $\mcal{A}$-module (in fact it 
is a constant sheaf supported on the closed set $\ol{  \{ x \} } \subseteq X$). 
The adjunction formula gives rise to a sheaf homomorphism 
$\psi_x : \mcal{M} \to  \mcal{I}_x$. 
Since the functor $g_x^* : \cat{Mod} \mcal{A} \to \cat{Mod} \mcal{A}_x$
is exact, the adjunction formula shows that $\mcal{I}_x$ is an injective
object of $\cat{Mod} \mcal{A}$.

Finally let $\mcal{J} := \prod_{x \in X} \mcal{I}_x$. This is an injective
$\mcal{A}$-module. There is a homomorphism 
$\psi :=  \prod_{x \in X} \psi_x : \mcal{M} \to \mcal{J}$
in $\cat{Mod} \mcal{A}$. This is a monomorphism, since for every point $x$, 
letting $\mcal{J}_x$ be the stalk of the sheaf $\mcal{J}$ at $x$, the 
composition $\mcal{M}_x \xar{\psi_x} \mcal{J}_x \xar{p_x} \mcal{I}_x$
is the embedding $\phi_x : \mcal{M}_x \to I_x$.
\end{proof}

\begin{rem} \label{rem:2510}
Let $A$ be a nonzero ring, and consider the abelian category 
$\cat{M} := \cat{Mod} A$, the category of $A$-modules. A reasonable question to 
ask is this: Are the abelian categories $\cat{M}$ and $\cat{M}^{\mrm{op}}$ 
equivalent? The answer is negative. 
In fact, P. Freyd, in \cite[Exercise 5.B.3]{Fre}, shows that 
$(\cat{Mod} A)^{\mrm{op}}$ is not equivalent, as an abelian category, to 
$\cat{Mod} B$ for any ring $B$. The argument involves a delicate study of 
countable coproducts and products, and properties of Grothendieck abelian 
categories.  
\end{rem}

Here is a special case of the previous remark, that might shed some more light 
on the issue. 

\begin{exa} \label{exa:2990}
Consider the category $\cat{Mod} \Z = \cat{Ab}$ of abelian groups. 
Here is a proof that there does not exist an 
additive equivalence 
$F : \cat{Ab}^{\mrm{op}} \to \cat{Ab}$.
Suppose we had such an equivalence. Consider the object 
$P := \Z \in \cat{Ab}$, and let $I := F(P) \in \cat{Ab}$.
Because $P$ is an indecomposable projective object, and 
$F : \cat{Ab} \to \cat{Ab}$ is a contravariant equivalence,
the object $I$ has to be an indecomposable injective. 
The endomorphism rings are 
$\opn{End}_{\cat{Ab}}(I) \cong \opn{End}_{\cat{Ab}}(P)^{\mrm{op}} = 
\Z^{\mrm{op}} = \Z$.
However, the structure theorem for injective modules over commutative 
noetherian rings (Theorem \ref{thm:2026}) says that the only indecomposable 
injectives in $\cat{Ab}$ are $I = \what{\Q}_p / \what{\Z}_p$ and 
$I = \Q$; and their endomorphism rings are $\what{\Z}_p$ and $\Q$ respectively.
\end{exa}








\cleardoublepage
\mysection{Differential Graded Algebra} \label{sec:DG-algebra}

\AYcopyright

Recall that according to Convention \ref{conv:2490} 
there is a nonzero commutative base ring $\K$. By default all  
rings are $\K$-central, all linear categories are $\K$-linear,
all linear operations (such as ring homomorphisms and linear functors)
are $\K$-linear, and $\ot$ means $\ot_{\K}$. 

Throughout ``DG'' stands for ``differential graded''. 
There is some material about DG algebra in a few published references, such as
the book \cite{Mac1} and the papers \cite{Kel1}, \cite{To1}, \cite{Sao} and 
\cite{Avr}. However, for our 
purposes we need a much more detailed understanding of this theory, and this is 
what the present section provides. 


\mysubsection{Graded Algebra} \label{subsec:gr-alg}

Before entering the DG world, it is good to understand the graded world. 

\begin{dfn} \label{dfn:4101}
A {\em cohomologically graded $\K$-module}%
\index{Cohomologically graded module}
is a $\K$-module $M$ equipped with a direct sum decomposition
$M = \bigoplus_{i \in \Z} M^i$
into $\K$-submodules. The $\K$-module $M^i$ is called the {\em homogeneous 
component of cohomological degree $i$} of $M$. The nonzero elements of $M^i$ 
are called {\em homogeneous elements of cohomological degree $i$}. 
\end{dfn}

From here until Section \ref{sec:alg-gra-rings} we are going to respect the 
following convention, which will simplify the discussion. 

\begin{conv} \label{conv:4125}
By ``graded $\K$-module'' we mean a {\em cohomologically graded $\K$-module}, 
as defined above.
\end{conv}

In Section \ref{sec:alg-gra-rings} we shall introduce {\em algebraically graded 
rings and modules}, and then we shall have to make a careful distinction 
between these notions. See Remark \ref{rem:2516} regarding commutativity in the 
two settings. 

Suppose $M$ and $N$ are graded $\K$-modules. For an integer $i$ let 
\[ (M \ot N)^i := \bigoplus_{j \in \Z} \, (M^j \ot N^{i - j}) . \]
Then 
\begin{equation} \label{eqn:1691}
M \ot N = \bigoplus_{i \in \Z} \, (M \ot N)^i  
\end{equation}
is a graded $\K$-module. 

A $\K$-linear homomorphism $\phi : M \to N$ is said to be {\em homogeneous of  
degree $i$} if $\phi(M^j) \sub N^{j + i}$ for all $j$. 
We denote by $\opn{Hom}_{\K}(M, N)^i$ the $\K$-module of degree $i$ 
homomorphisms $M \to N$. In other words 
\begin{equation} \label{eqn:4120}
\opn{Hom}_{\K}(M, N)^i := \prod_{j \in \Z} \, 
\opn{Hom}_{\K}(M^j, N^{j + i}) .
\end{equation}

\begin{dfn} \label{dfn:4120}
Let $M$ and $N$ be graded $\K$-modules.
\begin{enumerate}
\item The {\em module of graded $\K$-linear homomorphisms} from $M$ to $N$ is 
the graded $\K$-module
\[ \opn{Hom}_{\K}(M, N) := \bigoplus_{i \in \Z} \,
\opn{Hom}_{\K}(M, N)^i  . \]

\item A  degree $0$ homomorphism $\phi : M \to N$ is called a 
{\em strict homomorphism of graded $\K$-modules}.%
\index{Strict homomorphism!of graded modules} 
\end{enumerate}
\end{dfn}

If $M_0, M_1, M_2$ are graded $\K$-modules, and 
$\phi_k : M_{k} \to M_{k + 1}$ are $\K$-linear homomorphisms of degrees 
$i_k$, then $\phi_1 \circ \phi_0 : M_0 \to M_2$
is a $\K$-linear homomorphism of degree $i_0 + i_1$. 
The identity automorphism $\opn{id}_M : M \to M$ has degree $0$. 

\begin{dfn} \label{dfn:4100}
The {\em strict category of graded $\K$-modules}%
\index{Strict category! of graded modules} 
is the category $\dcat{G}_{\mrm{str}}(\K)$,
whose objects are the cohomologically graded $\K$-modules, and whose morphisms 
are the strict homomorphisms of cohomologically graded $\K$-mod\-ules.
\end{dfn}

It is easy to see that $\dcat{G}_{\mrm{str}}(\K)$ is a $\K$-linear abelian
category -- the kernels and cokernels are degreewise. 

\begin{rem} \label{rem:4120}
Let 
$\opn{Ungr} : \dcat{G}_{\mrm{str}}(\K) \to \dcat{M}(\K)$
be the functor that forgets the grading. It is faithful, but often not full. 
Namely the obvious homomorphism 
\[ \opn{Ungr} \bigl( \opn{Hom}_{\K}(M, N) \bigr) \to 
\opn{Hom}_{\K} \bigl( \opn{Ungr}(M), \opn{Ungr}(N) \bigr) \]
is injective but not bijective. 
See Remark \ref{rem:3795} for a discussion. 
\end{rem}

The tensor operation $(- \ot -)$ from (\ref{eqn:1691}) makes 
$\dcat{G}_{\mrm{str}}(\K)$ into a {\em monoidal $\K$-linear category},
with monoidal unit $\K$. This means that the bifunctor 
\[ (- \ot -) : \dcat{G}_{\mrm{str}}(\K) \times 
\dcat{G}_{\mrm{str}}(\K) \to \dcat{G}_{\mrm{str}}(\K) \]
satisfies the monoidal axioms (associativity up to a trifunctorial cocycle, 
etc., see \cite[Chapter XI]{Mac2}), and it is $\K$-bilinear. 
Moreover, given $M, N \in \dcat{G}_{\mrm{str}}(\K)$, let us define the 
{\em braiding isomorphism}% 
\index{Braiding isomorphism}
\begin{equation} \label{eqn:4100}
\opn{br}_{M, N} : M \ot N \iso N \ot M , 
\end{equation}
\begin{equation} \label{eqn:4101}
\opn{br}_{M, N}(m \ot n) := (-1)^{i \cd j} \cd n \ot m 
\end{equation}
for homogeneous elements $m \in M^i$ and $n \in N^j$.
Because 
$\opn{br}_{N, M} \circ \opn{br}_{M, N} = \opn{id}_{M \ot N}$,
this makes $\dcat{G}_{\mrm{str}}(\K)$ into a {\em symmetric 
monoidal $\K$-linear category}. The braiding isomorphism (\ref{eqn:4101})
is often called the {\em Koszul sign rule}% 
\index{Koszul sign rule}.
See Remark \ref{rem:2516} for a background discussion. 

\begin{exer} \label{exer:4100}
For an integer $l \geq 1$ let $\tau$ be a permutation of the set 
$\{ 1, \ldots, l \}$. Show that for every 
$M_1, \ldots, M_l \in \dcat{G}_{\mrm{str}}(\K)$ there is an isomorphism 
\[ \opn{br}_{\tau} : M_1 \ot \cdots \ot M_l \iso 
M_{\tau(1)} \ot \cdots \ot M_{\tau(l)} \]
in $\dcat{G}_{\mrm{str}}(\K)$, which is functorial in the sequence of objects
$\{ M_i \}_{i = 1, \ldots, l}$, functorial in the permutation $\tau$, monoidal
(for a partition $l = l_1 + l_2$), and 
$\opn{br}_{\tau} = \opn{br}_{M_1, M_2}$
when $l = 2$ and $\tau$ is the transposition.
See \cite[Theorem XI.1.1]{Mac2} or \cite{Berger} for the answer.
\end{exer}

\begin{dfn} \label{dfn:4103}
A {\em cohomologically graded central $\K$-ring}% 
\index{Cohomologically graded ring}
is a central $\K$-ring $A$, equipped with a direct sum decomposition
$A =  \bigoplus_{i \in \Z} A^i$
into $\K$-sub\-modules, such that $1_A \in A^0$, and 
$A^i \cd A^j \sub A^{i + j}$.
\end{dfn}

\begin{dfn} \label{dfn:4104}
Let $A$ and $B$ be cohomologically graded central $\K$-rings. 
A {\em homomorphism of cohomologically graded central $\K$-rings} 
is a $\K$-ring homomorphism $f : A \to B$ that respects the gradings, 
namely $f(A^i) \sub B^i$.

The category of cohomologically graded central $\K$-rings is denoted by 
$\catt{GRng} \centover \K$.  
\end{dfn}

As always for ring homomorphisms, $f$ must preserve units, i.e.\ 
$f(1_A) = 1_B$. Note that $\K$ itself is a graded ring, concentrated in degree 
$0$; and it is the initial object of $\catt{GRng} \centover \K$.  

To simplify the discussion, we shall follow the next convention (until Section 
\ref{sec:alg-gra-rings}). 

\begin{conv} \label{conv:4126}
By ``graded ring'' we mean a {\em cohomologically graded central $\K$-ring}, 
as defined above. 
\end{conv}

Recall that by Convention \ref{conv:2490}, all ring homomorphisms, and that 
includes graded rings, are over $\K$.

\begin{exa} \label{exa:1705}
Let $M$ be a graded $\K$-module. Then the graded module
\[ \opn{End}_{\K}(M) := \opn{Hom}_{\K}(M, M) = 
\bigoplus_{i \in \Z} \, \opn{Hom}_{\K}(M, M)^i , \]
with the operation of composition, is a graded $\K$-ring. 
\end{exa}

The prototypical manifestation of the Koszul sign rule is the next 
definition. 

\begin{dfn} \label{dfn:4102}
Let $A = \bigoplus_{i \in \Z} A^i$ be a graded $\K$-ring.
We say that $A$ is a {\em weakly commutative graded ring}% 
\index{Cohomologically graded ring! weakly commutative}
if 
$b \cd a = (-1)^{i \cd j} \cd a \cd b$
for all $a \in A^i$ and $b \in A^j$.
\end{dfn}

\begin{rem} \label{rem:4103}
Let us give a categorical explanation of Definition \ref{dfn:4102}, using the 
symmetric monoidal structure of $\dcat{G}_{\mrm{str}}(\K)$. 

The general categorical way to define a {\em ring object} in the symmetric
monoidal linear category
$\bigl( \dcat{G}_{\mrm{str}}(\K), \ot, \K, \opn{br} \big)$
is as an object $A \in \dcat{G}_{\mrm{str}}(\K)$ equipped with a 
multiplication morphism 
$\opn{m} : A \ot A \to A$
and a unit morphism $\opn{u} : \K \to A$, such that the triple 
$(A, \opn{m}, \opn{u})$ obeys the ring axioms. 
In our case it is precisely Definition \ref{dfn:4103}.

The {\em categorical commutativity condition} for the ring object 
$(A, \opn{m}, \opn{u})$  is that the diagram 
\[ \UseTips  \xymatrix @C=6ex @R=6ex {
A \ot A
\ar[dr]^{\opn{m}}
\ar[d]_{\opn{br}}
\\
A \ot A
\ar[r]^{\opn{m}}
&
A
} \]
in $\dcat{G}_{\mrm{str}}(\K)$ is commutative. This is what weak commutativity 
is, in Definition \ref{dfn:4102}.
\end{rem}

Here is a definition similar to Definition \ref{dfn:4102}.

\begin{dfn} \label{dfn:4105}
Let $A = \bigoplus_{i \in \Z} A^i$ be a graded ring.
\begin{enumerate}
\item Homogeneous elements $a \in A^i$ and $b \in A^j$ are said to {\em 
graded-commute} with each other if 
$b \cd a = (-1)^{i \cd j} \cd a \cd b$.

\item A homogeneous element $a \in A^i$ is called a {\em graded-central 
element} 
if it graded-commutes with all homogeneous elements of $A$. 

\item The {\em graded-center} of $A$ is the $\K$-submodule
$\opn{Cent}(A) \sub A$ generated by the homogeneous 
graded-central elements.
\end{enumerate}
\end{dfn}

\begin{exer} \label{exer:1706}
Let $A$ be a graded ring. Show that:
\begin{enumerate}
\item $\opn{Cent}(A)$ is a graded subring of $A$, it is weakly commutative, and 
it contains the image of the base ring $\K$. 

\item $A$ is weakly commutative iff $\opn{Cent}(A) = A$. 
\end{enumerate}
\end{exer}

Below are several sign formulas that are more subtle consequences of the Koszul 
sign rule. They can be traced -- with effort -- to the {\em biclosed monoidal 
structure} of $\dcat{G}_{\mrm{str}}(\K)$, namely to the interaction between 
the bifunctors $(- \ot -)$ and $\opn{Hom}_{\K}(-, -)$. 

Suppose that for $k = 0, 1$ we are given graded $\K$-module
homomorphisms $\phi_k : M_k \to N_k$ of degrees $i_k$. 
Then the homomorphism 
\[ \phi_0 \ot \phi_1 \in 
\opn{Hom}_{\K} \bigl( (M_0 \ot M_1, N_0 \ot N_1 \bigr)^{i_0 + i_1}  \]
acts on a tensor $m_0 \ot m_1 \in M_0 \ot M_1$, 
with $m_k \in M_k^{j_k}$, like this: 
\begin{equation} \label{eqn:4105}
(\phi_0 \ot \phi_1)(m_0 \ot m_1) := (-1)^{i_1 \cd j_0} \cd 
\phi_0(m_0) \ot \phi_1(m_1) \in N_0 \ot N_1 .
\end{equation}
The rule of thumb explaining this formula is that $\phi_1$ and $m_0$ were 
transposed. 

Suppose we are given graded $\K$-module
homomorphisms $\phi_0 : N_0 \to M_0$
and $\phi_1 : M_1 \to N_1$ of degrees $i_0$ and $i_1$. 
Then the homomorphism 
\[ \opn{Hom}(\phi_0, \phi_1) \in 
\opn{Hom}_{\K} \bigl( \opn{Hom}_{\K}(M_0, M_1), 
\opn{Hom}_{\K}(N_0, N_1) \bigr) ^{i_0 + i_1}  \]
acts on $\ga \in \opn{Hom}_{\K}(M_0, M_1)^j$
as follows: for an element $n_0 \in N^k_0$ we have
\begin{equation} \label{eqn:4106}
\opn{Hom}(\phi_0, \phi_1)(\ga)(n_0) := 
(-1)^{i_0 \cd (i_1 + j)}(\phi_1 \circ \ga \circ \phi_0)(n_0) \in 
N_1^{k + i_0 + i_1 + j} . 
\end{equation}
The sign is because $\phi_0$ jumped across $\phi_1$ and $\ga$. 

\begin{dfn} \label{dfn:4107} 
Let $A$ and $B$ be graded  rings. Then 
$A \ot B$ is a graded ring, with multiplication 
\[ (a_0 \ot b_0) \cd (a_1 \ot b_1) := 
(-1)^{i_1 \cd j_0} \cd (a_0 \cd a_1) \ot (b_0 \cd b_1)  \]
for elements $a_k \in A^{i_k}$ and  $b_k \in B^{j_k}$.
\end{dfn}

We shall require another notion of commutativity, that is not of categorical 
nature. 

\begin{dfn}  \label{dfn:3091} 
Let $A = \bigoplus_{i \in \Z} A^i$ be a graded ring.
\begin{enumerate}
\item The graded ring  $A$ is called 
{\em strongly commutative}%
\index{Cohomologically graded ring! strongly commutative}
if it is weakly commutative (Definition \ref{dfn:4102}), 
and also $a^2 = 0$ if $a \in A^i$ and $i$ is odd. 

\item The graded ring $A$ is called {\em nonpositive}%
\index{Cohomologically graded ring! nonpositive}
if $A^i = 0$ for all $i > 0$.

\item The graded ring $A$ is called a 
{\em commutative graded ring}%  
\index{Cohomologically graded ring! commutative}
if it is nonpositive and strongly commutative. 
\end{enumerate}
\end{dfn}

This definition is taken from \cite{Ye9}. In \cite{YeZh3} the term 
``super-commutative'' was used instead of ``strongly commutative''. 
The name ``strongly commutative'' was suggested to us by J. Palmieri. 

\begin{exa} \label{exa:4625}
By a {\em graded set}%
\index{Graded set}
we mean a set $X$ partitioned as 
$X = \coprod_{i \in \Z} X^i$. 
The elements of $X^i$ are called variables of degree $i$. The 
{\em noncommutative graded polynomial ring}% 
\index{Graded polynomial ring! noncommutative}
on the graded set $X$ is the graded ring $\K \bra{X}$, which is 
the free graded $\K$-module spanned by 
the monomials (i.e.\ words) $x_1 \cdots x_n$ in the elements of $X$, and its 
multiplication is by concatenation of monomials. 

The {\em strongly commutative graded polynomial ring}%
\index{Graded polynomial ring! commutative}
on the graded set $X$ is the graded ring 
$\K[X] := \K \bra{X} / I$,
where $I$ is the two-sided ideal of $\K \bra{X}$ generated by the elements 
$y \cd x - (-1)^{i \cd j} \cd x \cd y$ for all variables 
$x \in X^i$ and $y \in X^j$, together with the elements 
$z \cd z$ for all $z \in X^k$ and odd $k$. 
This is a strongly commutative graded ring. 
\end{exa}

\begin{exer} \label{exer:4105}
Show that if $A$ and $B$ are weakly (resp.\ strongly) commutative graded rings, 
then so is $A \ot B$. 
\end{exer}

\begin{rem} \label{rem:2516}
Weak commutativity is the obvious commutativity condition in the cohomologically
graded setting, when the Koszul sign rule%
\index{Koszul sign rule}
is imposed; see Remark 
\ref{rem:4103}. Of course there are many instances of commutative graded rings 
that do not involve the Koszul sign rule; see e.g.\ \cite{Eis} or \cite{Mats}. 
In our book we call them {\em algebraically graded rings}, and they are studied 
in Sections \ref{sec:alg-gra-rings}-\ref{sec:BDC}.

Strong commutativity has another reason. Its role is to guarantee that the 
strongly commutative graded polynomial ring $\K[X]$ from Example 
\ref{exa:4625} above will be a graded-free $\K$-module. Without this 
condition, the square of an odd variable $z$ would be a nonzero $2$-torsion 
element. 

Of course, if $2$ is invertible in $\K$ (e.g.\ if $\K$ contains $\Q$), then 
weak and strong commutativity of a graded central $\K$-ring coincide. Since 
most texts dealing with DG rings assume that $\Q \sub \K$, the subtle 
distinction we make is absent from them. 
\end{rem}

\begin{rem} \label{rem:2992}
Let $A = \bigoplus_{i \in \Z} A^i$ be a weakly commutative graded ring, and 
assume $A$ has nonzero odd elements. Then, after we forget the grading, the 
ring $A$ is no longer commutative (except in special cases, like in 
characteristic $2$). 

This phenomenon will reappear often in the study of DG algebra.
\end{rem}

\begin{rem} \label{rem:2991}
The origin of the Koszul sign rule%
\index{Koszul sign rule},
predating the work of J.L. Koszul, 
appears to be in the formula for the differential $\d$ of the tensor product of 
two complexes, that occurs in classical homological algebra (and is repeated 
here in Definition \ref{dfn:2992}). Without the sign we won't have 
$\d \circ \d = 0$. 

In his thesis \cite{Kosz2}, Koszul talks about the sign rule explicitly,
when discussing the Chevalley-Eilenberg complex. 
See \cite{Berger} for a detailed study of the Koszul sign rule and its 
group-theoretic interpretation.  

It should be noted that there are some sign inconsistencies in the book 
\cite{RD}. The Koszul sign rule could have prevented them. 

In our book we made a strenuous effort to have correct (namely consistent) 
signs. This did not always produce satisfactory outcomes -- for instance, in 
the context of the triangulated structure of the opposite homotopy category, as 
explained in Remark \ref{rem:2306}. 
\end{rem}

\begin{dfn} \label{dfn:4127}
Let $A$ be a graded ring. A 
{\em graded left $A$-module}%
\index{Graded module}
is a left $A$-module $M$, equipped with a $\K$-module decomposition 
$M =  \bigoplus_{i \in \Z} M^i$, such that 
$A^i \cd M^j \sub M^{i + j}$ for all $i, j$. 
\end{dfn}

We can also talk about graded right $A$-modules, and graded bimodules. 
But our default option (see Convention \ref{conv:2490}) is that modules are 
left modules.

\begin{exer} \label{exer:4110}
Let $M$ be a graded $\K$-module, $A$ a graded central $\K$-ring, and
$f : A \to \opn{End}_{\K}(M)$ a homomorphism in $\catt{GRng} \centover \K$.
\begin{enumerate}
\item Show that $M$ becomes a graded $A$-module, with action  
$a \cd m := f(a)(m)$.

\item Show that every graded $A$-module structure on $M$, that is consistent 
with the given graded $\K$-module structure, arises this way. 
\end{enumerate}  
\end{exer}

\begin{lem} \label{lem:1690}
Let $A$ be a graded ring, let $M$ be a graded right $A$-module, 
and let $N$ be a graded left $A$-module. Then the $\K$-module 
$M \ot_A N$ has a direct sum decomposition
\[ M \ot_A N = \bigoplus_{i \in \Z} \, (M \ot_A N)^i, \]
where $(M \ot_A N)^i$ is the $\K$-linear span of the tensors $m \ot n$, for 
all $j \in \Z$, $m \in M^j$ and  $n \in N^{i - j}$. 
\end{lem}

\begin{proof}
There is a canonical surjection of $\K$-modules 
$M \ot N \to M \ot_A N$.
Its kernel is the $\K$-submodule 
$L \subseteq M \ot N$ generated by the elements 
$(m \cd a) \ot n - m \ot (a \cd n)$,
for $m \in M^j$, $n \in N^k$ and $a \in A^l$.
So $L$ is a graded submodule of  $M \ot N$, and therefore so is the 
quotient. Finally, by formula (\ref{eqn:1691}) the $i$-th homogeneous component 
of $M \ot_A N$ is precisely  $(M \ot_A N)^i$. 
\end{proof}

\begin{dfn} \label{dfn:1690}
Let $A$ be a graded ring, and let $M, N$ be graded $A$-modules.
For each $i \in \Z$ define $\opn{Hom}_{A}(M, N)^i$ to be the subset of 
$\opn{Hom}_{\K}(M, N)^i$ consisting of the homomorphisms $\phi : M \to N$ such 
that
$\phi(a \cd m) = (-1)^{i \cd k} \cd a \cd \phi(m)$
for all $a \in A^k$. Next define the graded $\K$-module
\[ \opn{Hom}_{A}(M, N) := \bigoplus_{i \in \Z} \, \opn{Hom}_{A}(M, N)^i . \]
\end{dfn}

Suppose $\cat{C}$ is a $\K$-linear category (Definition \ref{dfn:1083}).
Since the composition of morphisms is $\K$-bilinear, for every  
triple of objects $M_0, M_1, M_2 \in \cat{C}$, composition can be expressed as
a $\K$-linear homomorphism 
% \[ \begin{aligned}
% &  \opn{Hom}_{\cat{C}}(M_1, M_2) \ot 
% \opn{Hom}_{\cat{C}}(M_0, M_1) \xar{} 
% \opn{Hom}_{\cat{C}}(M_0, M_2) ,  \\
% & \phi_1 \ot \phi_0 \mapsto \phi_1 \circ \phi_0 . 
% \end{aligned} \]
\[ \opn{Hom}_{\cat{C}}(M_1, M_2) \ot 
\opn{Hom}_{\cat{C}}(M_0, M_1) \xar{} 
\opn{Hom}_{\cat{C}}(M_0, M_2) , \]
$\phi_1 \ot \phi_0 \mapsto \phi_1 \circ \phi_0$.
We refer to it as the composition homomorphism. It will be used in the 
following definition. 

\begin{dfn} \label{dfn:1692}
A {\em graded $\K$-linear category}%
\index{Graded category}
is a $\K$-linear category $\cat{C}$, 
endowed with a grading on each of the $\K$-modules 
$\opn{Hom}_{\cat{C}}(M_0, M_1)$. The conditions are these:
\begin{enumerate}
\rmitem{a} For every object $M$, the identity automorphism $\opn{id}_M$ has 
degree $0$. 

\rmitem{b} For every triple of objects $M_0, M_1, M_2 \in \cat{C}$, the 
composition homomorphism 
\[ \opn{Hom}_{\cat{C}}(M_1, M_2) \ot \opn{Hom}_{\cat{C}}(M_0, M_1)
\xar{} \opn{Hom}_{\cat{C}}(M_0, M_2) \]
is a strict homomorphism of graded $\K$-modules. 
\end{enumerate}
\end{dfn}

In item (b) we use the graded module structure on a tensor product from 
equation (\ref{eqn:1691}).
A morphism $\phi \in \opn{Hom}_{\cat{C}}(M_0, M_1)^i$ is called a 
morphism of degree $i$. 
 
\begin{dfn} \label{dfn:1693} 
Let $\cat{C}$ be a graded $\K$-linear category.
The {\em strict subcategory of $\cat{C}$}%
\index{Strict subcategory! of a graded category}
is the subcategory $\opn{Str}(\cat{C})$ on all objects of $\cat{C}$, whose 
morphisms are the degree $0$ morphisms of $\cat{C}$.
\end{dfn}

\begin{exa} \label{exa:1696}
Let $A$ be a graded ring. Define $\cat{GMod} A$ to be 
the category whose objects are the graded $A$-modules. For 
$M, N \in \cat{GMod} A$, the set of morphisms is the graded $\K$-module 
$\opn{Hom}_{\cat{GMod} A}(M, N) := \opn{Hom}_{A}(M, N)$ 
from Definition \ref{dfn:1690}. Then $\cat{GMod} A$ is a graded $\K$-linear 
category. The morphisms in the subcategory
$\cat{GMod}_{\mrm{str}} A := \opn{Str}(\cat{GMod} A)$ are the 
strict homomorphisms of graded $A$-modules.
We often write
$\dcat{G}(A) := \cat{GMod} A$ and 
$\dcat{G}_{\mrm{str}}(A) := \cat{GMod}_{\mrm{str}} A$.  
In the special case $A = \K$, the category $\dcat{G}_{\mrm{str}}(\K)$
already appeared in Definition \ref{dfn:4100}. 
\end{exa}

\begin{rem} \label{rem:1190}
The name ``strict homomorphism of graded modules'', and the corresponding 
notations $\cat{GMod}_{\mrm{str}} A = \dcat{G}_{\mrm{str}}(A)$, 
are new. We introduced them to distinguish the 
abelian category $\cat{GMod}_{\mrm{str}} A$ from the graded category 
$\cat{GMod} A$ that contains it.
See Definitions \ref{dfn:1063} and
\ref{dfn:1212} for the DG versions of these notions.
\end{rem}

\begin{dfn} \label{dfn:1694}
Let $\cat{C}$ and $\cat{D}$ be graded $\K$-linear categories. 
A functor $F : \cat{C} \to \cat{D}$
is called a 
{\em graded $\K$-linear functor}%
\index{Graded functor}
if it satisfies this condition: 
\begin{itemize}
\item[$\vartriangleright$] For every pair of objects $M_0, M_1 \in \cat{C}$, 
the function 
\[ F : \opn{Hom}_{\cat{C}}(M_0, M_1) \to  
\opn{Hom}_{\cat{D}} \bigl( F(M_0), F(M_1) \bigr) \] 
is a strict homomorphism of graded $\K$-modules. 
\end{itemize}
\end{dfn}

\begin{conv} \label{conv:4625}
To simplify the terminology, we shall often use the expressions ``graded 
category'' and ``graded functor'' as abbreviations for ``graded $\K$-linear
category'' and ``graded $\K$-linear functor'', respectively. 
\end{conv}
 
\begin{exa} \label{exa:1694}
Let $A$ be a graded ring. We can view $A$ as a category 
$\cat{A}$ with a single object, and it is a graded category.
If $f : A \to B$ is a homomorphism of graded rings, then 
passing to single-object categories we get a 
graded functor $F : \cat{A} \to \cat{B}$. 
\end{exa}

Recall that ``morphism of functors'' is synonymous with ``natural 
transformation''. 

\begin{dfn} \label{dfn:1695}
Let  $F, G : \cat{C} \to \cat{D}$
be graded functors between graded categories, and let $i \in \Z$. 
A {\em degree $i$ morphism of graded functors}%
\index{Graded functor! morphism of {\indash}s}
$\eta : F \to G$ is a collection 
$\eta = \{ \eta_M \}_{M \in \cat{C}}$ of morphisms 
$\eta_M \in \opn{Hom}_{\cat{D}} \bigl( F(M), G(M) \bigr)^i$,
such that for every morphism 
$\phi \in \opn{Hom}_{\cat{C}} (M_0, M_1)^j$ there is equality 
\[ G(\phi) \circ \eta_{M_0} = (-1)^{i \cd j} \cd  \eta_{M_1} \circ F(\phi) \]
inside 
$\opn{Hom}_{\cat{D}} \bigl( F(M_0), G(M_1) \bigr)^{i + j}$.
\end{dfn}

If $i$ is odd in the definition above, then after we forget the grading, 
$\eta : F \to G$ is usually no longer a morphism of functors; this is another 
instance of the phenomenon mentioned in Remark \ref{rem:2992}.

\begin{dfn}  \label{dfn:1700}
Let $\cat{M}$ be an abelian category. 
A {\em graded object in $\cat{M}$}%
\index{Abelian category! graded object in}
is a collection 
$\{ M^i \}_{i \in \Z}$ of objects $M^i \in \cat{M}$.
\end{dfn}

Because we did not assume that $\cat{M}$ has countable direct sums, the graded 
objects are ``external'' to $\cat{M}$; cf.\ Exercise \ref{exer:4112}.

Suppose $M = \{ M^i \}_{i \in \Z}$ and $N = \{ N^i \}_{i \in \Z}$
are graded objects in $\cat{M}$. For an integer $i$ we define the $\K$-module
\begin{equation} \label{eqn:2495}
\opn{Hom}_{\cat{M}}(M, N)^i := \prod_{j \in \Z} 
\opn{Hom}_{\cat{M}}(M^j, N^{j + i}) .
\end{equation}
We get a graded $\K$-module 
\begin{equation} \label{eqn:1700}
\opn{Hom}_{\cat{M}}(M, N) := \bigoplus_{i \in \Z} \,
\opn{Hom}_{\cat{M}}(M, N)^i . 
\end{equation}

\begin{dfn}  \label{dfn:1701}
Let $\cat{M}$ be an abelian category. 
The {\em category of graded objects in $\cat{M}$} is the graded linear 
category $\dcat{G}(\cat{M})$, whose objects are the graded objects 
in $\cat{M}$, and the morphism sets are the graded modules 
\[ \opn{Hom}_{\dcat{G}(\cat{M})}(M, N) := \opn{Hom}_{\cat{M}}(M, N) \]
from equation (\ref{eqn:1700}). 
The composition operation is the obvious one. 
\end{dfn}

\begin{exer}  \label{exer:4112}
Suppose $\cat{M} = \dcat{M}(A)$, the category of modules over a 
central $\K$-ring $A$. For every 
$M = \{ M^i \}_{i \in \Z} \in \dcat{G}(\cat{M})$ let 
$F(M) := \bigoplus_{i \in \Z} M^i$.
Then $F(M)$ is a graded $A$-module, as discussed earlier, so $F(M)$ is an 
object of the category $\dcat{G}(A)$ from Example \ref{exa:1696}.
Prove that 
$F : \dcat{G}(\cat{M}) \to \dcat{G}(A)$
is an isomorphism of graded categories. 
\end{exer}

In the next definition we combine graded rings and linear categories, to 
concoct a new hybrid.

\begin{dfn} \label{dfn:1702}
Let $\cat{M}$ be an abelian category, and let $A$ be a graded ring. A 
{\em graded $A$-module in  $\cat{M}$}%
\index{Abelian category! graded $A$-module in}
is an object 
$M \in \dcat{G}(\cat{M})$, together with a graded ring homomorphism 
$f : A \to \opn{End}_{\cat{M}}(M)$. 
The set of graded $A$-modules in $\cat{M}$ is denoted by 
$\dcat{G}(A, \cat{M})$. 
\end{dfn}

What the definition says is that an element 
$a \in A^i$ gives rise to a degree $i$ endomorphism 
$f(a)$ of the graded object $M  = \{ M^j \}_{j \in \Z}$. In turn, this means 
that for every $j$, 
$f(a) : M^j \to M^{j + i}$ is a morphism in $\cat{M}$.  
The operation $f$ satisfies $f(1_A) = \opn{id}_M$ and 
$f(a_1 \cd a_2) = f(a_1) \circ f(a_2)$. 

\begin{exa}  \label{exa:1701}
If $A = \K$, then $\dcat{G}(A, \cat{M}) = \dcat{G}(\cat{M})$; 
and if $\cat{M} = \cat{Mod} \K$, then 
$\dcat{G}(A, \cat{M}) = \dcat{G}(A)$.
\end{exa}

The next definition is a variant of Definition \ref{dfn:1690}.

\begin{dfn} \label{dfn:1703}
Let $\cat{M}$ be an abelian category, and let $A$ be a 
graded ring. For $M, N \in \dcat{G}(A, \cat{M})$ and $i \in \Z$ 
we define 
$\opn{Hom}_{A, \cat{M}}(M, N)^i$ to be the subset of 
$\opn{Hom}_{\cat{M}}(M, N)^i$ consisting of the  morphisms 
$\phi : M \to N$ such that
\[ \phi \circ f_M(a) = (-1)^{i \cd k} \cd f_N(a) \circ \phi \]
for all $a \in A^k$. Next let 
\[ \opn{Hom}_{A, \cat{M}}(M, N) := 
\bigoplus_{i \in \Z} \, \opn{Hom}_{A, \cat{M}}(M, N)^i . \]
This is a graded $\K$-module.
\end{dfn}

\begin{dfn} \label{dfn:1704}
Let $\cat{M}$ be an abelian category, and let $A$ be a graded   
ring. The {\em category of graded $A$-modules in $\cat{M}$}
is the graded category $\dcat{G}(A, \cat{M})$
whose objects are the graded $A$-modules in 
$\cat{M}$, and the morphism graded $\K$-modules are 
\[ \opn{Hom}_{\dcat{G}(A, \cat{M})}(M_0, M_1) :=
\opn{Hom}_{A, \cat{M}}(M_0, M_1)  \]
from Definition \ref{dfn:1703}. 
The compositions are those of $\dcat{G}(\cat{M})$. 
\end{dfn}

In the special case that $A$ is a ring (i.e.\ a graded ring 
concentrated in degree $0$), this is \cite[Definition 8.5.1]{KaSc2}. 

Notice that forgetting the action of $A$ is a faithful graded 
functor $\dcat{G}(A, \cat{M}) \to \dcat{G}(\cat{M})$.
As in every graded category, there is the subcategory
\begin{equation} \label{eqn:2360}
\dcat{G}_{\mrm{str}}(A, \cat{M}) := \opn{Str}(\dcat{G}(A, \cat{M}))
\sub \dcat{G}(A, \cat{M})
\end{equation}
of strict (i.e.\ degree $0$) morphisms. 

\begin{exer} \label{exer:1705}
Show that $\dcat{G}_{\mrm{str}}(A, \cat{M})$ is an abelian category. 
\end{exer}

\begin{rem} \label{rem:2517}
The reader may have noticed that we can talk about the graded category 
$\dcat{G}(\cat{M})$ for every linear category $\cat{M}$, regardless if it 
is abelian or not. We chose to restrict attention to the abelian case for a 
pedagogical reason: this will hopefully reduce confusion between the many sorts 
of graded (and later DG) categories that occur in our discussion. 
\end{rem}



\mysubsection{DG \texorpdfstring{$\K$}{K}-Modules}

Recall that Conventions \ref{conv:2490} and \ref{conv:4125} are in force. 

\begin{dfn} \label{dfn:1070}
A {\em DG $\K$-module}%
\index{Differential graded module}
 is a graded $\K$-module 
$M = \bigoplus_{i \in \Z} M^i$, together with a $\K$-linear homomorphism 
$\d_M : M \to M$ of degree $1$, called the {\em differential}, satisfying 
$\d_M \circ \d_M = 0$.
\end{dfn}

When there is no danger of confusion, we may write $\d$ instead of $\d_M$. 

\begin{dfn} \label{dfn:1084}
Let $M$ and $N$ be DG $\K$-modules. A 
{\em strict homomorphism of DG $\K$-modules}%
\index{Strict homomorphism! of DG modules}
is a $\K$-linear homomorphism $\phi : M \to N$ 
of degree $0$ that commutes with the differentials. 
The resulting category is denoted by $\cat{DGMod}_{\mrm{str}} \K$, or by the 
abbreviated notation $\dcat{C}_{\mrm{str}}(\K)$. 
\end{dfn}

It is easy to see that $\dcat{C}_{\mrm{str}}(\K)$ is a $\K$-linear abelian 
category. There is a forgetful functor 
$\dcat{C}_{\mrm{str}}(\K) \to \dcat{G}_{\mrm{str}}(\K)$, 
$(M, \d_M) \mapsto M$.

\begin{dfn} \label{dfn:2992}
Suppose $M$ and $N$ are DG $\K$-modules. 
\begin{enumerate}
\item 
\index{Differential graded module! tensor product of {\indash}s}
The graded $\K$-module structure on the tensor product $M \ot N$
was given in equation (\ref{eqn:1691}). 
We put on it the differential 
\[ \d(m \ot n) := \d_M(m) \ot n + (-1)^i \cd m \ot \d_N(n) \]
for $m \in M^i$ and $n \in N^j$. In this way $M \ot N$ becomes a DG 
$\K$-module. We sometimes write $\d_{M \ot N}$ for this differential.

\item 
\index{Differential graded module! Hom between {\indash}s}
The graded module $\opn{Hom}_{\K}(M, N)$
was introduced in Definition \ref{dfn:4120}. We give it this 
differential: 
\[ \d(\phi) := \d_N \circ \phi - (-1)^i \cd \phi \circ \d_M \]
for $\phi \in \opn{Hom}_{\K}(M, N)^i$. 
Thus $\opn{Hom}_{\K}(M, N)$ becomes a DG $\K$-module. 
We sometimes denote this differential by 
$\d_{\opn{Hom}_{\K}(M, N)}$.
\end{enumerate}
\end{dfn}

\begin{dfn} \label{dfn:2993}
Let $M$ be a DG $\K$-module, and let $i$ be an integer. 
\begin{enumerate}
\item The module of degree $i$ {\em cocycles}%
\index{Cocycles! module of}%
\index{1-Zi(M)@$\opn{Z}^i(M)$}
of $M$ is
\[ \opn{Z}^i(M) := \opn{Ker} \bigl( M^i \xar{\d_M} M^{i + 1} \bigr) . \]

\item The module of degree $i$ {\em coboundaries}% 
\index{Coboundaries! module of}%
\index{1-Bi(M)@$\opn{B}^i(M)$}
of $M$ is 
\[ \opn{B}^i(M) := \opn{Im} \bigl( M^{i -1} \xar{\d_M} M^{i} \bigr) . \]

\item The module of degree $i$ {\em decocycles}%
\index{Decocycles! module of}%
\index{1-Yi(M)@$\opn{Y}^i(M)$}
of $M$ is 
\[ \opn{Y}^i(M) := \opn{Coker} \bigl( M^{i -1} \xar{\d_M} M^{i} \bigr) . \]

\item The $i$-th {\em cohomology}%
\index{Cohomology! module}%
\index{1-Hi(M)@$\opn{H}^i(M)$}
module of $M$ is 
\[ \opn{H}^i(M) := \opn{Z}^i(M) \, / \, \opn{B}^i(M)
\cong \opn{Coker} \bigl( M^{i - 1} \xar{\d_M} \opn{Z}^i(M) \bigr) . \]
\end{enumerate}
\end{dfn}

The definition above is standard, with the exception of item (3), which is 
a new contribution (both the notation $\opn{Y}^i(M)$ and the name 
``decocycles''). 
The fact that $\d_M \circ \d_M = 0$ implies that 
$\opn{B}^i(M) \sub \opn{Z}^i(M) \sub M^i$, so item (4) makes sense. 
On the other hand, since 
$\opn{Y}^i(M) = M^i / \opn{B}^i(M)$,
there is a canonical isomorphism 
\begin{equation} \label{eqn:2993}
\opn{H}^i(M) \cong \opn{Ker} \bigl( \opn{Y}^i(M) \xar{\d_M} M^{i + 1} \bigr) .
\end{equation}
The modules defined above are functorial in $M$; to be precise, these are 
$\K$-linear functors  
\begin{equation} \label{eqn:4127}
\opn{Z}^i, \opn{B}^i, \opn{Y}^i, \opn{H}^i  : 
\dcat{C}_{\mrm{str}}(\K) \to \dcat{M}(\K) . 
\end{equation}

We can rephrase Definition \ref{dfn:1084} using the notion of cocycles:  for 
DG $\K$-modules $M$ and $N$ there is equality
\begin{equation} \label{eqn:2520}
\opn{Hom}_{\dcat{C}_{\mrm{str}}(\K)}(M, N) = 
\opn{Z}^0 \bigl( \opn{Hom}_{\K}(M, N) \bigr)
\end{equation}
of submodules of $\opn{Hom}_{\K}(M, N)$. 



\mysubsection{DG Rings and Modules} \label{subsec:DG-rng-mod}

Recall that Conventions \ref{conv:2490}, \ref{conv:4125} and \ref{conv:4126} 
are in place. 

\begin{dfn} \label{dfn:1071}
A {\em DG central $\K$-ring}
\index{Differential graded ring}
is a graded central $\K$-ring 
$A =  \bigoplus_{i \in \Z} A^i$ (see Definition \ref{dfn:4103}), together with 
a $\K$-linear homomorphism $\d_A : A \to A$ of degree $1$, called the {\em 
differential}, satisfying  the equation
$\d_A \circ \d_A = 0$, and the graded Leibniz rule
\index{Graded Leibniz rule}
\[ \d_A(a \cd b) =  \d_A(a) \cd b + (-1)^i \cd a \cd \d_A(b) \]
for all $a \in A^i$ and $b \in A^j$. 
\end{dfn}

We sometimes write $\d$ instead of $\d_A$.

\begin{dfn} \label{dfn:1074}
Let $A$ and $B$ be DG central $\K$-rings. A 
{\em homomorphism of DG central $\K$-rings} 
\index{Differential graded ring! homomorphism of {\indash}s}
$f : A \to B$ is a graded central 
$\K$-ring homomorphism that 
commutes with the differentials of $A$ and $B$.
The resulting category is denoted by $\catt{DGRng} \centover \K$. 
\end{dfn}

Central $\K$-rings are viewed as DG central $\K$-rings concentrated in 
degree $0$ (and with trivial differentials). Thus the category $\catt{Rng} 
\centover \K$ is a full subcategory of $\catt{DGRng} \centover \K$. 

\begin{conv} \label{conv:4626}
To simplify terminology, we usually write ``DG ring'' instead of 
``DG central $\K$-ring''. 
\end{conv}

\begin{dfn} \label{dfn:3090}
\index{Differential graded ring! weakly commutative}
\index{Differential graded ring! strongly commutative}
\index{Differential graded ring! nonpositive}
\index{Differential graded ring! commutative}
A DG ring $A$ is called {\em weakly commutative}, {\em strongly 
commutative}, {\em nonpositive} or {\em commutative} if it is so, respectively, 
as a graded ring (after forgetting the differential), in the sense of 
Definition \ref{dfn:3091}. 

The corresponding full subcategories of $\catt{DGRng} \centover \K$ are 
$\catt{DGRng}_{\mrm{wc}} / \K$,  $\catt{DGRng}_{\mrm{sc}} / \K$,
$\catt{DGRng}^{\leq 0} \centover \K$ and 
$\catt{DGRng}^{\leq 0}_{\mrm{sc}} / \K$.
\end{dfn}
 
When $\K = \Z$ we can write $\catt{DGRng} := \catt{DGRng} \centover \Z$, etc.

Here are few examples of DG rings. First a silly example.

\begin{exa} \label{exa:1095}
Let $A$ be a graded ring. Then $A$, with the zero differential, 
is a DG ring.
\end{exa}

\begin{exa} \label{exa:1096}
Let $X$ be a differentiable (i.e.\ of type $\mrm{C}^{\infty}$) manifold over 
$\K := \R$. The de Rham complex $A$ of $X$ is a DG $\R$-ring, with the wedge 
product and the exterior differential. See \cite[Section 2.9.7]{KaSc1} for 
details. This is a strongly commutative DG ring, in the sense of Definition 
\ref{dfn:3090}. However, it is almost never nonpositive. 
\end{exa}

The next example  is the algebraic analogue of the previous one.
 
\begin{exa} \label{exa:1097}
Let $C$ be a commutative $\K$-ring. Then the algebraic de Rham complex
$A := \Omega_{C / \K} = \bigoplus_{p \geq 0} \Omega^p_{C / \K}$ is a
DG $\K$-ring. It is a strongly commutative DG ring, and it is 
almost never nonpositive. 
See \cite[Exercise 16.15]{Eis} or \cite[Section 25]{Mats} for details. 
\end{exa}

\begin{exa} \label{exa:1098}
Let $M$ be a DG $\K$-module. Consider the DG $\K$-module \lb 
$\opn{End}_{\K}(M) := \opn{Hom}_{\K}(M, M)$
from Definition \ref{dfn:2992}(2). 
Composition of endomorphisms is an associative 
multiplication on $\opn{End}_{\K}(M)$ that respects the grading, and the graded 
Leibniz rule holds. We see that 
$\opn{End}_{\K}(M)$ is a DG $\K$-ring.
It is usually neither weakly commutative nor nonpositive. 
\end{exa}

\begin{exa} \label{exa:1099}
Let $C$ be a commutative ring and let $c \in C$ be an element. 
The {\em Koszul complex}
\index{Koszul complex}
 of $c$ is the DG $C$-module $\opn{K}(C; c)$ defined as 
follows. In degree $0$ we let $\opn{K}^0(C; c) := C$. 
In degree $-1$, $\opn{K}^{-1}(C; c)$ is a free $C$-module of rank $1$, with 
basis element $x$. All other homogeneous components are trivial. The 
differential $\d$ is determined by what it does to the basis element
$x \in \opn{K}^{-1}(C; c)$, and we let 
$\d(x) := c \in \opn{K}^{0}(C; c)$. 

We want to make $\opn{K}(C; c)$ into a strongly commutative DG ring
(in the sense of Definition \ref{dfn:3090}). Since $x$ is an odd element,
we must define the graded ring structure by 
$\opn{K}(C; c) := C \ot \K[x]$,
where $\K[x]$ is the strongly commutative graded polynomial ring from Example 
\ref{exa:4625}, on the graded set 
$X = X^{-1} := \{ x \}$; i.e.\ $x^2 = 0$. 
\end{exa}

\begin{exa} \label{exa:1100}
Let $A$ and $B$ be DG rings. The graded ring 
$A \ot B$ from Definition \ref{dfn:4107}, with the differential
from Definition \ref{dfn:2992}(1),
is a DG ring.
\end{exa}

\begin{exa} \label{exa:1101}
Let $C$ be a commutative ring and let 
$\bsym{c} = (c_1, \ldots, c_n)$ be a sequence of elements in $C$. By combining 
Examples \ref{exa:1099} and \ref{exa:1100} we obtain the Koszul complex 
\index{Koszul complex}
\[ \opn{K}(C; \bsym{c}) := \opn{K}(C; c_1) \ot_{C} \cdots \ot_{C}
\opn{K}(C; c_n) . \]
This is a strongly commutative DG $C$-ring.
\end{exa}

\begin{rem} \label{rem:4685}
In the book \cite[Section 17, Exercises]{Eis} there are a few variants of the 
Koszul complex. Note that in most of the classical literature (including 
\cite{Eis} and \cite{Mats}), the multiplicative structure of $\opn{K}(C; 
\bsym{c})$ has been ignored. 
\end{rem}

\begin{dfn} \label{dfn:1080}
Let $A$ be a DG ring. The {\em opposite DG ring}
\index{Differential graded ring! opposite}
$A^{\mrm{op}}$
is the same DG $\K$-module as $A$, but the multiplication $\cdot^{\mrm{op}}$ of 
$A^{\mrm{op}}$ is the multiplication $\cd$ of $A$, reversed and twisted by 
signs: 
$a \cdot^{\mrm{op}}  b := (-1)^{i \cd j} \cd b \cd a$
for $a \in A^i$ and $b \in A^j$. 
\end{dfn}

\begin{exer} \label{exer:2500}
Verify that $A^{\mrm{op}}$ is a DG ring. 
\end{exer}

Note that $A$ is weakly commutative iff $A = A^{\mrm{op}}$. 

\begin{rem} \label{rem:4130}
In algebraic topology and homotopy theory it is customary to use lower indices 
for DG rings (and to call them DG algebras). In other words, they use {\em 
homological grading}, as opposed to our cohomological grading. Thus our 
nonpositive DG ring $A =  \bigoplus_{i \leq 0} A^i$ becomes a nonnegative 
graded algebra $A =  \bigoplus_{i \geq 0} A_i$ in their language. 

Another notion that is common in homotopy theory is that of a 
{\em connective DG algebra}; this is a DG ring 
$A =  \bigoplus_{i \in \Z} A_i$
whose homology $\opn{H}(A) = \bigoplus_{i \in \Z} \opn{H}_i(A)$
is nonnegative. In our cohomological language, the analogue would be a 
{\em coconnective DG ring}, that is a DG ring $A$ whose cohomology 
$\opn{H}(A) = \bigoplus_{i \in \Z} \opn{H}^i(A)$
is a nonpositive graded ring. Cf.\ Exercise \ref{exer:1100} below. 

If $A$ is a coconnective DG ring, then its smart truncation  
$A' := \opn{smt}^{\leq 0}(A)$, in the sense of Definition \ref{dfn:2320}, 
is a nonpositive DG ring, and the inclusion $A' \to A$ is a quasi-isomorphism. 
\end{rem}

\begin{dfn} \label{dfn:1081}
Let $A$ be a DG ring. A {\em DG left $A$-module}
\index{Differential graded module}
is a graded left $A$-module $M = \bigoplus_{i \in \Z} M^i$, together with a 
$\K$-linear homomorphism $\d_M : M \to M$ of degree $1$, called the 
differential, satisfying $\d_M \circ \d_M = 0$ and 
\[ \d_M(a \cd m) =  \d_A(a) \cd m + (-1)^i \cd a \cd \d_M(m) \]
for $a \in A^i$ and $m \in M^j$.
\end{dfn}

DG right $A$-modules are defined likewise, but we won't deal with them 
much. This is because DG right $A$-modules are DG left modules over the 
opposite DG ring $A^{\mrm{op}}$. More precisely, if $M$ is a DG right 
$A$-module, then the formula 
$a \cd m := (-1)^{i \cd j} \cd m \cd a$,
for $a \in A^i$ and $m \in M^j$, makes $M$ into a DG left 
$A^{\mrm{op}}$-module.

As implied by Convention \ref{conv:2490}(5), all DG modules are by default 
DG left modules. 
 
\begin{prop} \label{prop:1130}
Let $A$ be a DG $\K$-ring, and let $M$ be a DG $\K$-module. 
\begin{enumerate}
\item Suppose $f : A \to \opn{End}_{\K}(M)$ is a DG $\K$-ring homomorphism. 
Then the formula $a \cd m := f(a)(m)$, for $a \in A^i$ and $m \in M^j$,
makes $M$ into a DG $A$-module. 

\item Conversely, every DG $A$-module structure on $M$, that's compatible with 
the DG $\K$-module structure, arises in this way from a 
DG $\K$-ring homomorphism $f : A \to \opn{End}_{\K}(M)$.
\end{enumerate}
\end{prop}

\begin{exer} \label{exer:2520}
Prove this proposition. 
\end{exer}

\begin{dfn} \label{dfn:1285}
Let $M$ and $N$ be DG $A$-modules. A
{\em strict homomorphism of DG $A$-modules}
\index{Strict homomorphism! of DG modules}
 $\phi : M \to N$ 
is a strict homomorphism of DG $\K$-modules (Definition \ref{dfn:1084})
that respects the action of $A$. The resulting category is denoted by 
$\cat{DGMod}_{\mrm{str}} A$, or by the short notation 
$\dcat{C}_{\mrm{str}}(A)$%
\index{1-Cstr(A)@$\dcat{C}_{\mrm{str}}(A)$}.
\end{dfn}

The category $\dcat{C}_{\mrm{str}}(A)$ is abelian. See Proposition 
\ref{prop:4130} for a more general statement. 

\begin{exer} \label{exer:1100}
Let $A$ be a DG ring. Show that the the module of cocycles 
$\opn{Z}(A) := \bigoplus_{i \in \Z} \opn{Z}^i(A)$
is a graded subring of $A$, and the module of coboundaries
$\opn{B}(A) := \bigoplus_{i \in \Z} \opn{B}^i(A)$
is a two-sided graded ideal of $\opn{Z}(A)$. Conclude that the cohomology module
$\opn{H}(A) := \bigoplus_{i \in \Z} \opn{H}^i(A)$
is a graded ring. 

Let $f : A \to B$ be a homomorphism of DG rings. Show that 
$\opn{H}(f) : \opn{H}(A) \to \opn{H}(B)$
is a graded ring homomorphism. 
\end{exer}

\begin{exer} \label{exer:2521}
Let $A$ be a DG ring. Given a DG $A$-module $M$, show that its cohomology 
$\opn{H}(M)$ is a graded $\opn{H}(A)$-module. If $\phi : M \to N$
is a homomorphism in $\dcat{C}_{\mrm{str}}(A)$, then 
$\opn{H}(\phi) : \opn{H}(M) \to \opn{H}(N)$
is a homomorphism in $\dcat{G}_{\mrm{str}}(\opn{H}(A))$.
Conclude that 
\begin{equation} \label{eqn:4130}
\opn{H} : \dcat{C}_{\mrm{str}}(A) \to \dcat{G}_{\mrm{str}}(\opn{H}(A)) 
\end{equation}
is a linear functor.  
\end{exer}

\begin{dfn} \label{dfn:1100}
\index{Differential graded module! tensor product of {\indash}s}
Let $A$ be a DG $\K$-ring, let $M$ be a right DG $A$-module, and 
let $N$ be a left DG $A$-module. By Lemma \ref{lem:1690}, $M \ot_A N$
is a graded $\K$-module. We make it into a DG $\K$-module with the differential 
from Definition \ref{dfn:2992}(1). 
\end{dfn}

\begin{dfn} \label{dfn:1101}
\index{Differential graded module! Hom between {\indash}s}
Let $A$ be a DG $\K$-ring, and let $M$ and $N$ be left DG $A$-modules.
The graded $\K$-module $\opn{Hom}_{A}(M, N)$ from Definition \ref{dfn:1690}
is made into a DG $\K$-module with the differential from
Definition \ref{dfn:2992}(2). 
\end{dfn}

Generalizing formula (\ref{eqn:2520}), for DG $A$-modules $M$ and $N$ 
there is equality
\[ \opn{Hom}_{\dcat{C}_{\mrm{str}}(A)}(M, N) = 
\opn{Z}^0 \bigl( \opn{Hom}_{A}(M, N) \bigr) . \]

\begin{prop} \label{prop:4131}
Let $A$ be a DG ring. 
\begin{enumerate}
\item The category $\dcat{C}_{\mrm{str}}(A)$ has products. Given a collection 
$\{ M_x \}_{x \in X}$ of DG $A$-modules, their product 
$M  = \prod_{x \in X} M_x$ in $\dcat{C}_{\mrm{str}}(A)$
is $M := \bigoplus_{i \in \Z}  M^i$, where 
$M^i := \prod_{x \in X} M_x^i$.

\item The category $\dcat{C}_{\mrm{str}}(A)$ has direct sums. Given a 
collection $\{ M_x \}_{x \in X}$ of DG $A$-modules, their direct sum 
$M  = \bigoplus_{x \in X} M_x$ in $\dcat{C}_{\mrm{str}}(A)$
is $M := \bigoplus_{i \in \Z}  M^i$, where 
$M^i := \bigoplus_{x \in X} M_x^i$.

\item The functor $\opn{H}$ from \tup{(\ref{eqn:4130})} commutes with all 
products and direct sums. 
\end{enumerate}
\end{prop}

\begin{exer} \label{exer:4130}
Prove this proposition. 
\end{exer}





\mysubsection{DG Categories} \label{susec:DGCats}

In Definition \ref{dfn:1692} we saw graded categories. Here is the DG version. 
 
\begin{dfn} \label{dfn:1063}
A {\em DG $\K$-linear category}
\index{Differential graded category}
is a $\K$-linear category $\cat{C}$, endowed 
with a DG $\K$-module structure on each of the morphism $\K$-modules 
$\opn{Hom}_{\cat{C}}(M_0, M_1)$. The conditions are these:
\begin{itemize}
\rmitem{a} For every object $M$, the identity automorphism $\opn{id}_M$ is a 
degree $0$ cocycle in $\opn{Hom}_{\cat{C}}(M, M)$. 

\rmitem{b} For every triple of objects $M_0, M_1, M_2 \in \cat{C}$, the 
composition homomorphism 
\[ \opn{Hom}_{\cat{C}}(M_1, M_2) \ot \opn{Hom}_{\cat{C}}(M_0, M_1)
\xar{} \opn{Hom}_{\cat{C}}(M_0, M_2) \]
is a strict homomorphism of DG $\K$-modules. 
\end{itemize}
\end{dfn}

The differential in a DG $\K$-linear category $\cat{C}$ is sometimes denoted by 
$\d_{\cat{C}}$; for instance 
\begin{equation} \label{eqn:4135}
\d_{\cat{C}} : \opn{Hom}_{\cat{C}}(M_0, M_1)^{i} \to 
\opn{Hom}_{\cat{C}}(M_0, M_1)^{i + 1 } . 
\end{equation}

The next convention extends Convention \ref{conv:4626}. 

\begin{conv} \label{conv:4627}
We shall often write ``DG category'' instead of the longer expression
``DG $\K$-linear category''. 
\end{conv}

If $\cat{C}'$ is a full subcategory of a DG category $\cat{C}$, then of course
$\cat{C}'$ is itself a DG category. 

\begin{dfn} \label{dfn:1286}
Let $\cat{C}$ be a DG category.
\begin{enumerate}
\item A morphism $\phi \in \opn{Hom}_{\cat{C}}(M, N)^i$
is called a {\em degree $i$ morphism}. 

\item A morphism $\phi \in \opn{Hom}_{\cat{C}}(M, N)$
is called a {\em cocycle} if $\d_{\cat{C}}(\phi) = 0$. 

\item A morphism $\phi : M \to N$ in $\cat{C}$ is 
called a {\em strict morphism}%
\index{Differential graded category! strict morphism in}
if it is a degree $0$ cocycle.
\end{enumerate}
\end{dfn}

\begin{lem} \label{lem:1211}
Let $\cat{C}$ be a DG category, and for $i = 0, 1, 2$ let 
$\phi_i : M_{i} \to M_{i + 1}$ be a morphism in $\cat{C}$ of degree $k_i$.  
\begin{enumerate}
\item The morphism $\phi_1 \circ \phi_0$ has degree $k_0 + k_1$, and 
\[ \d_{\cat{C}}(\phi_1 \circ \phi_0) = 
\d_{\cat{C}}(\phi_1) \circ \phi_0 +
(-1)^{k_1} \cd \phi_1 \circ \d_{\cat{C}}(\phi_0) . \]

\item If $\phi_0$ and $\phi_1$ are cocycles, then so is $\phi_1 \circ \phi_0$.

\item If  $\phi_1$ is a coboundary, and $\phi_0$ and $\phi_2$ are cocycles, 
then $\phi_2 \circ \phi_1 \circ \phi_0$ is a coboundary.
\end{enumerate}
\end{lem}

\begin{proof}
(1) This is just a rephrasing of item (b) in Definition \ref{dfn:1063}. 

\medskip \noindent 
(2) This is immediate from (1). 

\medskip \noindent 
(3) Say $\phi_1 = \d_{\cat{C}}(\psi_1)$ for some degree $k_1 - 1$ morphism 
$\psi_1 : M_{1} \to M_2$. Then 
\[ \phi_2 \circ \phi_1 \circ \phi_0 = 
\d_{\cat{C}} \bigl( (-1)^{k_2} \cd \phi_2 \circ \psi_1 \circ \phi_0 \bigr) .
\qedhere \]
\end{proof}

The previous lemma makes the next definition possible. 

\begin{dfn} \label{dfn:1212} \mbox{ }
Let $\cat{C}$ be a DG category.
\begin{enumerate}
\item The {\em strict subcategory}%
\index{Differential graded category! strict subcategory of}%
\index{Strict subcategory! of a DG category}
of $\cat{C}$ is the category 
$\opn{Str}(\cat{C})$, with the same objects as 
$\cat{C}$, but with strict morphisms only. Thus 
\[ \opn{Hom}_{\opn{Str}(\cat{C})}(M, N) = 
\opn{Z}^0 \bigl( \opn{Hom}_{\cat{C}}(M, N) \bigr) . \]

\item The {\em homotopy category}%
\index{Differential graded category! homotopy category of}
of $\cat{C}$ is the category 
$\opn{Ho}(\cat{C})$, with the same objects as $\cat{C}$, 
and with morphism sets 
\[  \opn{Hom}_{\opn{Ho}(\cat{C})}(M, N) := 
\opn{H}^0 \bigl( \opn{Hom}_{\cat{C}}(M, N) \bigr) . \]

\item We denote by 
$\opn{P} : \opn{Str}(\cat{C}) \to \opn{Ho}(\cat{C})$
the functor which is  the identity on objects, and sends a strict morphism to 
its homotopy class. 
\end{enumerate}
\end{dfn}

The categories $\opn{Str}(\cat{C})$ and $\opn{Ho}(\cat{C})$ are linear, and 
the inclusion functor $\opn{Str}(\cat{C}) \to \cat{C}$ and the projection 
functor $\opn{P} : \opn{Str}(\cat{C}) \to \opn{Ho}(\cat{C})$
are linear. The first is faithful, and the second is full.

\begin{exa} \label{exa:1286}
If $\cat{A}$ is a DG category, then for every object $x \in \cat{A}$, 
its set of endomorphisms $A := \opn{End}_{\cat{A}}(x)$ is a DG ring.
Conversely, every DG ring $A$ can be viewed as a DG category $\cat{A}$ 
with a single object.
\end{exa}

\begin{exa} \label{exa:1285}
Let $A$ be a DG ring. The set of DG $A$-modules forms a 
DG category $\cat{DGMod} A$, in which the morphism DG $\K$-modules are
\[ \opn{Hom}_{\cat{DGMod} A}(M, N) := \opn{Hom}_{A}(M, N) \]
from Definition \ref{dfn:1101}. 
We shall often write $\dcat{C}(A) := \cat{DGMod} A$%
\index{1-C(A)@$\dcat{C}(A)$}. 
The strict subcategory here is 
\[ \opn{Str}(\cat{DGMod} A) = \cat{DGMod}_{\mrm{str}} A 
= \dcat{C}_{\mrm{str}}(A) ; \]
cf.\  Definition \ref{dfn:1285}. 
\end{exa}

Here is a result that will be used later. 

\begin{prop} \label{prop:1220}
Let $\phi : M \to N$ be a degree $i$ isomorphism in a DG category $\cat{C}$. 
Assume $\phi$ is a cocycle, namely 
$\d_{\cat{C}}(\phi) = 0$. 
Then its inverse $\phi^{-1} : N \to M$ is also a cocycle.  
\end{prop}

\begin{proof}
According the Leibniz rule (Lemma \ref{lem:1211}(1)), and the fact that 
$\opn{id}_M$ is a cocycle, we have
\[ \begin{aligned}
& 0 = \d_{\cat{C}}(\opn{id}_M) = \d(\phi^{-1} \circ \phi)
\\ & \ 
= \d_{\cat{C}}(\phi^{-1}) \circ \phi + 
(-1)^{-i} \cd \phi^{-1} \circ \, \d_{\cat{C}}(\phi) 
= \d_{\cat{C}}(\phi^{-1}) \circ \phi . 
\end{aligned} \]
Because $\phi$ is an isomorphism, we conclude that 
$\d_{\cat{C}}(\phi^{-1}) = 0$. 
\end{proof}

\begin{rem} \label{rem:1270}
The fact that the concept ``DG category'' refers both to DG rings 
(Example \ref{exa:1286}) and to the categories of DG modules over them (Example 
\ref{exa:1285}) is a source of frequent confusion. See Remarks \ref{rem:2517} 
and \ref{rem:1245}. 
\end{rem}

\begin{rem} \label{rem:1175}
For other accounts of DG categories see the relatively old references 
\cite{Kelly}, \cite{Kel1}, \cite{BoKa}, or the more recent \cite{To1}. 
An internet search can give plenty more information,
including the relation to simplicial and infinity categories. 

In this book we shall be exclusively concerned with the DG categories 
$\dcat{C}(A, \cat{M})$, to be introduced in Subsection \ref{subsec:DGModinM},
which have a lot more structure than other DG categories.

See Remark \ref{rem:1245} regarding the DG category
$\dcat{C}(\cat{A})$ of DG modules over a DG category $\cat{A}$, 
in the sense of \cite{Kel1}; and Remark \ref{rem:4960} regarding 
$\mrm{A}_{\infty}$ categories. 
\end{rem}



\mysubsection{DG Functors} \label{subsec:DGFunc}

Here $\cat{C}$ and $\cat{D}$ are DG $\K$-linear categories (see Definition 
\ref{dfn:1063}). When we forget differentials,  $\cat{C}$ and $\cat{D}$ become 
graded $\K$-linear categories. So we can talk about graded $\K$-linear
functors $\cat{C} \to \cat{D}$, as in Definition \ref{dfn:1694}. 

Recall the meaning of a strict homomorphism of DG $\K$-modules:
it has degree $0$ and commutes with the differentials. 

\begin{dfn} \label{dfn:1065}
Let $\cat{C}$ and $\cat{D}$ be DG $\K$-linear categories.
A functor  $F : \cat{C} \to \cat{D}$
is called a {\em DG $\K$-linear functor}
\index{Differential graded functor} 
if it satisfies this condition: 
\begin{itemize}
\item[$\vartriangleright$] For every pair of objects $M_0, M_1 \in \cat{C}$, 
the function 
\[ F : \opn{Hom}_{\cat{C}}(M_0, M_1) \to  
\opn{Hom}_{\cat{D}} \bigl( F(M_0), F(M_1) \bigr) \] 
is a strict homomorphism of DG $\K$-modules. 
\end{itemize}
\end{dfn}

In other words, $F$ is a DG functor if it is a graded functor, and 
$\d_{\cat{D}} \circ F = F \circ \d_{\cat{C}}$
as degree $1$ homomorphisms 
\[ \opn{Hom}_{\cat{C}}(M_0, M_1) \to  
\opn{Hom}_{\cat{D}} \bigl( F(M_0), F(M_1) \bigr) . \]

To match Convention \ref{conv:4627}, here is: 

\begin{conv} \label{conv:4628}
We shall usually write ``DG functor'' instead of the longer expression
``DG $\K$-linear functor''. 
\end{conv}

\begin{exa} \label{exa:1250}
Let $f : A \to B$ be a homomorphism of DG rings, and let 
$\cat{A}$ and $\cat{B}$ be the corresponding single-object DG  
categories. Then $f$ gives rise to a DG functor 
$F : \cat{A} \to \cat{B}$.
\end{exa}

Other examples of DG functors, more relevant to our study, will be given 
in Subsection \ref{subsec:exa-DGfun}. 

\begin{dfn} \label{dfn:1195}
Let $F, G : \cat{C} \to \cat{D}$ be DG functors.
\begin{enumerate}
\item A {\em degree $i$ morphism of DG functors}%
\index{Differential graded functor! morphism of {\indash}s}
$\eta : F \to G$ is a 
degree $i$ morphism of graded functors, as in Definition \ref{dfn:1695}.

\item Let $\eta : F \to G$ be a degree $i$ morphism of DG functors.
For each object $M \in \cat{C}$ there is a degree $i + 1$ morphism 
$\d_{\cat{D}}(\eta_M) : F(M) \to G(M)$
in $\cat{D}$. 
We define the collection of morphisms  
$\d_{\cat{D}}(\eta) := 
\{ \d_{\cat{D}}(\eta_M) \}_{M \in \cat{C}}$.

\item A {\em strict morphism of DG functors}
\index{Differential graded functor! strict morphism of {\indash}s}
is a degree $0$ morphism of graded functors  
$\eta : F \to G$ such that $\d_{\cat{D}}(\eta) = 0$. 
\end{enumerate}
\end{dfn}

Explanation of item (2) of the definition: for an object $M \in \cat{C}$, we 
have
\[ \eta_M \in \opn{Hom}_{\cat{D}} \bigl( F(M), G(M) \bigr)^i . \]
Now $\opn{Hom}_{\cat{D}} \bigl( F(M), G(M) \bigr)$
is a DG $\K$-module with differential $\d_{\cat{D}}$; see 
formula (\ref{eqn:4135}). Thus 
\[ \d_{\cat{D}}(\eta_M) \in 
\opn{Hom}_{\cat{D}} \bigl( F(M), G(M) \bigr)^{i + 1} . \]

\begin{prop} \label{prop:1295}
In the situation of Definition \tup{\ref{dfn:1195}}, the collection of 
morphisms $\d_{\cat{D}}(\eta)$  is a degree $i + 1$ morphism of DG 
functors $F \to G$. 
\end{prop}

\begin{exer} \label{exer:2522}
Prove this proposition.
\end{exer}

The categories $\opn{Str}(\cat{C})$ and 
$\opn{Ho}(\cat{C})$ were introduced in 
Definition \ref{dfn:1212}.

\begin{prop} \label{prop:1150}
Let $F : \cat{C} \to \cat{D}$ be a DG functor. Then $F$ induces linear 
functors
$\opn{Str}(F) : \opn{Str}(\cat{C}) \to \opn{Str}(\cat{D})$
and
$\opn{Ho}(F) : \opn{Ho}(\cat{C}) \to  \opn{Ho}(\cat{D})$.
\end{prop}

\begin{proof}
Because $F$ is a DG functor, it sends $0$-cocycles in 
$\opn{Hom}_{\cat{C}}(M_0, M_1)$ to $0$-cocycles in 
$\opn{Hom}_{\cat{D}}(F(M_0), F(M_1))$. 
The same for $0$-coboundaries.  
\end{proof}

By abuse of notation, and when there is no danger for confusion, we will 
sometimes write $F$ instead of $\opn{Str}(F)$ or $\opn{Ho}(F)$.

\begin{exer} \label{exer:1250}
Let $\cat{A}$ and $\cat{C}$ be DG categories, and assume 
$\cat{A}$ is small. Define 
$\cat{DGFun}(\cat{A}, \cat{C})$ to be the set of DG functors 
$F : \cat{A} \to \cat{C}$.
Show that $\cat{DGFun}(\cat{A}, \cat{C})$ is a DG category, 
where the morphisms are from Definition \tup{\ref{dfn:1195}(1)}, and their 
differentials are from Definition \tup{\ref{dfn:1195}(2)}.
\end{exer}


\mysubsection{Complexes in Abelian Categories} \label{subsec:complexes}

Here we recall facts about complexes from the classical homological theory, and 
place them within our context. In this subsection $\cat{M}$ is a $\K$-linear 
abelian category. 

A {\em complex of objects of $\cat{M}$}, or a 
{\em complex in $\cat{M}$}%
\index{Abelian category! complex in},
is a diagram 
\begin{equation} \label{eqn:1065}
\bigl( \cdots \to M^{-1} \xar{\d_M^{-1}}  M^0 \xar{\d_M^0}  
M^1 \xar{\d_M^1} M^2 \to \cdots \bigr)
\end{equation}
of objects and morphisms in $\cat{M}$, such that $\d_M^{i+1} \circ \d_M^i = 0$.
The collection of objects $M := \{ M^i \}_{i \in \Z}$ 
is nothing but a graded object of $\cat{M}$, as defined in Subsection 
\ref{subsec:gr-alg}. The collection of morphisms 
$\d_M := \{ \d^i_M \}_{i \in \Z}$ 
is called a {\em differential}, or a {\em coboundary operator}.
Thus a complex is a pair $(M, \d_M)$ made up of a graded object $M$ and a 
differential $\d_M$ on it. 
We sometimes write $\d$ instead of $\d_M$ or $\d_M^i$. 
At other times we leave the differential implicit, and just refer to the 
complex as $M$. 

Let $N$ be another complex in $\cat{M}$. A {\em strict morphism of complexes} 
$\phi : M \to N$ is a collection 
$\phi = \{ \phi^i \}_{i \in \Z}$ of morphisms $\phi^i : M^i \to N^i$ in
$\cat{M}$, such that
\begin{equation} \label{eqn:1254}
\d^i_N \circ \phi^i = \phi^{i+1} \circ \d^i_M .
\end{equation}
Note that a strict morphism $\phi : M \to N$ can be viewed as a commutative 
diagram 
\[ \UseTips  \xymatrix @C=8ex @R=6ex {
\cdots 
\ar[r]
&
M^i
\ar[r]^{\d_M^i}
\ar[d]_{\phi^i}
&
M^{i+1}
\ar[r]
\ar[d]_{\phi^{i+1}}
&
\cdots
\\
\cdots 
\ar[r]
&
N^i
\ar[r]^{\d_N^i}
&
N^{i+1}
\ar[r]
&
\cdots
} \]
in $\cat{M}$. 
The identity automorphism $\opn{id}_M$ of the complex $M$ is a strict 
morphism.  

\begin{rem} \label{rem:1210}
In most textbooks, what we call ``strict morphism of complexes'' is simply 
called a ``morphism of complexes''. See Remark \ref{rem:1190} for an 
explanation. 
\end{rem}

Let us denote by 
$\dcat{C}_{\mrm{str}}(\cat{M})$%
\index{1-Cstr(M)@$\dcat{C}_{\mrm{str}}(\cat{M})$}
the category of complexes in 
$\cat{M}$, with strict morphisms. This is a $\K$-linear abelian category. 
Indeed, the direct sum of complexes is the degree-wise direct sum, i.e.\ 
$(M \oplus N)^i = M^i \oplus N^i$.
The same for kernels and cokernels.
If  $\cat{N}$ is a full abelian subcategory of $\cat{M}$, then  
$\dcat{C}_{\mrm{str}}(\cat{N})$ is a full abelian subcategory of 
$\dcat{C}_{\mrm{str}}(\cat{M})$.

A single object $M^0 \in \cat{M}$ can be viewed as a complex 
\[ M :=  \bigl( \cdots \to 0 \to M^0 \to 0 \to \cdots \bigr) , \] 
where $M^0$ is in degree $0$; the differential of this complex is of course 
zero. The assignment $M^0 \mapsto M$ is a fully faithful $\K$-linear functor 
$\cat{M} \to \dcat{C}_{\mrm{str}}(\cat{M})$.

Let $M$ and $N$ be complexes in $\cat{M}$. As in (\ref{eqn:1700}) there is a 
graded $\K$-module $\opn{Hom}_{\cat{M}}(M, N)$. It is a DG $\K$-module with 
differential $\d$ given by the formula 
\begin{equation} \label{eqn:1130}
\d(\phi) := \d_N \circ \phi - (-1)^i \cd \phi \circ \d_M 
\end{equation}
for $\phi \in \opn{Hom}_{\cat{M}}(M, N)^i$. 
It is easy to check that $\d \circ \d = 0$. 
We sometimes denote this differential by
$\d_{\opn{Hom}}$ or $\d_{\opn{Hom}_{\cat{M}}(M, N)}$.

Thus, an element $\phi \in \opn{Hom}_{\cat{M}}(M, N)^i$
is a collection $\phi = \{ \phi^j \}_{j \in \Z}$ of morphisms 
$\phi^j : M^j \to N^{j + i}$. In a diagram, for $i = 2$, it looks like this:
\[ \UseTips  \xymatrix @C=7ex @R=8ex {
\cdots 
\ar[r]
&
M^j
\ar[r]^{\d}
\ar[drr]^(0.6){\phi^j}
&
M^{j+1}
\ar[r]^{\d}
\ar[drr]^(0.6){\phi^{j+1}}
&
M^{j + 2}
\ar[r]^{\d}
&
M^{j + 3}
\ar[r]
&
\cdots
\\
\cdots 
\ar[r]
&
N^j
\ar[r]^{\d}
&
N^{j + 1}
\ar[r]^{\d}
&
N^{j + 2}
\ar[r]^{\d}
&
N^{j + 3}
\ar[r]
&
\cdots 
} \]
Warning: since $\phi$ does not have to commute with the differentials, this is 
usually not a commutative diagram! 

For a triple of complexes $M_0, M_1, M_2$ and degrees $i_0, i_1$ there are 
$\K$-linear homomorphisms
\[ \begin{aligned}
& 
\opn{Hom}_{\cat{M}}(M_1, M_2)^{i_1} \otimes_{} 
\opn{Hom}_{\cat{M}}(M_0, M_1)^{i_0} \xar{} 
\opn{Hom}_{\cat{M}}(M_0, M_2)^{i_0 + i_1} ,
\\
&
\phi_1 \ot \phi_0 \mapsto \phi_1 \circ \phi_0 . 
\end{aligned} \]

\begin{lem} \label{1213}
The composition homomorphism 
\[ \opn{Hom}_{\cat{M}}(M_1, M_2) \otimes_{} 
\opn{Hom}_{\cat{M}}(M_0, M_1) \xar{} 
\opn{Hom}_{\cat{M}}(M_0, M_2) \]
is a strict homomorphism of DG $\K$-modules. 
\end{lem}

\begin{exer}
Prove the lemma. 
\end{exer}

The lemma justifies the next definition.

\begin{dfn}
Let $\dcat{C}(\cat{M})$% 
\index{1-C(M)@$\dcat{C}(\cat{M})$}
be the DG $\K$-linear category whose objects are the 
complexes in $\cat{M}$, and whose morphism DG $\K$-modules are 
$\opn{Hom}_{\cat{M}}(M, N)$, from formulas (\ref{eqn:1700}) and 
(\ref{eqn:1130}).  
\end{dfn}

It is clear, from comparing formulas (\ref{eqn:1130}) and (\ref{eqn:1254}), 
that the strict morphisms of complexes defined at the top of this subsection 
are the same as those from Definition \ref{dfn:1212}(1). In other words, 
$\opn{Str}(\dcat{C}(\cat{M})) = \dcat{C}_{\mrm{str}}(\cat{M})$.

\begin{rem}
A possible ambiguity could arise in the meaning of 
$\opn{Hom}_{\cat{M}}(M, N)$ 
if $M, N \in \cat{M}$. Does it mean the $\K$-module of morphisms in the 
category $\cat{M}$~? Or, if we view $M$ and $N$ as complexes by the canonical
embedding $\cat{M} \subseteq \dcat{C}(\cat{M})$, does 
$\opn{Hom}_{\cat{M}}(M, N)$
mean the complex of morphisms in the DG category $\dcat{C}(\cat{M})$~? 
It turns out that there is no actual difficulty: since the complex of 
$\K$-modules $\opn{Hom}_{\cat{M}}(M, N)$ is concentrated in degree $0$, 
we may view it as a single $\K$-module, and this is precisely the $\K$-module 
of 
morphisms in the category $\cat{M}$.
\end{rem}

When $\cat{M} = \cat{Mod} A$ for a ring $A$, there is no essential distinction 
between complexes and DG modules. The next proposition is the DG version of 
Exercise \ref{exer:4112}. 

\begin{prop} \label{prop:1065}
Let $A$ be a ring. Given a complex 
$M \in \dcat{C}(\cat{Mod} A)$, with notation as in \tup{(\ref{eqn:1065})}, 
define the DG $A$-module 
$F(M) := \bigoplus_{i \in \Z} M^i$,
with differential $\d := \sum_{i \in \Z} \d_M^i$. Then the functor 
$F : \dcat{C}(\cat{Mod} A) \to \cat{DGMod} A$
is an isomorphism of DG categories. 
\end{prop}

\begin{exer}
Prove this proposition. (Hint: choose good notation.)
\end{exer}




\mysubsection{The Long Exact Cohomology Sequence}
\label{subsec:long-cohom-seq}

As in the previous subsection, $\cat{M}$ is a $\K$-linear 
abelian category. Here we give a detailed proof of the existence and 
functoriality of the long exact cohomology sequence, using our {\em sheaf 
tricks} from Subsection \ref{subsec:sheaf-tricks}. This allows us to avoid the 
Freyd-Mitchell Theorem. 

Consider a short exact sequence
\begin{equation} \label{eqn:3640}
\bsym{E} = \bigl( 0 \to L \xar{\phi} M \xar{\psi} N \to 0 \bigr)
\end{equation}
in the abelian category $\dcat{C}_{\mrm{str}}(\cat{M})$. 
This means that $\phi$ and $\psi$ are morphisms in 
$\dcat{C}_{\mrm{str}}(\cat{M})$, and in each degree $i$ there is a short 
exact sequence 
\[ 0 \to L^i \xar{\phi^i} M^i \xar{\psi^i} N^i \to 0 \]
in $\cat{M}$. The next definitions and lemmas refer to the short exact 
sequence $\bsym{E}$. In them we shall make use the notation from 
Subsection \ref{subsec:sheaf-tricks}.

Let us denote by
$\pi_N : \opn{Z}^i(N) \surj \opn{H}^i(N)$
the canonical epimorphism in $\cat{M}$; and likewise for $L$ and $M$. 

\begin{dfn} \label{dfn:3645}
Let $\bar{n} \in \Ga(U, \opn{H}^i(N))$ for some $U \in \cat{M}$,
and let $V \surj U$ be a covering of $U$. 
By a {\em connecting triple} for $\bar{n}$ over $V$
we mean a triple $(n, m, l)$ of sections
$n \in \Ga(V, \opn{Z}^i(N))$, 
$m \in \Ga(V, M^i)$ and $l \in \Ga(V, \opn{Z}^{i + 1}(L))$,
such that
$\pi_N(n) = \bar{n} \in \Ga(V, \opn{H}^i(N))$, 
$\psi(m) = n \in \Ga(V, N^i)$
and
$\phi(l) = \d(m) \in  \Ga(V, M^{i + 1})$.
\end{dfn}

Here is the relevant diagram in $\cat{M}$~:
\[ \UseTips \xymatrix @C=6ex @R=6ex {
\opn{H}^{i + 1}(L)
&
\opn{Z}^{i + 1}(L)
\ar@{->>}[l]_{\pi_L}
\ar[r]^{\phi}
&
M^{i + 1}
&
\\
&
&
M^{i}
\ar[u]^{\d}
\ar[r]^{\psi}
&
\opn{Z}^{i}(N)
\ar@{->>}[r]^{\pi_N}
&
\opn{H}^{i}(N)
} \]

It is possible that some coverings do not admit a connecting triple. But some 
do:

\begin{lem} \label{lem:3650}
Let $\bar{n} \in \Ga(U, \opn{H}^i(N))$. 
There exists a covering $\rho : V \surj U$ that admits a connecting triple 
$(n, m, l)$ for $\bar{n}$.
\end{lem}

\begin{proof}
Because $\pi_N : \opn{Z}^{i}(N) \to \opn{H}^{i}(N)$ is an epimorphism, by 
the first sheaf trick there is a covering $V' \surj U$ with 
a section $n \in \Ga(V', \opn{Z}^{i}(N))$
such that $\pi_N(n) = \bar{n}$. 

Because $\psi : M^i \to N^i$ is an epimorphism, there is a covering 
$V'' \surj V'$ with 
a section $m \in \Ga(V'', M^i)$ such that $\psi(m) = n$. 

Consider the section $\d(m) \in \Ga(V'', M^{i + 1})$. 
We have 
$\psi(\d(m)) = \d(\psi(m)) = \d(n) = 0$.
By exactness at $M^{i + 1}$, there is a covering $V \surj V''$ with 
a section $l \in \Ga(V, L^{i + 1})$ such that $\phi(l) = \d(m)$
in $\Ga(V, M^{i + 1})$. 

Now 
$\phi(\d(l)) = \d(\phi(l)) = \d(\d(m))) = 0$.
Because $\phi$ is a monomorphism, it follows that $\d(l) = 0$. Thus 
$l \in \Ga(V, \opn{Z}^{i + 1}(L))$. 
\end{proof}

\begin{lem} \label{lem:3651} 
Let $\bar{n} \in \Ga(U, \opn{H}^i(N))$. Suppose $V \surj U$
is a covering that admits connecting triples 
$(n, m, l)$ and $(n', m', l')$ for $\bar{n}$.
Then 
$\pi_L(l') = \pi_L(l)$ in $\Ga(V, \opn{H}^{i + 1}(L))$.
\end{lem}

\begin{proof}
We know that $\pi_N(n) = \bar{n} = \pi_N(n')$ in $\Ga(V, \opn{H}^{i}(N))$. 
Thus 
$n - n' \in \opn{Ker}(\pi_N) \sub \Ga(V, \opn{Z}^{i}(N))$.
Looking at the  exact sequence 
\[ N^{i - 1} \xar{\d} \opn{Z}^{i}(N) \xar{\pi_N} \opn{H}^{i}(N) \to 0 \]
and invoking the first sheaf trick, we see that there's a covering 
$W \surj V$ with a section $\til{n} \in \Ga(W, N^{i - 1})$
such that $\d(\til{n}) = n - n'$. 

Again invoking the first sheaf trick, there's a covering 
$W' \surj W$ with a section $\til{m} \in \Ga(W', M^{i - 1})$
such that $\psi(\til{m}) = \til{n}$. 

Consider the section 
$(m - m') - \d(\til{m}) \in \Ga(W', M^{i})$.
It satisfies 
\[ \psi \bigl( (m - m') - \d(\til{m}) \bigr) = 
(n - n') - \d(\til{n}) = 0 . \]
By the first sheaf trick, there's a covering 
$W'' \surj W'$ with a section $\til{l} \in \Ga(W'', L^{i})$
such that 
$\phi(\til{l}) = (m - m') - \d(\til{m})$.
Now 
\[ \phi \bigl( \d(\til{l}) - (l - l') \bigr) = 
\d(m - m') - \phi(l - l') = 0 . \]
Because $\phi$ is a monomorphism, it follows that 
$l - l' = \d(\til{l})$. 
Therefore $\pi_L(l) = \pi_L(l')$ in $\Ga(W'', \opn{H}^{i + 1}(L))$. 
But $W'' \surj V$ is a covering, so 
$\pi_L(l) = \pi_L(l')$ in $\Ga(V, \opn{H}^{i + 1}(L))$. 
\end{proof}

Lemma \ref{lem:3651} justifies the next definition. 

\begin{dfn} \label{dfn:3650}
Let $\bar{n} \in \Ga(U, \opn{H}^i(N))$. Suppose $V \surj U$ is a covering 
that admits a connecting triple. We define 
$\de_V(\bar{n}) \in \Ga(V, \opn{H}^{i + 1}(L))$
to be the unique section such that 
$\de_V(\bar{n}) = \pi_L(l)$ for every 
connecting triple $(n, m, l)$ for $\bar{n}$ over $V$.
\end{dfn}

\begin{lem} \label{lem:3640}
Let $\bar{n} \in \Ga(U, \opn{H}^i(N))$, and let  
\[ \UseTips \xymatrix @C=6ex @R=6ex {
V'
\ar@{->>}[r]^{\rho'}
\ar[d]_{\tau}
&
U'
\ar[d]^{\si}
\\
V
\ar@{->>}[r]^{\rho}
&
U
} \]
be a commutative diagram in $\cat{M}$, in which the horizontal arrows are 
coverings. Suppose $V$ admits a connecting triple for $\bar{n}$.
Then $V'$ admits a connecting triple for 
$\bar{n}' := \si^*(\bar{n})\in \Ga(U', \opn{H}^i(N))$,
and there is equality 
$\de_{V'}(\bar{n}') = \tau^*(\de_V(\bar{n}))$
in $\Ga(V', \opn{H}^{i + 1}(L))$.
\end{lem}

\begin{proof}
If $(n, m, l)$ is a connecting triple for $\bar{n}$ over $V$, then 
$\bigl( \tau^*(n), \tau^*(m), \tau^*(l) \bigr)$
is a connecting triple for $\bar{n}'$ over $V'$. Thus
$\de_{V'}(\bar{n}') = \pi_L(\tau^*(l)) = \tau^*(\pi_L(l)) = 
\tau^*(\de_V(\bar{n}))$.
\end{proof}

\begin{lem} \label{lem:3653}
Let $\bar{n} \in \Ga(U, \opn{H}^i(N))$. Suppose $\rho : V \surj U$
is a covering that admits a connecting triple for $\bar{n}$.
Then $\de_V(\bar{n})$ lies in the subgroup 
$\Ga(U, \opn{H}^{i + 1}(L))$ of $\Ga(V, \opn{H}^{i + 1}(L))$.
\end{lem}

\begin{proof}
This is a ``descent'' argument, using the second sheaf trick.
Let $V_1, V_2 := V$ and  $\rho_1, \rho_2 := \rho$.
In the notation of Proposition \ref{prop:3636}, applied to the object 
$Q := \opn{H}^{i + 1}(L) \in \cat{M}$, there is an exact sequence 
\[ 0 \to \Ga(U, Q) \xar{(\rho_1^*, \rho_2^*)} \Ga(V_1, Q) \times \Ga(V_2, Q)
\xar{(\si_1^*, - \si_2^*)} \Ga(W, Q) .  \]
Here 
\[ \rho_1 \circ \si_1 = \rho_2 \circ \si_2 : 
W = V_1 \times_{U} V_2 \to U \]
is a covering too. 
According to Lemma \ref{lem:3640}, the section 
\[ \bigl( \de_{V_1}(\bar{n}), \de_{V_2}(\bar{n}) \bigr) \in 
\Ga(V_1, Q) \times \Ga(V_2, Q) \]
satisfies 
\[ (\si_1^*, - \si_2^*)\bigl( \de_{V_1}(\bar{n}), \de_{V_2}(\bar{n}) \bigr) =
\de_W(\bar{n}) - \de_W(\bar{n}) = 0 . \]
Hence $\de_V(\bar{n})$ is in the image of $\rho^*$. 
\end{proof}

\begin{lem} \label{lem:3641}
There is a unique morphism 
$\de_{\bsym{E}}^i : \opn{H}^{i}(N) \to \opn{H}^{i + 1}(L)$
in $\cat{M}$ with the following property\tup{:}
\begin{itemize}
\item[($\dag$)] For every $U \in \cat{M}$, every section 
$\bar{n} \in \Ga(U, \opn{H}^i(N))$, and every
covering $V \surj U$ that admits a connecting triple for $\bar{n}$, there is 
equality 
$\de_{\bsym{E}}^i(\bar{n}) = \de_V(\bar{n})$
in $\Ga(U, \opn{H}^{i + 1}(L))$.
\end{itemize}
\end{lem}

\begin{proof}
We look at the universal section of $\opn{H}^{i}(N)$, namely we take 
$U_0 := \opn{H}^{i}(N)$ and 
$\bar{n}_0 := \opn{id}_{U_0} \in \Ga(U_0, \opn{H}^i(N))$.
Let $\rho_0 : V_0 \surj U_0$ be a covering that admits a connecting triple for 
$\bar{n}_0$; see Lemma \ref{lem:3650}. According to Lemma \ref{lem:3653} we 
know that 
\[ \de_{V_0}(\bar{n}_0) \in \Ga(U_0, \opn{H}^{i + 1}(L)) =
\opn{Hom}_{\cat{M}} \bigl( \opn{H}^i(N), \opn{H}^{i + 1}(L) \bigr) . \]
Define 
$\de_{\bsym{E}}^i := \de_{V_0}(\bar{n}_0)$.

We need to prove that $\de_{\bsym{E}}^i$ has property ($\dag$). 
So let $U \in \cat{M}$, let $\bar{n} \in \Ga(U, \opn{H}^i(N))$, and let
$\rho : V \surj U$ a covering that admits a connecting triple for $\bar{n}$.
Define the morphism $\si := \bar{n} : U \to U_0$. 
We get a morphism 
$\si \circ \rho : V \to U_0$. 
Define $V' := V \times_{U_0} V_0$. By Lemma \ref{lem:3630}(2) the induced 
morphism $V' \to V$ is a covering. Therefore, by composing it with $\rho$  we 
get a covering $\rho' : V' \surj U$. 

Consider the commutative diagram 
\[ \UseTips \xymatrix @C=6ex @R=6ex {
V'
\ar@{->>}[r]^{\rho'}
\ar[d]_{\tau}
&
U
\ar[d]^{\si}
\\
V_0
\ar@{->>}[r]^{\rho_0}
&
U_0
} \]
in $\cat{M}$. The horizontal arrows are coverings. Now 
$\bar{n} = \si = \si \circ \opn{id}_{U_0} = \si^*(\bar{n}_0)$.
So, according to Lemma \ref{lem:3640}, there is equality
\[ \de_{V'}(\bar{n}) = \tau^*(\de_{V_0}(\bar{n}_0)) =
\tau^*(\de_{\bsym{E}}^i \circ \rho_0) = 
\de_{\bsym{E}}^i \circ \rho_0 \circ \tau = 
\de_{\bsym{E}}^i \circ \si \circ \rho' = \de_{\bsym{E}}^i(\bar{n}) \]
in $\Ga(V', \opn{H}^{i + 1}(L))$. 
But $\rho' : V' \surj U$ is a covering, and hence there is equality 
$\de_{V'}(\bar{n}) = \de_{\bsym{E}}^i(\bar{n})$
in $\Ga(U, \opn{H}^{i + 1}(L))$.
\end{proof}

\begin{dfn} \label{dfn:3640}
The morphism 
$\de_{\bsym{E}}^i : \opn{H}^{i}(N) \to \opn{H}^{i + 1}(L)$
is called the {\em $i$-th connecting morphism}
\index{Connecting morphism}
 of the short exact sequence
$\bsym{E}$ from (\ref{eqn:3640}). 
\end{dfn}

\begin{thm}[Long Exact Sequence in Cohomology] \label{thm:3640}
\index{Long exact sequence in cohomology}
Let $\cat{M}$ be an abelian category. 
Given a short exact sequence 
\[ \bsym{E} = \bigl( 0 \to L \xar{\phi} M \xar{\psi} N \to 0 \bigr) \]
in $\dcat{C}_{\mrm{str}}(\cat{M})$, 
the sequence 
\[ \cdots \to \opn{H}^{i}(L) \xar{\opn{H}^{i}(\phi)}
\opn{H}^{i}(M) \xar{\opn{H}^{i}(\psi)}
\opn{H}^{i}(N) \xar{\de_{\bsym{E}}^i} \opn{H}^{i + 1}(L) \to \cdots \]
in $\cat{M}$ is exact. 
\end{thm}

\begin{proof}
We need to check exactness at $\opn{H}^{i}(N)$, $\opn{H}^{i + 1}(L)$  and
$\opn{H}^{i}(M)$. This will be done using the first sheaf trick (Proposition 
\ref{prop:3635}) and the second sheaf trick (Proposition \ref{prop:3636}). 

\medskip \noindent 
$\triangleright$  Step 1. 
Exactness at $\opn{H}^{i}(N)$. This is done in two substeps. 
For simplification we shall write $\bar{\psi} := \opn{H}^{i}(\psi)$. 

\medskip \noindent 
$\triangleright$  $\triangleright$ Substep 1.a.
We start by proving that 
$\de_{\bsym{E}}^i \circ \bar{\psi} = 0$.
According to Proposition \ref{prop:3635}(1) it suffices 
to show that given an element 
$\bar{m} \in \Ga(U, \opn{H}^{i}(N))$, 
with image 
$\bar{n} := \bar{\psi}(\bar{m}) \in \Ga(U, \opn{H}^{i}(N))$, 
the element 
$\de_{\bsym{E}}^i(\bar{n})$ is zero. 

By Proposition \ref{prop:3635}(3) there is a covering $V \surj U$, and an 
element $m \in \Ga(V, \opn{Z}^{i}(M))$, 
such that 
$\pi_M(m) = \bar{m} \in  \Ga(V, \opn{H}^{i}(M))$. 
Letting $n := \psi(m) \in  \Ga(V, \opn{Z}^{i}(N))$, 
we see that $\pi_N(n) = \bar{n} \in  \Ga(V, \opn{H}^{i}(N))$.
Now $\d_M(m) = 0$, and hence, taking $l := 0$, the triple $(n, m, l)$
is a connecting triple for $\bar{n}$ over $V$. Therefore
$\de_{\bsym{E}}^i(\bar{n}) = \pi_L(l) = 0$.

\medskip \noindent 
$\triangleright$  $\triangleright$  Substep 1.b.
Now we prove that the morphism 
$\bar{\psi} : \opn{H}^{i}(M) \to \opn{Ker}(\de_{\bsym{E}}^i)$
is an epimorphism. 
We shall use Proposition \ref{prop:3635}(3). Let 
$\bar{n} \in \Ga(U, \opn{Ker}(\de_{\bsym{E}}^i))$; 
so $\bar{n} \in \Ga(U, \opn{H}^{i}(N))$ and
$\de_{\bsym{E}}^i(\bar{n}) = 0$. 
It suffices to find a covering $V' \surj U$ and a section 
$\bar{m} \in \Ga(V', \opn{H}^{i}(M))$ such that 
$\bar{\psi}(\bar{m}) = \bar{n}$ in $\Ga(V', \opn{H}^{i}(N))$. 

Take a covering $V \surj U$  that admits a connecting triple $(n, m, l)$ for 
$\bar{n}$. Then 
$\pi_L(l) = \de_{V}(\bar{n}) = \de_{\bsym{E}}^i(\bar{n}) = 0$
in $\Ga(V, \opn{H}^{i + 1}(L))$. 
By Proposition \ref{prop:3635}(3) there is some covering 
$V' \surj V$ and a section $l' \in \Ga(V', L^{i})$ s.t.\ 
$\d_L(l') = l$ in $\Ga(V', L^{i + 1})$. 
Define $m' := \phi(l') \in \Ga(V', M^{i})$. 
Now 
$\d_M(m - m') = \d_M(m) - \d_M(m') = \phi(l) - \phi(l) = 0$,
so $m - m'$ belongs to $\Ga(V', \opn{Z}^{i}(M))$. 
Because
$\psi(m - m') = \psi(m) = n$,
the cohomology class 
$\bar{m} := \pi_M(m - m')$ in $\Ga(V', \opn{H}^{i}(M))$
satisfies $\bar{\psi}(\bar{m}) = \bar{n}$
in $\Ga(V', \opn{H}^{i}(N))$. 


\medskip \noindent 
$\triangleright$ Step 2. Exactness at $\opn{H}^{i + 1}(L)$. This is also done 
in two substeps. We will use the abbreviation 
$\bar{\phi} := \opn{H}^{i + 1}(\phi)$. 

\medskip \noindent 
$\triangleright$  $\triangleright$  Substep 2.a.
We start by proving that 
$\bar{\phi} \circ \de_{\bsym{E}}^i  = 0$.
Take $\bar{n} \in \Ga(U, \opn{H}^{i}(N))$, and let 
$\bar{l} := \de_{\bsym{E}}^i(\bar{n})$. We need to prove that 
$\bar{\phi}(\bar{l}) = 0$. 
Let $V \surj U$ be a covering that admits a connecting triple $(n, m, l)$ for 
$\bar{n}$. This says that $\bar{l} = \pi_L(l)$, and hence 
$\bar{\phi}(\bar{l}) = \pi_L(\phi(l)) = \pi_M(\d_M(m)) = 0$.

\medskip \noindent 
$\triangleright$  $\triangleright$  Substep 2.b.
Now we prove that 
$\de_{\bsym{E}}^i : \opn{H}^{i}(N) \to \opn{Ker}(\bar{\phi})$
is an epimorphism. 
Let $\bar{l} \in \Ga(U, \opn{H}^{i + 1}(L))$ be such that 
$\bar{\phi}(\bar{l}) = 0$. 
Take a covering $V \surj U$ for which there's a section 
$l \in \Ga(V, \opn{Z}^{i + 1}(L))$
s.t.\ $\pi_L(l) = \bar{l}$. The section 
$\phi(l) \in \Ga(V, \opn{Z}^{i + 1}(M))$
satisfies 
$\pi_M(\phi(l)) = \bar{\phi}(\bar{l}) = 0$
in $\Ga(V, \opn{H}^{i + 1}(M))$, and therefore there is some covering 
$V' \surj V$ and a section 
$m \in \Ga(V', M^{i})$ s.t.\ 
$\d_M(m) = \phi(l)$ in $\Ga(V', M^{i + 1})$.

Define $n := \psi(m) \in \Ga(V', N^{i})$. Then 
$\d_N(n) = \d_N(\psi(m)) = \psi(\d_M(m)) = \psi(\phi(l)) = 0$,
so in fact $n \in \Ga(V', \opn{Z}^{i}(N))$. 
Let $\bar{n} := \pi_N(n) \in \Ga(V', \opn{H}^{i}(N))$.
We see that $(n, m, l)$ is a connecting triple for $\bar{n}$ over $V'$. 
Therefore 
$\bar{l} = \pi_L(l) = \de_{V'}(\bar{n}) = \de_{\bsym{E}}^i(\bar{n})$
in $\Ga(V', \opn{H}^{i + 1}(L))$. 

\medskip \noindent 
$\triangleright$ Step 3. Exactness at $\opn{H}^{i}(M)$: 
let's use the abbreviations $\bar{\phi} := \opn{H}^{i}(\phi)$ and 
$\bar{\psi} := \opn{H}^{i}(\psi)$. 
It is clear that $\bar{\psi} \circ \bar{\phi} = 0$. 
It remains to prove that 
$\bar{\phi} : \opn{H}^{i}(L) \to \opn{Ker}(\bar{\psi})$
is an epimorphism. 

Let $\bar{m} \in \Ga(U, \opn{H}^{i}(M))$ be such that 
$\bar{\psi}(\bar{m}) = 0$. 
We are going to find a covering $W \to U$ and a section 
$\bar{l} \in \Ga(W, \opn{H}^{i}(L))$
such that $\bar{\phi}(\bar{l}) = \bar{m}$ in 
$\Ga(W, \opn{H}^{i}(M))$.

Take a covering $V \surj U$ for which there's a section 
$m \in \Ga(V, \opn{Z}^{i}(M))$
s.t.\ $\pi_M(m) = \bar{m}$. 
The section $n := \psi(m) \in \Ga(V, \opn{Z}^{i}(N))$
satisfies
$\pi_N(n) = \bar{\psi}(\bar{m}) = 0$
in $\Ga(V, \opn{H}^{i}(N))$, and therefore there is some covering 
$V' \surj V$ and a section 
$n' \in \Ga(V', N^{i - 1})$ s.t.\ 
$\d(n') = n$ in $\Ga(V', N^{i})$.

Take a covering $V'' \surj V'$ for which there is a section 
$m'' \in \Ga(V'', M^{i - 1})$ s.t.\ 
$\psi(m'') = n'$. The section 
$m - \d_M(m'') \in \Ga(V'', \opn{Z}^{i}(M))$
satisfies $\pi_M(m - \d(m'')) = \bar{m}$ in 
$\Ga(V'', \opn{H}^{i}(M)$, and also 
$\psi(m - \d_M(m'')) = \psi(m) - \psi(\d_M(m'')) = n - n = 0$.

There is a covering $W \surj V''$ with a section 
$l \in \Ga(W, L^i)$ s.t.\ $\phi(l) = m - \d_M(m'')$. 
Because $\phi$ is a monomorphism, and $m - \d_M(m'')$ is a cocycle, it follows 
that the section $l$ belongs to 
$\Ga(W, \opn{Z}^{i}(L))$. Its cohomology class 
$\bar{l} := \pi_L(l)$ in $\Ga(W, \opn{H}^{i}(L))$
satisfies 
$\bar{\phi}(\bar{l}) = \pi_M(\phi(l)) = \pi_M(m - \d_M(m'')) = \bar{m}$
in $\Ga(W, \opn{H}^{i}(M))$.  
\end{proof}

\begin{prop} \label{prop:3640}
Let $\chi : \bsym{E} \to \bsym{E}'$ be a morphism of short exact sequences in 
$\dcat{C}_{\mrm{str}}(\cat{M})$. Namely 
$\chi = (\chi_L, \chi_M, \chi_N)$ in the commutative diagram with exact rows
\[ \UseTips \xymatrix @C=8ex @R=6ex {
0
\ar[r]
&
L
\ar[r]^{\phi}
\ar[d]_{\chi_L}
&
M
\ar[r]^{\psi}
\ar[d]_{\chi_M}
&
N
\ar[r]
\ar[d]_{\chi_N}
&
0
\\
0
\ar[r]
&
L'
\ar[r]^{\phi'}
&
M'
\ar[r]^{\psi'}
&
N'
\ar[r]
&
0
} \]
in $\dcat{C}_{\mrm{str}}(\cat{M})$. Then for every $i$ the diagram 
\[ \UseTips \xymatrix @C=8ex @R=6ex {
\opn{H}^{i}(N) 
\ar[r]^{\de_{\bsym{E}}^i}
\ar[d]_{\opn{H}^{i}(\chi_N)}
&
\opn{H}^{i + 1}(L)
\ar[d]^{\opn{H}^{i + 1}(\chi_L)}
\\
\opn{H}^{i}(N') 
\ar[r]^{\de_{\bsym{E}'}^i}
&
\opn{H}^{i + 1}(L')
} \]
is commutative. 
\end{prop}

\begin{exer} \label{exer:3640}
Prove Proposition \ref{prop:3640}. (Hint: study the proof of the previous 
theorem.)
\end{exer}
 




\mysubsection{The DG Category 
\texorpdfstring{$\dcat{C}(A, \cat{M})$} {C(A,M)}} 
\label{subsec:DGModinM}

We now combine material from previous subsections. The concept introduced in 
the definition below is new. It is the DG version of Definition 
\ref{dfn:1702}. 

Conventions \ref{conv:2490}, \ref{conv:4626}, \ref{conv:4627} and 
\ref{conv:4628} are in place. Thus all linear structures and operations 
are $\K$-linear, and rings and bimodules are central over $\K$. 

Recall that given an  abelian category $\cat{M}$, the category of 
complexes $\dcat{C}(\cat{M})$ is a DG category. For a complex 
$M \in \dcat{C}(\cat{M})$ we have the DG ring 
$\opn{End}_{\cat{M}}(M) = \opn{Hom}_{\cat{M}}(M, M)$.
The multiplication in this ring is composition. 

\begin{dfn} \label{dfn:1130}
Let $\cat{M}$ be an abelian category, and let $A$ be a DG ring. A {\em DG 
$A$-module in  $\cat{M}$}
\index{Abelian category! DG $A$-module in}
is an object 
$M \in \dcat{C}(\cat{M})$, together with a DG ring homomorphism 
$f : A \to \opn{End}_{\cat{M}}(M)$. 
\end{dfn}

If $M$ is a DG $A$-module in $\cat{M}$, then after forgetting the 
differentials, $M$ becomes a graded $A$-module in $\cat{M}$.

\begin{dfn} \label{dfn:1131}
\index{Differential graded module! Hom between {\indash}s}
Let $\cat{M}$ be an abelian category, let $A$ be a DG ring, and let $M, N$ 
be DG $A$-modules in $\cat{M}$. In Definition \ref{dfn:1703} we introduced the 
graded $\K$-module $\opn{Hom}_{A, \cat{M}}(M, N)$. This is made into a DG 
$\K$-module with differential 
\[ \d(\phi) := \d_N \circ \phi - (-1)^i \cd \phi \circ \d_M \]
for $\phi \in \opn{Hom}_{A, \cat{M}}(M, N)^i$. 
\end{dfn}

When we have to be specific, we denote the differential of 
$\opn{Hom}_{A, \cat{M}}(M, N)$ by 
$\d_{\mrm{Hom}}$, $\d_{A, \cat{M}}$, or some similar expression. 

As we have seen before (in Lemmas \ref{1213} and \ref{lem:1211}), given 
morphisms
\[ \phi_k \in \opn{Hom}_{A, \cat{M}}(M_{k}, M_{k + 1})^{i_k} \]
for $k \in \{ 0,  1 \}$, we have  
\[ \phi_1 \circ \phi_0 \in 
\opn{Hom}_{A, \cat{M}}(M_0, M_2)^{i_0 + i_1} , \]
and 
\[ \d(\phi_1 \circ \phi_0) = 
\d(\phi_1) \circ \phi_0 + (-1)^{i_1} \cd \phi_1 \circ \d(\phi_0) . \]
Also the identity automorphism $\opn{id}_M$ belongs to 
$\opn{Hom}_{A, \cat{M}}(M, M)^{0}$, and $\d(\opn{id}_M) = 0$.
Therefore the next definition is legitimate. 

\begin{dfn} \label{dfn:1213}
\index{Differential graded category! of DG $A$-modules in $\cat{M}$}
Let $\cat{M}$ be an abelian category, and let $A$ be a DG ring. The 
DG category of DG $A$-modules in $\cat{M}$ is 
denoted by $\dcat{C}(A, \cat{M})$%
\index{1-C(A,M)@$\dcat{C}(A, \cat{M})$}.
The morphism DG $\K$-modules are
\[ \opn{Hom}_{\dcat{C}(A, \cat{M})}(M_0, M_1) :=
\opn{Hom}_{A, \cat{M}}(M_0, M_1) \]
from Definition \ref{dfn:1131}.
The composition is that of $\dcat{C}(\cat{M})$.
\end{dfn}

Notice that forgetting the action of $A$ is a faithful DG
functor $\dcat{C}(A, \cat{M}) \to \dcat{C}(\cat{M})$.

\begin{exa} \label{exa:1820}
If $A = \K$, then $\dcat{C}(A, \cat{M}) = \dcat{C}(\cat{M})$; 
and if $\cat{M} = \cat{Mod} \K$, then 
$\dcat{C}(A, \cat{M}) = \dcat{C}(A) = \cat{DGMod} A$.
\end{exa}

\begin{dfn} \label{dfn:1132}
In the situation of Definition \ref{dfn:1213}:
\begin{enumerate}
\item The strict category of $\dcat{C}(A, \cat{M})$%
\index{Strict category! of DG $A$-modules in $\cat{M}$}
(see Definition \ref{dfn:1212}(1)) is denoted by \lb 
$\dcat{C}_{\mrm{str}}(A, \cat{M})$%
\index{1-Cstr(A,M)@$\dcat{C}_{\mrm{str}}(A, \cat{M})$}.

\item The homotopy category of $\dcat{C}(A, \cat{M})$%
\index{Homotopy category! of DG $A$-modules in $\cat{M}$}%
\index{1-K(A,M)@$\dcat{K}(A, \cat{M})$}
(see Definition \ref{dfn:1212}(2)) is denoted by 
$\dcat{K}(A, \cat{M})$.
\end{enumerate}
\end{dfn}

To say this in words: a morphism $\phi : M \to N$ in 
$\dcat{C}(A, \cat{M})$ is strict if and only if it has degree $0$ and 
$\phi \circ \d_M = \d_N \circ \phi$. 
The morphisms in $\dcat{K}(A, \cat{M})$ are the homotopy classes of strict 
morphisms. 

An additive functor $F : \cat{M} \to \cat{N}$ between abelian categories is 
called {\em faithfully exact} if for every sequence $\bsym{E}$ in 
$\dcat{C}_{\mrm{str}}(A, \cat{M})$, the sequence $\bsym{E}$ is exact if and 
only if the sequence $F(\bsym{E})$ in $\cat{N}$ is exact.

\begin{prop} \label{prop:4130}
The category $\dcat{C}_{\mrm{str}}(A, \cat{M})$ is a $\K$-linear abelian 
category, and the forgetful functors 
\[ \dcat{C}_{\mrm{str}}(A, \cat{M}) \xar{} \dcat{C}_{\mrm{str}}(\cat{M}) 
\xar{\opn{Und}} \dcat{G}_{\mrm{str}}(\cat{M}) \]
are faithfully exact. 
\end{prop}

\begin{exer} \label{exer:4145}
Prove the proposition above. 
\end{exer}

Like Definition \ref{dfn:2993}, given a DG module 
$M \in \dcat{C}(A, \cat{M})$ and an 
integer $i$, we can consider the objects of degree $i$ cocycles 
$\opn{Z}^i(M)$,  decocycles $\opn{Y}^i(M)$, coboundaries $\opn{B}^i(M)$ and 
cohomology $\opn{H}^i(M)$. These are all objects of $\cat{M}$. As $M$ varies we 
get linear functors 
\begin{equation} \label{eqn:4145}
\opn{Z}, \opn{Y}, \opn{B}, \opn{H} : \dcat{C}_{\mrm{str}}(A, \cat{M}) \to 
\dcat{G}_{\mrm{str}}(\cat{M}) . 
\end{equation}
These objects are related by the following exact sequences in $\cat{M}$~:
\begin{equation} \label{eqn:4148}
0 \to \opn{Z}^i(M) \to M^i \xar{\d^i_M} M^{i + 1} , 
\end{equation}
\begin{equation} \label{eqn:4149}
M^{i - 1} \xar{\d^{i - 1}_M} M^{} \to \opn{Y}^i(M) \to 0 , 
\end{equation}
\begin{equation} \label{eqn:4146}
0 \to \opn{B}^i(M) \to \opn{Z}^i(M) \to \opn{H}^i(M) \to 0 , 
\end{equation}
\begin{equation} \label{eqn:4147}
0 \to \opn{H}^i(M) \to \opn{Y}^i(M) \to \opn{B}^{i + 1}(M) \to 0 . 
\end{equation}


\begin{rem} \label{rem:3980}
Actually, the cohomology functor $\opn{H}$ from equation (\ref{eqn:4145})
factors through the category $\dcat{G}_{\mrm{str}}(\opn{H}(A),  \cat{M})$. 
See formula (\ref{eqn:2360}). We will not need this fact. 
\end{rem}

\begin{prop} \label{prop:4145}
The functor $\opn{Z}$ in equation \tup{(\ref{eqn:4145})}
is left exact, and the functor
$\opn{Y}$ in equation \tup{(\ref{eqn:4145})}
is right exact.
\end{prop}

\begin{proof}
It is enough to consider each degree separately. And by Proposition 
\ref{prop:4130} we may ignore the DG ring $A$.
Let $F^i : \dcat{G}_{\mrm{str}}(\cat{M}) \to \cat{M}$
be the functor that sends a graded module $M$ to its degree $i$ component 
$M^i$. This is an exact functor. 

Formula (\ref{eqn:4148}) exhibits $\opn{Z}^i$ as the kernel of the homomorphism 
of functors $\d^i : F^i \to F^{i + 1}$. According to Proposition 
\ref{prop:3620}(2) the functor $\opn{Z}^i$ is left exact. 

Similarly, formula (\ref{eqn:4149}) exhibits $\opn{Y}^i$ as the cokernel of the 
homomorphism of functors 
$\d^{i - 1} : F^{i - 1} \to F^{i}$. According to Proposition 
\ref{prop:3620}(2) the functor $\opn{Y}^i$ is right exact. 
\end{proof}

\begin{prop} \label{prop:3140}
Let $\{ M_x \}_{x \in X}$ be a collection of objects of
$\dcat{C}(A, \cat{M})$, indexed by a set $X$. Assume that for every $i \in \Z$ 
the direct sum $M^i := \bigoplus_{x \in X} M^i_x$ 
exists in $\cat{M}$. Then the graded object 
$M := \{ M^i \}_{i \in \Z} \in \dcat{G}(\cat{M})$ 
has a canonical structure of DG $A$-module in $\cat{M}$, 
and $M$ is a direct sum of the collection $\{ M_x \}_{x \in X}$
in $\dcat{C}_{\mrm{str}}(A, \cat{M})$. 
\end{prop}

\begin{exer} \label{exer:3140}
Prove Proposition \ref{prop:3140}. 
(Hint: look at the proof of Theorem \ref{thm:2495}.)
\end{exer}

\begin{exa} \label{exa:2367}
Since $\dcat{M}(\K)$ has infinite direct sums, the proposition above shows that 
$\dcat{C}_{\mrm{str}}(A)$ has infinite direct sums. 
\end{exa}

\begin{prop} \label{prop:2365}
Given a DG ring $A$, let $A^{\natural}$ be the graded ring gotten by forgetting 
the differential. Consider the forgetful functor 
$\opn{Und} : \dcat{C}(A, \cat{M}) \to \dcat{G}(A^{\natural}, \cat{M})$
that forgets the differentials of DG modules, i.e.\ 
$\opn{Und}(M, \d_M) :=  M$.
\begin{enumerate}
\item $\opn{Und}$ is a fully faithful graded functor.

\item On the strict categories, 
$\opn{Str}(\opn{Und}) :  \dcat{C}_{\mrm{str}}(A, \cat{M}) \to 
\dcat{G}_{\mrm{str}}(A, \cat{M})$
is a faithfully exact additive functor.
\end{enumerate} 
\end{prop}

\begin{exer} \label{exer:2367}
Prove Proposition \ref{prop:2365}.
\end{exer}

A graded object $M = \{ M^i \}_{i \in \Z}$ in $\cat{M}$ is said to be {\em
bounded above} if the set $\{ i \mid M^i \neq 0 \}$ is bounded above. 
Likewise we define {\em bounded below} and {\em bounded} graded objects. 

\begin{dfn} \label{dfn:3220}
\index{1-C(A,M)@$\dcat{C}^{\star}(A, \cat{M})$}
\index{Boundedness! condition}
We define $\dcat{C}^-(A, \cat{M})$, $\dcat{C}^+(A, \cat{M})$ and 
$\dcat{C}^{\mrm{b}}(A, \cat{M})$ to be the full subcategories of 
$\dcat{C}(A, \cat{M})$ consisting of bounded above, bounded below and bounded
DG modules respectively. 
\end{dfn}

The symbols ``$-$'', ``$+$'' and ``$\mrm{b}$'' and 
``$\bra{\mrm{empty}}$'' (the empty symbol) are called {\em boundedness 
indicators}. We shall usually use the symbol ``$\star$'' to denote an 
unspecified boundedness indicator. Thus the expression 
$\dcat{C}^{\star}(A, \cat{M})$ can refer to any of the four full subcategories 
of $\dcat{C}(A, \cat{M})$ that are mentioned in the definition above.

\begin{rem} \label{rem:1245}
Here is a generalization of Definition \ref{dfn:1213}. Instead of a DG central
$\K$-ring $A$ we can take a small $\K$-linear DG category $\cat{A}$. 
We then define the $\K$-linear DG category 
$\dcat{C}(\cat{A}, \cat{M}) := \cat{DGFun}(\cat{A}, \dcat{C}(\cat{M}))$
as in Exercise \ref{exer:1250}.

This is indeed a generalization of Definition \ref{dfn:1213}: when $\cat{A}$ 
has a single object $x$, and we write $A := \opn{End}_{\cat{A}}(x)$, then 
the functor $M \mapsto M(x)$ is an isomorphism of DG categories
$\dcat{C}(\cat{A}, \cat{M}) \iso \dcat{C}(A, \cat{M})$. 

In the special case of $\cat{M} = \dcat{M}(\K) = \cat{Mod} \K$, the DG category 
$\dcat{C}(\cat{A}, \cat{M})$ is what Keller \cite{Kel1} calls the DG category 
of {\em DG $\cat{A}$-modules}. 

Almost everything we do in this book for $\dcat{C}(A, \cat{M})$
holds in the more general context of $\dcat{C}(\cat{A}, \cat{M})$.
However, in the more general context a lot of the intuition is lost, and some 
aspects become pretty cumbersome. This is the reason we have decided to stick 
with the less general context. 
\end{rem}

\begin{rem} \label{rem:4960}
The theory of DG categories has a generalization -- this is the theory of {\em 
$\mrm{A}_{\infty}$ categories}%
\index{Ainfty@$\mrm{A}_{\infty}$ category},
which will be outlined in this remark.
All set theoretical issues are going to be ignored.

The concept of an {\em $\mrm{A}_{\infty}$ algebra}, namely a single-object 
$\mrm{A}_{\infty}$ category, was introduced by J. Stasheff \cite{Sta1}, 
\cite{Sta2} in the 1960's, in relation with his work on {\em higher 
operations on H-spaces} in topology. 

$\mrm{A}_{\infty}$ categories are at the core of  {\em Homological Mirror 
Symmetry}, conceived by K. Fukaya and M. Kontsevich during the 1990's.
Indeed, a symplectic manifold $Y$ gives rise to the {\em Fukaya category} 
$\cat{F}(Y)$, which is an $\mrm{A}_{\infty}$ category, and is a categorification 
of {\em Floer cohomology}. See \cite{Fuk}, \cite{Kon}, \cite{KoSo}, \cite{Seid} 
and \cite{HMS}. 

Further developments on $\mrm{A}_{\infty}$ algebras and categories were made, 
among others, by T.V. Kadeishvili \cite{Kad}, K. Lef\`evre-Hasegawa \cite{Lef}, 
B. Keller \cite{Kel6}, \cite{Kel5}, and Canonaco-Ornaghi-Stellari \cite{COS}. 
For the analogous theory of {\em $\mrm{L}_{\infty}$ algebras},
generalizing DG Lie algebras, see \cite{Ye5.5}.

We shall only discuss strictly unital $\mrm{A}_{\infty}$ categories, giving 
very few details (full details can be found in \cite{Kel5} and \cite{COS}), and 
our notation deviates somewhat from the references
(because we want it to be compatible with the rest of the book). 
The base ring $\K$ is assumed to be a field. 
An $\mrm{A}_{\infty}$ category $\cat{A}$ over $\K$
has a set of objects $\opn{Ob}(\cat{A})$. For each pair of objects $x, y \in 
\opn{Ob}(\cat{A})$ there is a graded $\K$-module 
$\cat{A}(x, y) = \opn{Hom}_{\cat{A}}(x, y)$. 
For each $x \in \opn{Ob}(\cat{A})$ there is a degree 
$0$ element $\opn{id}_x \in \cat{A}(x, x)$.
For each $i \geq 1$ and each sequence 
$x_0, \ldots, x_i \in \opn{Ob}(\cat{A})$, there is a 
degree $2 - i$ homomorphism 
\begin{equation} \label{4960}
\opn{m}_i : \cat{A}(x_{i - 1}, x_{i}) \ot \cdots \ot  
\cat{A}(x_{0}, x_{1}) \to \cat{A}(x_{0}, x_{i}) 
\end{equation}
$\dcat{G}(\K)$. 
The operations $\opn{m}_i$ satisfy certain relations. 
Roughly speaking, the operation $\opn{m}_1$ is a differential, the operation 
$\opn{m}_2$ is a multiplication, the element $\opn{id}_x$ is a unit for the 
multiplication $\opn{m}_2$, the multiplication $\opn{m}_2$ is associative 
up to the hotomopy $\opn{m}_3$, etc.
The relations imply that the cohomology $\opn{H}(\cat{A})$, w.r.t.\ the 
differential $\opn{m}_1$, is a graded category, in the sense of 
Definition \ref{dfn:1692}.
If the higher (i.e.\ $i > 2$) operations $\opn{m}_i$ vanish, then $\cat{A}$ 
is a DG category, in the sense 
of Definition \ref{dfn:1063}.

Suppose $\cat{A}$ and $\cat{B}$ are $\mrm{A}_{\infty}$ categories. 
An {\em $\mrm{A}_{\infty}$ functor} $F : \cat{A} \to \cat{B}$
consists of a function
$F_{\mrm{ob}} : \opn{Ob}(\cat{A}) \to \opn{Ob}(\cat{B})$, 
and for each $i \geq 1$ and each sequence 
$x_0, \ldots, x_i \in \opn{Ob}(\cat{A})$, a degree $1 - i$ homomorphism
\begin{equation} \label{4961}
F_i : \cat{A}(x_{i - 1}, x_{i}) \ot \cdots \ot  \cat{A}(x_{0}, x_{1}) \to 
\cat{B} \bigl( F_{\mrm{ob}}(x_{0}), F_{\mrm{ob}}(x_{i}) \bigr) 
\end{equation}
in $\dcat{G}(\K)$. 
There are relations that the $F_i$ have to satisfy.
Given another $\mrm{A}_{\infty}$ functor $G : \cat{B} \to \cat{C}$, 
there is a composed $\mrm{A}_{\infty}$ functor
$G \circ F : \cat{A} \to \cat{C}$, which is hard to describe directly (it is 
usually defined in terms of {\em cocategories}, via the {\em bar 
construction}). 
We denote by $\cat{A_{\infty}Cat}$ the category whose objects are the 
$\mrm{A}_{\infty}$ categories and whose morphisms are the $\mrm{A}_{\infty}$ 
functors.  

For a pair $\cat{A}, \cat{B} \in \opn{Ob}(\cat{A_{\infty}Cat})$
we denote by  $\cat{A_{\infty}Fun}(\cat{A}, \cat{B})$ 
the set of $\mrm{A}_{\infty}$ functors from $\cat{A}$ to $\cat{B}$, so 
that \begin{equation} \label{eqn:4970}
\opn{Hom}_{\cat{A_{\infty}Cat}}(\cat{A}, \cat{B}) =
\cat{A_{\infty}Fun}(\cat{A}, \cat{B}) . 
\end{equation}
An important fact is that 
$\cat{A_{\infty}Fun}(\cat{A}, \cat{B})$
is itself the set of objects of an $\mrm{A}_{\infty}$ category, i.e.\ 
$\cat{A_{\infty}Fun}(\cat{A}, \cat{B}) \in \opn{Ob}(\cat{A_{\infty}Cat})$. 
If $\cat{B}$ is a DG category (i.e.\ its higher operations vanish), then 
$\cat{A_{\infty}Fun}(\cat{A}, \cat{B})$
is also a DG category.

We shall discuss more aspects of $\mrm{A}_{\infty}$ categories in Remarks 
\ref{rem:4970} and \ref{rem:4971}. 
\end{rem}





\mysubsection{Contravariant DG Functors} \label{subsec:contrvar-dg-func}

In this subsection we address, in a systematic fashion, the issue of reversing 
arrows in DG categories. As always, we work over a commutative base ring $\K$. 
 
\begin{dfn} \label{dfn:2495}
Let $\cat{C}$ and $\cat{D}$ be DG categories.
A {\em contravariant DG functor}%
\index{Differential graded functor! contravariant}
$F : \cat{C} \to \cat{D}$ consists of a function 
$F : \opn{Ob}(\cat{C}) \to \opn{Ob}(\cat{D})$,
and for each pair of objects $M_0, M_1 \in \opn{Ob}(\cat{C})$ a homomorphism 
\[ F : \opn{Hom}_{\cat{C}}(M_0, M_1) \to 
\opn{Hom}_{\cat{D}} \bigl( F(M_1), F(M_0) \bigr) \]
in $\dcat{C}_{\mrm{str}}(\K)$. The conditions are:
\begin{itemize}
\rmitem{a} Units: $F(\opn{id}_M) = \opn{id}_{F(M)}$. 

\rmitem{b} Graded reversed composition: suppose that for $k \in \{ 0, 1 \}$
we are given morphisms 
$\phi_k \in \opn{Hom}_{\cat{C}}(M_{k}, M_{k + 1})^{i_k}$.
Then there is equality 
\[ F(\phi_1 \circ \phi_0) = (-1)^{i_0 \cd i_1} \cd F(\phi_0) \circ F(\phi_1) \]
inside 
$\opn{Hom}_{\cat{D}} \bigl( F(M_2), F(M_0) \bigr)^{i_0 + i_1}$.  
\end{itemize}
\end{dfn}

Warning: a contravariant DG functor does not remain a contravariant functor
after the grading is forgotten; cf.\ Remark \ref{rem:2992}. 

Here is the categorical version of Definition \ref{dfn:1080}.

\begin{dfn} \label{dfn:1272}
Let $\cat{C}$ be a DG category. The 
{\em opposite DG category}% 
\index{Differential graded category! opposite category of}
$\cat{C}^{\mrm{op}}$ has the same set of objects. The morphism DG $\K$-modules 
are
\[ \opn{Hom}_{\cat{C}^{\mrm{op}}} (M_0, M_1) := 
\opn{Hom}_{\cat{C}} (M_1, M_0) . \]
The composition $\circ^{\mrm{op}}$ of $\cat{C}^{\mrm{op}}$ is the composition 
$\circ$ of $\cat{C}$, reversed and multiplied by signs: 
\[ \phi_0 \circ^{\mrm{op}} \phi_1 := 
(-1)^{i_0 \cd i_1} \cd \phi_1 \circ \phi_0  \]
for morphisms 
$\phi_k \in \opn{Hom}_{\cat{C}} (M_k, M_{k + 1})^{i_k}$.
\end{dfn}

One needs to verify that this is indeed a DG category. 
This is basically the same verification as in Exercise \ref{exer:2500}.

As before, we define the operation 
$\opn{Op} : \cat{C} \to \cat{C}^{\mrm{op}}$
to be the identity on objects, and the identity on morphisms in reversed 
order, i.e.\ 
\[ \opn{Op} = \opn{id} : \opn{Hom}_{\cat{C}} (M_0, M_1) \iso 
\opn{Hom}_{\cat{C}^{\mrm{op}}} (M_1, M_0) . \]
Note that $(\cat{C}^{\mrm{op}})^{\mrm{op}} = \cat{C}$, 
and we denote the inverse operation $\cat{C}^{\mrm{op}} \to \cat{C}$
also by $\opn{Op}$. 

\begin{prop} \label{prop:2500}
Let $\cat{C}$, $\cat{D}$ and $\cat{E}$ be DG categories.
\begin{enumerate}
\item The operations $\opn{Op} : \cat{C} \to \cat{C}^{\mrm{op}}$
and $\opn{Op} : \cat{C}^{\mrm{op}} \to \cat{C}$
are contravariant DG functors. 

\item If $F : \cat{C} \to \cat{D}$ is a contravariant DG functor, 
then the composition $F \circ \opn{Op} : \cat{C}^{\mrm{op}} \to \cat{D}$
is a DG functor; and vice versa. 

\item If $F : \cat{C} \to \cat{D}$ and $G : \cat{D} \to \cat{E}$
are contravariant DG functors, 
then the composition $G \circ F : \cat{C} \to \cat{E}$
is a  DG functor. 
\end{enumerate}
\end{prop}

\begin{exer}  \label{exer:2495}
Prove the previous proposition. 
\end{exer}

Definitions \ref{dfn:1272} and \ref{dfn:2495} make sense for graded categories, 
by forgetting differentials. Thus for graded categories $\cat{C}$ and 
$\cat{D}$ we can talk about contravariant graded functors
$\cat{C} \to \cat{D}$, and about the graded category $\cat{C}^{\mrm{op}}$.

We already met $\dcat{G}(\cat{M})$, the category of graded objects in an
abelian category $\cat{M}$; see Definition \ref{dfn:1701}. It is a 
graded category. 
Its objects are collections $M = \{ M^i \}_{i \in \Z}$ of objects 
$M^i \in \cat{M}$. 

Let $\cat{M}$ and $\cat{N}$ be abelian categories, and let 
$F : \cat{M} \to \cat{N}$ be a contravariant linear functor. 
For a graded object 
$M = \{ M^i \}_{i \in \Z} \in \dcat{G}(\cat{M})$
let us define the graded object 
\begin{equation} \label{eqn:2496}
\dcat{G}(F)(M) := \{ N^i \}_{i \in \Z} \in \dcat{G}(\cat{N}) , \quad 
N^i := F(M^{-i}) \in \cat{N} .
\end{equation}

Next consider a pair of objects $M_0, M_1 \in \dcat{G}(\cat{M})$
and a degree $i$ morphism 
$\phi : M_0 \to M_1$ in $\dcat{G}(\cat{M})$. Thus 
\[ \phi = \{ \phi^{\lms j} \}_{j \in \Z} \in 
\opn{Hom}_{\dcat{G}(\cat{M})}(M_0, M_1)^i , \]
where, as in formula (\ref{eqn:2495}), the morphism $\phi^{\lms j}$ belongs to 
$\opn{Hom}_{\cat{M}}(M_0^j, M_1^{j + i})$.
We have objects 
$N_k := \dcat{G}(F)(M_k) \in \dcat{G}(\cat{N})$,
for $k \in \{ 0, 1 \}$, defined by (\ref{eqn:2496}). Explicitly, 
$N_k = \{ N_k^i \}_{i \in \Z}$ and $N_k^i = F(M_k^{-i})$. 
For each $j \in \Z$ define the morphism 
\begin{equation} \label{eqn:2505}
\psi^{\lms j} := (-1)^{i \cd j} \cd F(\phi^{-j - i}) \in
\opn{Hom}_{\cat{N}} \bigl( N_1^j, N_0^{j + i} \bigr) .
\end{equation}
Collecting them we obtain a morphism 
\begin{equation} \label{eqn:2497}
\dcat{G}(F)(\phi) := \{ \psi^{\lms j} \}_{j \in \Z}  \in 
\opn{Hom}_{\dcat{G}(\cat{N})}(N_1, N_0)^i . 
\end{equation}

\begin{lem} \label{lem:2495}
The assignments \tup{(\ref{eqn:2496})} and \tup{(\ref{eqn:2497})}
produce a contravariant graded functor 
$\dcat{G}(F) : \dcat{G}(\cat{M}) \to \dcat{G}(\cat{N})$.
\end{lem}

\begin{proof}
Since for morphisms of degree $0$ there is no sign twist, the identity 
automorphism 
$\opn{id}_M = \{ \opn{id}_{M^i} \}_{i \in \Z}$
of $M = \{ M^i \}_{i \in \Z}$ in $\dcat{G}(\cat{M})$
is sent to the identity automorphism of $\dcat{G}(F)(M)$ in 
$\dcat{G}(\cat{N})$. 

Next we look at morphisms 
$\phi_k = \{ \phi_k^{\lms j} \}_{j \in \Z}$
in 
$\opn{Hom}_{\dcat{G}(\cat{M})}(M_k, M_{k + 1})^{i_k}$ 
for $k = 0, 1$. 
The composition $\phi_1 \circ \phi_0$ has degree $i_0 + i_1$, and the $j$-th 
component of $\phi_1 \circ \phi_0$ is 
$\phi_1^{\lms j + i_0} \circ \phi_0^{\lms j}$.
Therefore the $j$-th component of $\dcat{G}(F)(\phi_1 \circ \phi_0)$ 
is 
\begin{equation} \label{eqn:2510}
\begin{aligned}
& \dcat{G}(F)(\phi_1 \circ \phi_0)^{\lms j} 
= (-1)^{\lms j \cd (i_0 + i_1)} \cd 
F \bigl( \phi_1^{-j - i_1} \circ \phi_0^{-j - (i_0 + i_1)} \bigr)
\\
& \quad = (-1)^{\lms j \cd (i_0 + i_1)} \cd 
F(\phi_0^{-j - (i_0 + i_1)}) \circ F(\phi_1^{-j - i_1}) .
\end{aligned}
\end{equation}
The  $j$-th component of $\dcat{G}(F)(\phi_k)$ is 
$\dcat{G}(F)(\phi_k)^{\lms j} = 
(-1)^{\lms j \cd i_k} \cd F(\phi_k^{-j - i_k})$. 
So the $j$-th component of 
$(-1)^{i_0 \cd i_1} \cd \dcat{G}(F)(\phi_0) \circ \dcat{G}(F)(\phi_1)$ 
is 
\begin{equation} \label{eqn:2511}
\begin{aligned}
& (-1)^{i_0 \cd i_1} \cd 
\bigl( \dcat{G}(F)(\phi_0) \circ \dcat{G}(F)(\phi_1) \bigr)^j 
\\
& \quad = (-1)^{i_0 \cd i_1} \cd 
(-1)^{(j + i_1) \cd i_0} \cd F(\phi_0^{-(j + i_1) - i_0})
\circ (-1)^{\lms j \cd i_1} \cd F(\phi_1^{-j - i_1}) . 
\end{aligned}
\end{equation}
We see that the morphisms (\ref{eqn:2510}) and (\ref{eqn:2511}) are equal. 
\end{proof}

Now we consider a complex $(M, \d_M) \in \dcat{C}(\cat{M})$. 
This is made up of a graded object 
$M  = \{ M^i \}_{i \in \Z} \in \dcat{G}(\cat{M})$ together 
with a differential $\d_M = \{ \d_M^i \}_{i \in \Z}$, 
where $\d_M^i : M^i \to M^{i + 1}$. 
We can view $\d_M$ as an element of 
$\opn{End}_{\dcat{G}(\cat{M})}(M)^1 = 
\opn{Hom}_{\dcat{G}(\cat{M})}(M, M) ^1$.
We specify a differential 
$\d_{\dcat{C}(F)(M)}$ on the graded object 
$\dcat{G}(F)(M) \in \dcat{G}(\cat{N})$
as follows: 
\begin{equation} \label{eqn:2506}
\d_{\dcat{C}(F)(M)} := - \dcat{G}(F)(\d_M) \in 
\opn{End}_{\dcat{G}(\cat{N})} \bigl( \dcat{G}(F)(M) \bigr)^1 . 
\end{equation}
To be explicit, the component 
\[ \d_{\dcat{C}(F)(M)}^i : \dcat{G}(F)(M)^i = F(M^{-i}) \to 
F(M^{-i - 1}) = \dcat{G}(F)(M)^{i + 1} \]
of $\d_{\dcat{C}(F)(M)}$ is, by (\ref{eqn:2505}), 
\begin{equation} \label{eqn:2365}
\d_{\dcat{C}(F)(M)}^i = (-1)^{i + 1} \cd F(\d_M^{-i - 1}) :
F(M^{-i}) \to F(M^{-i - 1}) . 
\end{equation}
This shows that our formula coincides with the one in 
\cite[Remark 1.8.11]{KaSc1}. 

The full subcategory $\dcat{C}^{\star}(\cat{N}) \sub \dcat{C}^{}(\cat{N})$, for 
a given boundedness condition $\star$, was introduced in Definition 
\ref{dfn:3220}. Each boundedness condition $\star$ has a {\em reversed 
boundedness condition} $-\star$, with the obvious meaning; e.g.\ if 
$\star$ is $+$ then $-\star$ is $-$. 

\begin{thm} \label{thm:3015} 
Let $\cat{M}$ and $\cat{N}$ be abelian categories, and let 
$F : \cat{M} \to \cat{N}$ be a contravariant linear functor. 
The assignments \tup{(\ref{eqn:2496})}, \tup{(\ref{eqn:2497})} and 
\tup{(\ref{eqn:2506})} produce a contravariant DG functor 
\[ \dcat{C}(F) : \dcat{C}(\cat{M}) \to \dcat{C}(\cat{N}) . \]
For every boundedness condition $\star$ we have 
$\dcat{C}(F) \bigl( \dcat{C}^{\star}(\cat{M}) \bigr) \sub
\dcat{C}^{- \star}(\cat{N})$,
where $-\star$ is the reversed boundedness condition. 
\end{thm}

\begin{proof}
We already know, by Lemma \ref{lem:2495}, that $\dcat{G}(F)$ is a graded 
functor. We need to prove that for a pair of DG modules 
$(M_0, \d_{M_0})$ and $(M_1, \d_{M_1})$ in $\dcat{C}(\cat{M})$ 
the strict homomorphism of graded $\K$-modules 
\[ \dcat{G}(F) : \opn{Hom}_{\dcat{G}(\cat{M})}(M_0, M_1) \to 
\opn{Hom}_{\dcat{G}(\cat{N})} \bigl( \dcat{G}(F)(M_1), \dcat{G}(F)(M_0) \bigr) 
\]
respects differentials. Take any 
$\phi \in \opn{Hom}_{\dcat{G}(\cat{M})}(M_0, M_1)^i$.
By definition we have 
$\d(\phi) = \d_{M_1} \circ \phi - (-1)^i \cd \phi \circ \d_{M_0}$.
Using the fact that $\dcat{G}(F)$ is a contravariant graded functor, we obtain 
these equalities: 
\[ \begin{aligned}
& \dcat{G}(F)(\d(\phi)) 
\\
& \ = (-1)^{i} \cd \dcat{G}(F)(\phi) \circ \dcat{G}(F)(\d_{M_1}) 
- (-1)^i \cd (-1)^i \cd \dcat{G}(F)(\d_{M_0}) \circ \dcat{G}(F)(\phi)
\\
& \ = \d_{\dcat{C}(F)(M_0)} \circ \dcat{G}(F)(\phi)
- (-1)^{i} \cd \dcat{G}(F)(\phi) \circ \d_{\dcat{C}(F)(M_1)} 
= \d(\dcat{G}(F)(\phi)) . 
\end{aligned} \]

The claim about the boundedness conditions is immediate from equation 
(\ref{eqn:2496}). 
\end{proof}

The sign appearing in formula (\ref{eqn:2506}) might seem arbitrary. 
Besides being the only sign for which Theorem \ref{thm:3015} holds, 
there is another explanation, which can be seen in the next exercise. 

\begin{exer} \label{exer:2505}
Take $\cat{M} = \cat{N} := \cat{Mod} \K$, and consider the contravariant 
additive functor $F := \opn{Hom}_{\K}(-, \K)$ from $\cat{M}$ to itself.
Let $M \in \dcat{C}(\cat{M})$; we can view $M$ as a complex of $\K$-modules or 
as a DG $\K$-module, as done in Proposition \ref{prop:1065}. 
Show that 
$\dcat{C}(F)(M) \cong \opn{Hom}_{\K}(M, \K)$
in $\dcat{C}_{\mrm{str}}(\K)$, where the second object is the graded module 
from formula (\ref{eqn:1700}), with the differential 
$\d$ from Definition \ref{dfn:2992}(2). 
\end{exer}

The next definition and theorem will help us later when studying contravariant 
triangulated functors.

\begin{dfn} \label{dfn:3015}
Let $A$ be a DG ring and let $\cat{M}$ be an abelian 
category. The {\em flipped category}
\index{Flipped category}
of $\dcat{C}(A, \cat{M})$ is the DG category 
$\dcat{C}(A, \cat{M})^{\mrm{flip}} := 
\dcat{C}(A^{\mrm{op}}, \cat{M}^{\mrm{op}})$. 
\end{dfn}

\begin{thm} \label{thm:2495}
Let $A$ be a DG ring and let $\cat{M}$ be an abelian category. Then\tup{:}
\begin{enumerate}
\item There is a canonical isomorphism of DG categories  
\[ \opn{Flip} : 
\dcat{C}(A, \cat{M})^{\mrm{op}} \iso 
\dcat{C}(A^{\mrm{op}}, \cat{M}^{\mrm{op}}) = \dcat{C}(A, \cat{M})^{\mrm{flip}} . 
\]

\item For every boundedness condition $\star$ we have 
\[ \opn{Flip} \bigl( \dcat{C}^{\star}(A, \cat{M})^{\mrm{op}} \bigr) = 
\dcat{C}^{-\star}(A^{\mrm{op}}, \cat{M}^{\mrm{op}}) , \]
where $-\star$ is the reversed boundedness condition. 

\item The induced isomorphism on the strict categories
\[ \opn{Str}(\opn{Flip}) : \dcat{C}_{\mrm{str}}(A, \cat{M})^{\mrm{op}} \iso 
\dcat{C}_{\mrm{str}}(A^{\mrm{op}}, \cat{M}^{\mrm{op}}) \]
is an exact functor, for the respective abelian category structures.

\item For every integer $i$, the functor $\opn{Str}(\opn{Flip})$ makes the 
diagram 
\[ \UseTips  \xymatrix @C=12ex @R=6ex {
\dcat{C}_{\mrm{str}}(A, \cat{M})^{\mrm{op}}
\ar[r]^{\opn{Str}(\opn{Flip})}_{\cong}
\ar[dr]_{\opn{H}^i}
&
\dcat{C}_{\mrm{str}}(A^{\mrm{op}}, \cat{M}^{\mrm{op}})
\ar[d]^{\opn{H}^{-i}}
\\
&
\cat{M}^{\mrm{op}}
} \]
commutative, up to an isomorphism of linear functors. 
\end{enumerate}
\end{thm}

\begin{proof} \mbox{}

\smallskip \noindent 
(1) According to Proposition \ref{prop:2500} there is a contravariant DG 
functor 
$\opn{Op} : \lb \dcat{C}(A, \cat{M})^{\mrm{op}} \to \dcat{C}(A, \cat{M})$.
It is bijective on objects and morphisms. We are going to construct a
contravariant DG functor
$E : \dcat{C}(A, \cat{M}) \to \dcat{C}(A^{\mrm{op}}, \cat{M}^{\mrm{op}})$
which is also bijective on objects and morphisms. The composed DG functor 
\[ \opn{Flip} := E \circ \opn{Op} : \dcat{C}(A, \cat{M})^{\mrm{op}} \to 
\dcat{C}(A^{\mrm{op}}, \cat{M}^{\mrm{op}}) \]
will have the desired properties. 

Let us construct $E$. We start with the contravariant additive functor 
$F := \opn{Op} : \cat{M} \to \cat{M}^{\mrm{op}}$.
Theorem \ref{thm:3015}   says that 
$\dcat{C}(F) : \dcat{C}(\cat{M}) \to \dcat{C}(\cat{M}^{\mrm{op}})$
is a contravariant DG functor. 
Recall that an object of $\dcat{C}(A, \cat{M})$ is a triple 
$(M, \d_M, f_M)$, where $M \in \dcat{G}(\cat{M})$; $\d_M$ is a differential on 
the graded object $M$; and 
$f_M : A \to \opn{End}_{\dcat{C}(\cat{M})}(M)$
is a DG ring homomorphism. See Definitions \ref{dfn:1703}, \ref{dfn:1130} and 
\ref{dfn:1213}.
Define 
$(N, \d_N) := \dcat{C}(F)(M, \d_M) \in \dcat{C}(\cat{M}^{\mrm{op}})$.
Since 
$\dcat{C}(F) : \opn{End}_{\dcat{C}(\cat{M})}(M, \d_M) \to 
\opn{End}_{\dcat{C}(\cat{M}^{\mrm{op}})}(N, \d_N)$
is a DG ring anti-homomorphism (by which we mean the single object version of a 
contravariant DG functor), and 
$\opn{Op} : A^{\mrm{op}} \to A$ 
is also such an anti-homomorphism, it follows that 
\[ f_N := \dcat{C}(F) \circ f_M \circ \opn{Op} : A^{\mrm{op}} \to 
\opn{End}_{\dcat{C}(\cat{M}^{\mrm{op}})}(N, \d_N) \]
is a DG ring homomorphism. Thus 
$E(M, \d_M, f_M) := (N, \d_N, f_N)$ 
is an object of $\dcat{C}(A^{\mrm{op}}, \cat{M}^{\mrm{op}})$.
In this way we have a function 
\[ E : \opn{Ob} \bigl( \dcat{C}(A, \cat{M}) \bigr) \to 
\opn{Ob} \bigl( \dcat{C}(A^{\mrm{op}}, \cat{M}^{\mrm{op}}) \bigr) , \]
and it is clearly bijective.  

The operation of $E$ on morphisms is of course that of $\dcat{C}(F)$. 
It remains to verify that the resulting morphisms in 
$\dcat{C}(\cat{M}^{\mrm{op}})$ respect the action of 
elements of $A^{\mrm{op}}$. Namely that the condition in Definition 
\ref{dfn:1703} is satisfied. Take any morphism 
\[ \phi \in \opn{Hom}_{\dcat{C}(A, \cat{M})} 
\bigl( (M_0, \d_{M_0}, f_{M_0}), (M_1, \d_{M_1}, f_{M_1}) \bigr)^i \]
and any element $a \in (A^{\mrm{op}})^j$; and write 
$(N_k, \d_{N_k}, f_{N_k}) := E(M_k, \d_{M_k}, f_{M_k})$
and
\[ \psi := \dcat{G}(F)(\phi) \in 
\opn{Hom}_{\dcat{C}(\cat{M}^{\mrm{op}})} 
\bigl( (N_1, \d_{N_1}), (N_0, \d_{N_0}) \bigr)^i . \]
We have to prove that 
\begin{equation} \label{eqn:2366}
\psi \circ f_{N_1}(a) = (-1)^{i \cd j} \cd f_{N_0}(a) \circ \psi . 
\end{equation}
This is done using Lemma \ref{lem:2495}, like in the proof of 
Theorem \ref{thm:3015}; and we leave this to the reader.

\medskip \noindent 
(2) Clear from Theorem \ref{thm:3015}.

\medskip \noindent 
(3) Exactness in the categories 
$\dcat{C}_{\mrm{str}}(A, \cat{M})^{\mrm{op}}$
and 
$\dcat{C}_{\mrm{str}}(A^{\mrm{op}}, \cat{M}^{\mrm{op}})$
is checked in each degree separately, and both are exactness in the abelian 
category $\cat{M}^{\mrm{op}}$. The functor
$\opn{Str}(\opn{Flip})$ is $\pm \opn{Id}_{\cat{M}^{\mrm{op}}}$
in each degree, so it is exact. 

\medskip \noindent 
(4) As complexes in $\cat{M}^{\mrm{op}}$, $M$ and $\opn{Flip}(M)$ are equal, up 
to the renumbering of degrees and the signs of the differentials in various 
degrees. So the cohomology objects satisfy 
$\opn{H}^i(M) \cong \opn{H}^{-i}(\opn{Flip}(M))$
in $\cat{M}^{\mrm{op}}$. 
\end{proof}

\begin{exer} \label{exer:2365}
Prove formula (\ref{eqn:2366}) above. 
\end{exer}



\begin{rem} \label{rem:2505}
Theorem \ref{thm:2495} will be used to introduce a triangulated structure on 
the category $\dcat{K}(A, \cat{M})^{\mrm{op}}$. This will be done in Subsection 
\ref{subsec:opp-hom-triang}.  

Combined with Proposition \ref{prop:2500}, Theorem \ref{thm:2495} allows us to 
replace a contravariant DG functor 
$F : \dcat{C}(A, \cat{M}) \to \cat{D}$
with a usual, covariant, DG functor 
$F \circ \opn{Flip}^{-1} : \dcat{C}(A^{\mrm{op}}, \cat{M}^{\mrm{op}})
\to \cat{D}$.
This replacement is going to be very useful when discussing formal properties, 
such as existence of derived functors etc. 

However, in practical terms (e.g.\ for producing resolutions of DG modules),
the category $\dcat{C}(A^{\mrm{op}}, \cat{M}^{\mrm{op}})$ is not very 
helpful. The reason is that the opposite abelian category 
$\cat{M}^{\mrm{op}}$ is almost always a synthetic construction (it does not 
``really exist in concrete terms''). See Remark \ref{rem:2510} and 
Example \ref{exa:2990}. 

We are going to maneuver between the two approaches for reversal of morphisms, 
each time choosing the more useful approach. 
\end{rem}







\cleardoublepage
\mysection{Translations and Standard Triangles} \label{sec:stand-tri}

\AYcopyright

As before, we fix a $\K$-linear abelian category $\cat{M}$,  and a
DG central $\K$-ring $A$. In this section we study the translation functor and 
the standard cone of a strict morphism, all in the context of the DG category 
$\dcat{C}(A, \cat{M})$. 

We then study properties of DG functors 
$F : \dcat{C}(A, \cat{M}) \to \dcat{C}(B, \cat{N})$
between such DG categories. In view of Theorem \ref{thm:2495} it suffices to 
look at covariant DG functors (and not to worry about contravariant DG 
functors). 


\mysubsection{The Translation Functor} \label{subsec:translation}

The translation functor goes back to the beginnings of derived categories -- 
see Remark \ref{rem:1275}. The treatment in this subsection, with the
``little $\opn{t}$ operator'',  is taken from \cite[Section 1]{Ye9}. 
Here and later, the word ``operator'' is used as a 
synonym for ``morphism in a linear category''. 

\begin{dfn} \label{dfn:1260}
Let $M = \{ M^i \}_{i \in \Z}$ be a graded object in $\cat{M}$, i.e.\ an 
object of $\dcat{G}(\cat{M})$. The {\em translation}
\index{Translation! of a graded module}
of $M$ is the object 
\[ \opn{T}(M) = \bigl\{ \opn{T}(M)^i \bigr\}_{i \in \Z} \in \dcat{G}(\cat{M}) \]
defined as follows: the graded component of degree $i$ of $\opn{T}(M)$ is
$\opn{T}(M)^i := M^{i + 1}$.
\end{dfn}

\begin{dfn}[The little t operator] \label{dfn:1261}
Let $M = \{ M^i \}_{i \in \Z}$ be an object of $\dcat{G}(\cat{M})$. We define 
\[ \opn{t}_M : M \to \opn{T}(M) \]
\index{Little t operator}%
\index{1-TM@$\opn{t}_{M}$}
to be the degree $-1$ morphism of graded objects of $\cat{M}$, that for every 
degree $i$ is the identity morphism
${\opn{t}_M}|_{M^{i}} := \opn{id}_{M^{i}} :  M^{i} \iso \opn{T}(M)^{i - 1}$
of the object $M^{i}$ in $\cat{M}$.
\end{dfn}

Note that the morphism 
$\opn{t}_M : M \to \opn{T}(M)$ is a degree $-1$ isomorphism in the graded 
category $\dcat{G}(\cat{M})$. 

\begin{dfn} \label{dfn:2300}
For $M \in \dcat{G}(\cat{M})$ we define the morphism 
$\opn{t}_M^{-1} : \opn{T}(M) \to M$
in $\dcat{G}(\cat{M})$ to be the inverse of $\opn{t}_M$. 
\end{dfn}

Of course the morphism $\opn{t}_M^{-1}$ has degree $+1$. The reason for stating 
this definition is to avoid the potential confusion between the morphism 
$\opn{t}_M^{-1}$ in $\dcat{G}(\cat{M})$ and the degree $-1$ component of the 
morphism $\opn{t}_M$, which we denote by ${\opn{t}_M}|_{M^{-1}}$, as in 
Definition \ref{dfn:1261} above. 

\begin{dfn} \label{dfn:1170}
Let $M = \{ M^i \}_{i \in \Z}$ be a DG $A$-module in $\cat{M}$, i.e.\ an 
object of $\dcat{C}(A, \cat{M})$.
The {\em translation} of $M$ 
\index{Translation! of a DG module}
\index{1-T(M)@$\opn{T}(M)$}
is the object 
\[ \opn{T}(M) \in \dcat{C}(A, \cat{M}) \]
defined as follows. 
\begin{enumerate}
\item As a graded object of $\cat{M}$, it is as specified in 
Definition \ref{dfn:1260}. 

\item The differential $\d_{\opn{T}(M)}$
is defined by the formula
\[ \d_{\opn{T}(M)} := - \opn{t}_M \circ \, \d_M^{} \circ 
\opn{t}_M^{-1} . \]

\item Let $f_M : A \to \opn{End}_{\cat{M}}(M)$ be the DG ring homomorphism that 
determines the action of $A$ on $M$. Then 
\[ f_{\lms \opn{T}(M)} : A \to \opn{End}_{\cat{M}}(\opn{T}(M)) \]
is defined by
\[ f_{\lms \opn{T}(M)}(a) := (-1)^{\lms j} \cd 
\opn{t}_M \circ \, f_{M}(a) \circ \opn{t}_M^{-1} \]
for $a \in A^j$.
\end{enumerate}
\end{dfn}

Thus, the differential 
$\d_{\lms \opn{T}(M)}  = \{ \d_{\lms \opn{T}(M)}^i \}_{i \in \Z}$ makes this 
diagram in $\cat{M}$ commutative for every $i$~: 
\[ \UseTips  \xymatrix @C=8ex @R=6ex {
\opn{T}(M)^i 
\ar[r]^{\d_{\lms \opn{T}(M)}^i}
&
\opn{T}(M)^{i + 1}
\\
M^{i + 1}
\ar[u]^{\opn{t}_M}
\ar[r]^{- \d_M^{i + 1}}
&
M^{i + 2}
\ar[u]_{\opn{t}_M}
} \]
And the left $A$-module structure makes this diagram in $\cat{M}$ 
commutative for every $i$ and every $a \in A^j$~:
\[ \UseTips  \xymatrix @C=12ex @R=6ex {
\opn{T}(M)^i 
\ar[r]^{f_{\lms \opn{T}(M)}(a)}
&
\opn{T}(M)^{i + j}
\\
M^{i + 1}
\ar[u]^{\opn{t}_M}
\ar[r]^{ (-1)^{\lms j} \cd f_{\lms M}(a)}
&
M^{i + j + 1}
\ar[u]_{\opn{t}_M}
} \]

\begin{prop}  \label{prop:1176}
The morphisms $\opn{t}_M$ and $\opn{t}_M^{-1}$ are cocycles, in the DG 
$\K$-modules 
$\opn{Hom}_{A, \cat{M}} \bigl( M, \opn{T}(M) \bigr)$ and 
$\opn{Hom}_{A, \cat{M}} \bigl( \opn{T}(M), M \bigr)$
respectively.
\end{prop}

\begin{proof}
We use the notation $\d_{\mrm{Hom}}$ for the differential in the DG module \lb 
$\opn{Hom}_{A, \cat{M}} \bigl( M, \opn{T}(M) \bigr)$. 
Let us calculate. Because $\opn{t}_M$ has degree $-1$, we have
\[ \begin{aligned}
& \d_{\mrm{Hom}}(\opn{t}_M) = 
\d_{\opn{T}(M)} \circ  \opn{t}_M + \opn{t}_M \circ \, \d_{M}
\\ & \quad 
=  (- \opn{t}_M \circ \, \d_M^{} \circ \opn{t}_M^{-1}) 
\circ  \opn{t}_M + \opn{t}_M \circ \, \d_{M} = 0 . 
\end{aligned} \]

As for $\opn{t}_M^{-1}$~: this is done using the graded Leibniz rule, just like 
in the proof  Proposition \ref{prop:1220}.
\end{proof}

\begin{dfn} \label{dfn:1171}
\index{Translation! of a graded morphism}
Given a  morphism 
$\phi \in \opn{Hom}_{A, \cat{M}} (M, N)^{i}$
we define the morphism 
\[ \opn{T}(\phi)  \in 
\opn{Hom}_{A, \cat{M}} \bigl( \opn{T}(M), \opn{T}(N) \bigr)^i \]
to be
\[ \opn{T}(\phi) := (-1)^i \cd \opn{t}_N \circ \, \phi \circ 
\opn{t}_M^{-1}  .  \]
\end{dfn}

To clarify this definition, let us write 
$\phi = \{ \phi^{\lms j} \}_{j \in \Z}$, 
so that $\phi^{\lms j} : M^j \to N^{j + i}$
is a morphism in $\cat{M}$. Then 
$\opn{T}(\phi)^{\lms j} : \opn{T}(M)^{\lms j} \to \opn{T}(N)^{\lms j + i}$
is  
$\opn{T}(\phi)^{\lms j} = (-1)^i \cd \opn{t}_N \circ \, \phi^{\lms j + 1} 
\circ \opn{t}_M^{-1}$.
The corresponding commutative diagram in $\cat{M}$, for each $i, j$, is: 
\begin{equation} \label{eqn:1162}
 \UseTips  \xymatrix @C=12ex @R=6ex {
\opn{T}(M)^{\lms j}
\ar[r]^{\opn{T}(\phi)^{\lms j}}
&
\opn{T}(N)^{\lms j + i}
\\
M^{j + 1}
\ar[u]^{\opn{t}_M}
\ar[r]^{ (-1)^i \cd \phi^{\lms j + 1} }
&
N^{j + i + 1}
\ar[u]_{\opn{t}_N}
}
\end{equation}

\begin{thm} \label{thm:1260} 
Let $\cat{M}$ be an abelian category and let $A$ be a DG ring.
\begin{enumerate}
\item The assignments $M \mapsto \opn{T}(M)$ and 
$\phi \mapsto \opn{T}(\phi)$ are a  DG functor
\[ \opn{T} : \dcat{C}(A, \cat{M}) \to \dcat{C}(A, \cat{M}) . \]

\item The collection 
$\opn{t} := \{ \opn{t}_M \}_{M \in \dcat{C}(A, \cat{M})}$ 
is a degree $-1$ isomorphism 
\[ \opn{t} : \opn{Id} \to  \opn{T} \]
of DG functors from $\dcat{C}(A, \cat{M})$ to itself. 
\end{enumerate}
\end{thm}

\begin{proof} \mbox{}

\smallskip \noindent
(1) Take morphisms
$\phi_1 : M_0 \to M_1$ and $\phi_2 : M_1 \to M_2$, of degrees $i_1$ and $i_2$ 
respectively. Then 
\[ \begin{aligned}
& \opn{T}(\phi_2 \circ \phi_1) = 
(-1)^{i_1 + i_2} \cd \opn{t}_{M_2} \circ \, (\phi_2 \circ \phi_1) \circ 
\opn{t}_{M_0}^{-1}
\\ 
& \qquad  = (-1)^{i_1 + i_2} \cd \opn{t}_{M_2} \circ \, \phi_2 \circ 
(\opn{t}_{M_1}^{-1} \circ \opn{t}_{M_1}) \circ \, \phi_1 \circ  
\opn{t}_{M_0}^{-1} 
\\ 
& \qquad = \bigl( (-1)^{i_2} \cd \opn{t}_{M_2} \circ \, \phi_2 \circ 
\opn{t}_{M_1}^{-1}\bigr) \circ 
\bigl( (-1)^{i_1} \cd  \opn{t}_{M_1} \circ \, \phi_1 \circ \opn{t}_{M_0}^{-1} 
\bigr)
\\ 
& \qquad =
\opn{T}(\phi_2) \circ \opn{T}(\phi_1) . 
\end{aligned} \]
Clearly $\opn{T}(\opn{id}_{M}) = \opn{id}_{\opn{T}(M)}$, and 
$\opn{T}(\la \cd \phi + \psi) = \la \cd \opn{T}(\phi) + \opn{T}(\psi)$
for all $\la \in \K$ and 
$\phi, \psi \in \opn{Hom}_{A, \cat{M}} (M_0, M_1)^{i}$. 
So $\opn{T}$ is a $\K$-linear graded functor. 

By Proposition \ref{prop:1176} we know that 
$\d \circ \opn{t} = - \opn{t} \circ \, \d$
and 
$\d \circ \opn{t}^{-1} = - \opn{t}^{-1} \circ \, \d$.
This implies that for every morphism 
$\phi$ in $\dcat{C}(A, \cat{M})$, we have 
$\opn{T}(\d(\phi)) = \d(\opn{T}(\phi))$.
So $\opn{T}$ is a DG functor. 

\medskip \noindent
(2) Take some $\phi \in \opn{Hom}_{A, \cat{M}} (M_0, M_1)^{i}$. 
We have to prove that 
$\opn{t}_{M_1} \circ \, \phi = \lb (-1)^{i} \cd \opn{T}(\phi) \circ 
\opn{t}_{M_0}$ 
as elements of 
$\opn{Hom}_{A, \cat{M}} \bigl( M_0, \opn{T}(M_1) \bigr)^{i + 1}$.
But by Definition \ref{dfn:1171} we have 
\[ \opn{T}(\phi) \circ \opn{t}_{M_0} = 
\bigl( (-1)^{i} \cd \opn{t}_{M_1} \circ \, \phi \circ \opn{t}_{M_0}^{-1} \bigr) 
\circ \opn{t}_{M_0} =  (-1)^{i} \cd \opn{t}_{M_1} \circ \, \phi . \qedhere \]
\end{proof}

\begin{dfn} \label{dfn:1275}
\index{Translation functor! of the DG category $\dcat{C}(A, \cat{M})$}
We call $\opn{T}$ the {\em translation functor} of the DG category 
$\dcat{C}(A, \cat{M})$.  
\end{dfn}

\begin{cor} \label{cor:1300} \mbox{}
\begin{enumerate}
\item The functor $\opn{T}$ is an automorphism of the DG category 
$\dcat{C}(A, \cat{M})$.

\item For every $k, l \in \Z$ there is an equality of functors
$\opn{T}^{l} \circ \opn{T}^{k} = \opn{T}^{l + k}$.
\end{enumerate}
\end{cor}

\begin{proof}
(1) By the theorem, $\opn{T}$ is a DG functor. By definition, $\opn{T}$
is bijective on the set of objects of $\dcat{C}(A, \cat{M})$ and on the sets of 
morphisms. 

\medskip \noindent 
(2) By part (1) of this corollary, the inverse $\opn{T}^{-1}$ is a uniquely 
defined functor (not just up to an isomorphism of functors). 
Thus we can define $\opn{T}^k$, the $k$-th power of $\opn{T}$, and the equality 
stated holds. 
\end{proof}

\begin{prop} \label{prop:1935}
Consider any $M \in \dcat{C}(A, \cat{M})$.
\begin{enumerate}
\item There is equality 
$\opn{t}_{\opn{T}(M)} = - \opn{T}(\opn{t}_M)$
of degree $-1$ morphisms $\opn{T}(M) \to \opn{T}^2(M)$
in  $\dcat{C}(A, \cat{M})$.

\item There is equality 
$\opn{t}_{\opn{T}^{-1}(M)} = - \opn{T}^{-1}(\opn{t}_M)$
between these degree $-1$ morphisms  
$\opn{T}^{-1}(M) \to \opn{T}(\opn{T}^{-1}(M)) = M = \opn{T}^{-1}(\opn{T}(M))$
in  $\dcat{C}(A, \cat{M})$.
\end{enumerate}
\end{prop}

\begin{proof}
(1) This is an easy calculation, using Definition \ref{dfn:1171}:
$\opn{T}(\opn{t}_M) = \lb 
- \opn{t}_{\opn{T}(M)} \circ \, \opn{t}_{M} \circ \opn{t}_{M}^{-1} = 
- \opn{t}_{\opn{T}(M)}$.

\medskip \noindent 
(2) A similar calculation. 
\end{proof}

\begin{prop} \label{prop:4155}
Let $M \in \dcat{C}(A, \cat{M})$ and $i \in \Z$. There is equality 
$\d_{\opn{T}^i(M)} = \opn{T}^i(\d_M)$
in 
$\opn{Hom}_{A, \cat{M}} \bigl(\opn{T}^i(M) , \opn{T}^i(M) \bigr)^{1}$. 
\end{prop}

\begin{proof}
We start with $i = 1$. The differential $\d_M$ is an element of 
$\opn{Hom}_{\cat{M}}(M , M)^1$. 
By Definitions \ref{dfn:1170} and \ref{dfn:1171} we get 
$\d_{\opn{T}(M)} = - \opn{t}_M \circ \, \d_M \circ \opn{t}_M^{-1} 
=  \opn{T}(\d_{M})$.
Using induction on $i$ the assertion holds for all $i \geq 0$. 

For $i \leq 0$ we use descending induction on $i$. 
We assume that the assertion holds for $i$. 
Let us define $N := \opn{T}^{i - 1}(M)$, so 
$\opn{T}^{}(N) = \opn{T}^{i}(M)$. 
By the previous paragraph, with the DG module $N$, we know that 
$\d_{\opn{T}^{i}(M)} = \d_{\opn{T}^{}(N)} = \opn{T}(\d_{N}) = 
\opn{T}(\d_{\opn{T}^{i - 1}(M)})$.
Applying the functor $\opn{T}^{-1}$ to this equality we get 
$\opn{T}^{i - 1}(\d_{M}) = 
\opn{T}^{-1} \bigl( \opn{T}^{i}(\d_M) \bigr) = 
\opn{T}^{-1}(\d_{\opn{T}^{i}(M)})
= \d_{\opn{T}^{i - 1}(M)}$. % \qedhere \]
\end{proof}

\begin{rem} \label{rem:1275}
There are several names in the literature for the translation functor
$\opn{T}$~: {\em twist}, {\em shift} and {\em suspension}. There are also 
several notations: $\opn{T}(M) = M[1] = \Sigma M$.
In the later part of this book we shall use the notation 
$M[k] := \opn{T}^{k}(M)$ for the $k$-th translation. 
\end{rem}


\mysubsection{The Standard Triangle of a Strict Morphism} \label{subsec:cone}

As before,  we fix an abelian category $\cat{M}$ and a DG ring $A$.
Here is the cone construction in $\dcat{C}(A, \cat{M})$, as it looks 
using the operator $\opn{t}$. 

\begin{dfn} \label{dfn:1172}
Let $\phi : M \to N$ be a strict morphism in $\dcat{C}(A, \cat{M})$.
The {\em standard cone of $\phi$}%
\index{Cone! standard {\indash} of strict morphism in $\dcat{C}(A, \cat{M})$}
is the object 
$\opn{Cone}(\phi) \in \dcat{C}(A, \cat{M})$
defined as follows. As a graded $A$-module in $\cat{M}$ we let
\[ \opn{Cone}(\phi) := N \oplus \opn{T}(M) . \]
The differential $\d_{\mrm{Cone}}$ is this: if we express 
the graded module as a column
\[ \opn{Cone}(\phi) = \bmat{N \\[0.1em] \opn{T}(M)} , \]
then $\d_{\mrm{Cone}}$ is left multiplication by the matrix
\[ \d_{\mrm{Cone}} :=
\bmat{\d_N & \phi \circ \opn{t}_M^{-1} \\[0.1em] 0 & \d_{\opn{T}(M)}} \]
of degree $1$ morphisms of graded $A$-modules in $\cat{M}$.
\end{dfn}

In other words, the morphism 
$\d_{\mrm{Cone}}^i : \opn{Cone}(\phi)^i \to \opn{Cone}(\phi)^{i + 1}$
is 
$\d_{\mrm{Cone}}^i = \d^i_N + \d^i_{\opn{T}(M)} + \phi^{i + 1} \circ 
\opn{t}_M^{-1}$,
where $\phi^{i + 1} \circ \opn{t}_M^{-1}$ is the composed morphism 
$\opn{T}(M)^i \xar{\opn{t}_M^{-1}} M^{i + 1} \xar{\phi^{i + i}}  N^{i + 1}$.

In the situation of Definition \ref{dfn:1172}, let us denote by 
\begin{equation} \label{eqn:1195}
e_{\phi} : N \to N \oplus \opn{T}(M) 
\end{equation}
the embedding, and by
\begin{equation} \label{eqn:1196}
p_{\phi} : N \oplus \opn{T}(M) \to \opn{T}(M)
\end{equation}
the projection. 
Thus, as matrices we have 
\[ e_{\phi} = \bmat{ \opn{id}_N  \\[0.1em] 0 } \quad \tup{and} \quad 
p_{\phi} = \bmat{ 0 &  \opn{id}_{\opn{T}(M)} } . \]
The standard cone of $\phi$ sits in the exact sequence 
\begin{equation} \label{eqn:1815}
0 \to N \xar{e_{\phi}} \opn{Cone}(\phi) \xar{p_{\phi}} \opn{T}(M) \to 0
\end{equation}
in the abelian category $\dcat{C}_{\mrm{str}}(A, \cat{M})$. 

\begin{dfn} \label{dfn:1140}
Let $\phi : M \to N$ be a morphism in $\dcat{C}_{\mrm{str}}(A, \cat{M})$.
The diagram 
\[ M \xar{ \, \phi \, } N \xar{ \, e_{\phi} \, } \opn{Cone}(\phi) 
\xar{ \, p_{\phi} \, } \opn{T}(M) \]
in $\dcat{C}_{\mrm{str}}(A, \cat{M})$ is called the {\em standard triangle} 
\index{Triangle! standard {\indash} of a strict morphism in 
$\dcat{C}(A, \cat{M})$}
associated to $\phi$. 
\end{dfn}

The cone construction is functorial, in the following sense. 

\begin{prop}\label{prop:1162}
Let 
\[ \UseTips  \xymatrix @C=8ex @R=6ex {
M_0 
\ar[r]^{\phi_0}
\ar[d]_{\psi}
& 
N_0 
\ar[d]^{\chi}
\\
M_1
\ar[r]^{\phi_1}
& 
N_1 
} \]
be a commutative diagram in $\dcat{C}_{\mrm{str}}(A, \cat{M})$. 
Then 
\begin{equation} \label{eqn:1113}
\bigl( \chi, \opn{T}(\psi) \bigr) : \opn{Cone}(\phi_0) \to \opn{Cone}(\phi_1)
\end{equation}
is a morphism in $\dcat{C}_{\mrm{str}}(A, \cat{M})$, and the diagram 
\[ \UseTips  \xymatrix @C=10ex @R=6ex {
M_0 
\ar[r]^{\phi_0}
\ar[d]_{\psi}
& 
N_0 
\ar[r]^(0.4){e_{\phi_0}}
\ar[d]_{\chi}
& 
\opn{Cone}(\phi_0)
\ar[r]^{p_{\phi_0}}
\ar[d]_{(\chi, \opn{T}(\psi))}
& 
\opn{T}(M_0)
\ar[d]_{\opn{T}(\psi)}
\\
M_1
\ar[r]^{\phi_1}
& 
N_1
\ar[r]^(0.4){e_{\phi_1}}
& 
\opn{Cone}(\phi_1)
\ar[r]^{p_{\phi_1}}
& 
\opn{T}(M_1)
} \]
in $\dcat{C}_{\mrm{str}}(A, \cat{M})$ is commutative. 
\end{prop}

\begin{proof}
This is a simple consequence of the definitions. 
\end{proof}



\mysubsection{The Gauge of a Graded Functor} \label{subsec:gauge}

Recall that we have an abelian category $\cat{M}$ and a DG ring $A$.
The next definition is new. 

\begin{dfn} \label{dfn:1214}
Let 
$F : \dcat{C}(A, \cat{M}) \to \dcat{C}(B, \cat{N})$
be a graded functor. For every object 
$M \in \dcat{C}(A, \cat{M})$ let 
\[ \ga_{F, M} := \d_{F(M)} - F(\d_M)  \in 
\opn{Hom}_{B, \cat{N}} \bigl( F(M), F(M) \bigr)^1 . \]
The collection of morphisms 
\[ \ga_F := \{ \ga_{F, M} \}_{M \in \dcat{C}(A, \cat{M})}  \]
is called the {\em gauge of $F$}.
\index{Gauge of a graded functor} 
\end{dfn}

The next theorem is due to R. Vyas. 

\begin{thm} \label{thm:1225}
The following two conditions are equivalent for a graded functor 
$F : \dcat{C}(A, \cat{M}) \to \dcat{C}(B, \cat{N})$.
\begin{enumerate}
\rmitem{i} $F$ is a DG functor.
 
\rmitem{ii} The gauge $\ga_F$ is a degree $1$ morphism of graded functors 
$\ga_F : F \to F$.  
\end{enumerate}
\end{thm}

\begin{proof}
Recall that $F$ is a DG functor (condition (i)) iff 
\begin{equation}  \label{eqn:1231}
(F \circ \d_{A, \cat{M}})(\phi) = (\d_{B, \cat{N}} \circ F)(\phi)
\end{equation}
for every 
$\phi \in \opn{Hom}_{A, \cat{M}}(M_{0}, M_{1})^{i}$.
And $\ga_F$ is a degree $1$ morphism of graded functors  (condition (ii))
iff 
\begin{equation}  \label{eqn:1232}
\ga_{F, M_1} \circ  F(\phi) = (-1)^i \cd F(\phi) \circ \ga_{F, M_0} 
\end{equation}
for every such $\phi$. 

Here is the calculation. 
Because $F$ is a graded functor, we get
\begin{equation}  \label{eqn:1225}
\begin{aligned}
& F \bigl( \d_{A, \cat{M}}(\phi) \bigr) = 
F \bigl( \d_{M_1} \circ \phi - (-1)^i \cd \phi \circ \d_{M_0} \bigr) 
\\
& \qquad = F(\d_{M_1}) \circ F(\phi) - (-1)^i \cd F(\phi) \circ F(\d_{M_0})
\end{aligned} 
\end{equation}
and 
\begin{equation}  \label{eqn:1226}
\d_{B, \cat{N}} \bigl( F(\phi)  \bigr) = 
\d_{F(M_1)} \circ F(\phi) - (-1)^i \cd F(\phi) \circ \d_{F(M_0)} .
\end{equation}
Using equations  (\ref{eqn:1225}) and (\ref{eqn:1226}),
and the definition of $\ga_F$, we obtain 
\begin{equation}  \label{eqn:1227}
\begin{aligned}
& (F \circ \d_{A, \cat{M}} - \d_{B, \cat{N}} \circ F)(\phi) = 
F \bigl( \d_{A, \cat{M}}(\phi) \bigr) - \d_{B, \cat{N}} \bigl( F(\phi)  \bigr)
\\
& \qquad = \bigl( F(\d_{M_1}) - \d_{F(M_1)} \bigr) \circ F(\phi) 
- (-1)^i \cd F(\phi) \circ \bigl( F(\d_{M_0}) - \d_{F(M_0)} \bigr)
\\
& \qquad = - \ga_{F, M_1} \circ  F(\phi) 
+ (-1)^i \cd F(\phi) \circ \ga_{F, M_0} . 
\end{aligned} 
\end{equation}

Finally, the vanishing of the first expression in (\ref{eqn:1227}) 
is the same as equality in (\ref{eqn:1231});  whereas the vanishing of the last 
expression in (\ref{eqn:1227}) is the same as equality in (\ref{eqn:1232}).
\end{proof}




\mysubsection{The Translation Isomorphism of a DG Functor}

Here we consider \lb belian categories $\cat{M}$ and $\cat{N}$, 
and DG rings $A$ and $B$.
The translation functor of the DG category $\dcat{C}(A, \cat{M})$ will be 
denoted here by $\opn{T}_{A, \cat{M}}$. 
For an object $M \in \dcat{C}(A, \cat{M})$, we have the little t operator 
\[ \opn{t}_M \in \opn{Hom}_{A, \cat{M}} \bigl( M, \opn{T}_{A, \cat{M}}(M) 
\bigr)^{-1} . \]
This is an isomorphism in the DG category $\dcat{C}(A, \cat{M})$. 
Likewise for the DG category $\dcat{C}(B, \cat{N})$.

\begin{dfn} \label{dfn:1150}
Let 
$F : \dcat{C}(A, \cat{M}) \to \dcat{C}(B, \cat{N})$
be a DG functor.
For an object $M \in \dcat{C}(A, \cat{M})$, let 
\[ \tau_{F, M} : F(\opn{T}_{A, \cat{M}}(M)) \to \opn{T}_{B, \cat{N}}(F(M)) \]
be the degree $0$ isomorphism
\[ \tau_{F, M} := \opn{t}_{F(M)} \circ \, F(\opn{t}_M)^{-1}  \]
in $\dcat{C}(B, \cat{N})$, called the {\em translation isomorphism}
\index{Translation isomorphism! of a DG functor}
of the functor $F$ at the object $M$. 
\end{dfn}

The isomorphism $\tau_{F, M}$ sits in the following commutative diagram
\[ \UseTips  \xymatrix @C=8ex @R=6ex {
F(\opn{T}_{A, \cat{M}}(M))
\ar[r]^{\tau_{F, M}}
&
\opn{T}_{B, \cat{N}}(F(M))
\\
F(M)
\ar[u]^{ F(\opn{t}_M) } 
\ar[ur]_{ \opn{t}_{F(M)} }
} \]
of isomorphisms in the category $\dcat{C}(B, \cat{N})$.

\begin{prop} \label{prop:1179}
$\tau_{F, M}$ is an isomorphism in $\dcat{C}_{\mrm{str}}(B, \cat{N})$.
\end{prop}

\begin{proof}
We know that $\tau_{F, M}$ is an isomorphism in $\dcat{C}(B, \cat{N})$.
It suffices to prove that both $\tau_{F, M}$ and its inverse  
$\tau_{F, M}^{-1}$ are 
strict morphisms. Now by Proposition \ref{prop:1176}, $\opn{t}_M$ and 
$\opn{t}_M^{-1}$ are cocycles. Therefore, $F(\opn{t}_M)$ and 
$F(\opn{t}_M)^{-1} = F(\opn{t}_M^{-1})$ are cocycles. For the same reason,
$\opn{t}_{F(M)}$ and $\opn{t}_{F(M)}^{-1}$  are cocycles. But 
$\tau_{F, M} = \opn{t}_{F(M)} \circ \, F(\opn{t}_M)^{-1}$, and 
$\tau_{F, M}^{-1} = F(\opn{t}_M) \, \circ \opn{t}_{F(M)}^{-1}$.
\end{proof}

\begin{thm} \label{thm:1150}
Let 
$F : \dcat{C}(A, \cat{M}) \to \dcat{C}(B, \cat{N})$
be a DG functor. Then the collection
\[ \tau_F := \{ \tau_{F, M} \}_{M \in \dcat{C}(A, \cat{M})} \] 
is an isomorphism 
\[ \tau_F : F \circ \opn{T}_{A, \cat{M}} \iso \opn{T}_{B, \cat{N}} \circ \, F \]
of functors 
$\dcat{C}_{\mrm{str}}(A, \cat{M}) \to \dcat{C}_{\mrm{str}}(B, \cat{N})$.
\end{thm}

The slogan summarizing this theorem is ``A DG functor commutes with 
translations''. 

\begin{proof}
In view of Proposition \ref{prop:1179}, all we need to prove is that 
$\tau_{F}$ is a morphism of functors (i.e.\ it is a natural transformation). 

Let $\phi : M_0 \to M_1$ be a morphism in 
$\dcat{C}_{\mrm{str}}(A, \cat{M})$. We must prove that the diagram 
\[ \UseTips  \xymatrix @C=8ex @R=6ex {
(F \circ \opn{T}_{A, \cat{M}})(M_0)
\ar[r]^{ \tau_{F, M_0} } 
\ar[d]_{ (F \circ \opn{T}_{A, \cat{M}})(\phi) }
&
(\opn{T}_{B, \cat{N}} \circ \, F)(M_0)
\ar[d]^{ (\opn{T}_{B, \cat{N}} \circ \, F)(\phi) }
\\
(F \circ \opn{T}_{A, \cat{M}})(M_1)
\ar[r]^{ \tau_{F, M_1} } 
&
(\opn{T}_{B, \cat{N}} \circ \, F)(M_1)
} \]
in $\dcat{C}_{\mrm{str}}(B, \cat{N})$ is commutative.
This will be true if the next diagram 
\[ \UseTips  \xymatrix @C=10ex @R=6ex {
(F \circ \opn{T}_{A, \cat{M}})(M_0)
\ar[d]_{ (F \circ \opn{T}_{A, \cat{M}})(\phi) }
&
F(M_0)
\ar[l]_(0.4){ F(\opn{t}_{M_0}) }
\ar[r]^(0.4){ \opn{t}_{F(M_0)} }
\ar[d]_{F(\phi)}
&
(\opn{T}_{B, \cat{N}} \circ \, F)(M_0)
\ar[d]^{ (\opn{T}_{B, \cat{N}} \circ \, F)(\phi) }
\\
(F \circ \opn{T}_{A, \cat{M}})(M_1)
&
F(M_1)
\ar[l]_(0.4){ F(\opn{t}_{M_1}) }
\ar[r]^(0.4){ \opn{t}_{F(M_1)} }
&
(\opn{T}_{B, \cat{N}} \circ \, F)(M_1)
} \]
in $\dcat{C}(B, \cat{N})$, whose horizontal arrows are isomorphisms, is 
commutative. For this to be true, it is enough to prove that both squares in 
this diagram are commutative. This is true by 
Theorem \ref{thm:1260}(2). 
\end{proof}

Recall that the translation $\opn{T}$ and all its powers are DG functors. 
To finish this subsection, we calculate their translation isomorphisms. 

\begin{prop} \label{prop:1300}
For every integer $k$, the translation isomorphism of the DG functor 
$\opn{T}^k$ is 
$\tau_{\opn{T}^k} = (-1)^k \cd \opn{id}_{\lms \opn{T}^{k + 1}}$,
where $\opn{id}_{\lms \opn{T}^{k + 1}}$ is the identity automorphism of the 
functor $\opn{T}^{k + 1}$. 
\end{prop}

\begin{proof}
By Definition \ref{dfn:1150} and Proposition \ref{prop:1935}(1), for $k = 1$ 
the formula is
$\tau_{\lms \opn{T}, M} = \opn{t}_{\opn{T}(M)} \circ \opn{T}(\opn{t}_M)^{-1} = 
- \opn{id}_{\lms \opn{T}^2(M)}$,
where $\opn{id}_{\lms \opn{T}^2(M)}$ is the identity automorphism of the DG 
module 
$\opn{T}^2(M)$. Hence $\tau_{\lms \opn{T}} =  - \opn{id}_{\opn{T}^2}$.
For other integers $k$ the calculation is similar.
\end{proof}



\mysubsection{Standard Triangles and DG Functors} \label{subsec:std-triangles}

In Subsection \ref{subsec:cone} we defined the standard triangle associated to 
a strict morphism; and in Subsection \ref{subsec:DGFunc} we defined DG 
functors. Now we show how these notions interact with each other. 
As before, we consider abelian categories $\cat{M}$ and $\cat{N}$, 
and DG rings $A$ and $B$.

Let 
$F : \dcat{C}(A, \cat{M}) \to \dcat{C}(B, \cat{N})$
be a DG functor. Given a morphism 
$\phi : M_0 \to M_1$ in $\dcat{C}_{\mrm{str}}(A, \cat{M})$,
we have a morphism 
$F(\phi) : F(M_0) \to F(M_1)$
in $\dcat{C}_{\mrm{str}}(B, \cat{N})$, 
and objects 
$F(\opn{Cone}_{A, \cat{M}}(\phi))$ and 
$\opn{Cone}_{B, \cat{N}}(F(\phi))$ 
in $\dcat{C}(B, \cat{N})$.
By definition, and using the fully faithful graded functor
$\opn{Und}$ from Proposition \ref{prop:2365},
there is a canonical isomorphism  
\begin{equation}  \label{eqn:1197}
\opn{Cone}_{A, \cat{M}}(\phi) \cong M_1 \oplus \opn{T}_{A, \cat{M}}(M_0)
\end{equation}
in $\dcat{G}_{\mrm{str}}(A, \cat{M})$.
Since $F$ is a linear functor, it commutes with finite direct sums,
and therefore there is a canonical isomorphism 
\begin{equation}  \label{eqn:1185}
F(\opn{Cone}_{A, \cat{M}}(\phi)) \cong F(M_1) \oplus 
F(\opn{T}_{A, \cat{M}}(M_0))
\end{equation}
in $\dcat{G}_{\mrm{str}}(A, \cat{M})$. 
And by definition there is a canonical isomorphism 
\begin{equation}  \label{eqn:1186}
\opn{Cone}_{B, \cat{N}}(F(\phi)) \cong F(M_1) \oplus 
\opn{T}_{B, \cat{N}}(F(M_0))
\end{equation}
in $\dcat{G}_{\mrm{str}}(A, \cat{M})$.
Warning: the isomorphisms (\ref{eqn:1197}), (\ref{eqn:1185}) and 
(\ref{eqn:1186}) might not commute with the differentials. 
The differentials on the right sides are  diagonal matrices, but on the left 
sides they are upper-triangular matrices (see Definition \ref{dfn:1172}). 

\begin{lem} \label{lem:1200}
Let 
$F , G : \dcat{G}(A, \cat{M}) \to \dcat{G}(B, \cat{N})$
be graded functors, and let $\eta : F \to G$ be a degree $j$ 
morphism of graded functors. Suppose 
$M \cong M_0 \oplus M_1$ in $\dcat{G}_{\mrm{str}}(A, \cat{M})$, with embeddings 
$e_i : M_i \to M$ and  projections $p_i : M \to M_i$. 
Then 
\[ \eta_M = \bigl( G(e_0), G(e_1) \bigr) \circ (\eta_{M_0}, \eta_{M_1}) \circ 
\bigl( F(p_0), F(p_1) \bigr) , \]
as degree $j$ morphisms $F(M) \to G(M)$ in 
$\dcat{G}(B, \cat{N})$. 
\end{lem}

The lemma says that the diagram 
\[ \UseTips  \xymatrix @C=14ex @R=6ex {
F(M) 
\ar[r]^(0.4){ (F(p_0), F(p_1)) }
\ar[d]_{ \eta_M }
& 
F(M_0) \oplus F(M_1)
\ar[d]^{ (\eta_{M_0}, \eta_{M_1}) }
\\
G(M)
& 
G(M_0) \oplus G(M_1)
\ar[l]_(0.6){ (G(e_0), G(e_1)) }
} \]
in $\dcat{G}(B, \cat{N})$ is commutative. 

\begin{proof}
It suffices to prove that the diagram below is commutative for $i = 0, 1$~:
\[ \UseTips  \xymatrix @C=10ex @R=6ex {
F(M_i)
\ar[r]^{ F(e_i) }
\ar[d]_{ \eta_{M_i} }
\ar@(u,u)[rr]^{ \opn{id} }
&
F(M) 
\ar[r]^{ F(p_i) }
\ar[d]_{ \eta_M }
& 
F(M_i) 
\ar[d]_{ \eta_{M_i} }
\\
G(M_i)
\ar[r]^{ G(e_i) }
\ar@(d,d)[rr]_{ \opn{id} }
&
G(M)
\ar[r]^{ G(p_i) }
& 
G(M_i) 
} \]
This is true because $\eta$ is a morphism of functors (a natural 
transformation). 
\end{proof}

\begin{thm} \label{thm:1185}
Let 
$F : \dcat{C}(A, \cat{M}) \to \dcat{C}(B, \cat{N})$
be a DG functor, and let 
$\phi : M_0 \to M_1$ be a morphism in $\dcat{C}_{\mrm{str}}(A, \cat{M})$.
Define the isomorphism 
\[ \opn{cone}(F, \phi) : F (\opn{Cone}_{A, \cat{M}}(\phi)) \iso 
\opn{Cone}_{B, \cat{N}} (F(\phi))  \]
in  $\dcat{G}_{\mrm{str}}(A, \cat{M})$ to be 
\[ \opn{cone}(F, \phi) := 
\bigl( \opn{id}_{F(M_1)}, \tau_{F, M_0} \bigr) . \]
Then\tup{:}
\begin{enumerate}
\item The isomorphism $\opn{cone}(F, \phi)$ commutes with the differentials, 
so it is an isomorphism in $\dcat{C}_{\mrm{str}}(A, \cat{M})$. 

\item The diagram 
\[ \UseTips  \xymatrix @C=7ex @R=6ex {
F(M_0) 
\ar[r]^{F(\phi)}
\ar[d]_{\opn{id}}
& 
F(M_1) 
\ar[r]^(0.4){F(e_{\phi})}
\ar[d]_{\opn{id}}
& 
F(\opn{Cone}_{A, \cat{M}}(\phi))
\ar[r]^{F(p_{\phi})}
\ar[d]_{\opn{cone}(F, \phi)}
& 
F(\opn{T}_{A, \cat{M}}(M_0))
\ar[d]_{ \tau_{F, M_0}}
\\
F(M_0)
\ar[r]^{F(\phi)}
& 
F(M_1)
\ar[r]^(0.4){e_{F(\phi)}}
& 
\opn{Cone}_{B, \cat{N}}(F(\phi))
\ar[r]^{p_{F(\phi)}}
& 
\opn{T}_{B, \cat{N}}(F(M_0))
} \]
in $\dcat{C}_{\mrm{str}}(B, \cat{N})$ is commutative. 
\end{enumerate}
\end{thm}

When defining $\opn{cone}(F, \phi)$ above, we are using the decompositions 
(\ref{eqn:1185}) and  (\ref{eqn:1186}) in the category 
$\dcat{G}_{\mrm{str}}(A, \cat{M})$, 
and the isomorphism $\tau_{F, M_0}$ from Definition 
\ref{dfn:1150}.

The slogan summarizing this theorem is ``A DG functor sends standard triangles 
to standard triangles''. 

\begin{proof} \mbox{}

\smallskip \noindent
(1) To save space let us write $\th := \opn{cone}(F, \phi)$. 
We have to prove that $\d_{B, \cat{N}}(\th) = 0$. 
Let's write 
$P := \opn{Cone}_{A, \cat{M}}(\phi)$
and 
$Q := \opn{Cone}_{B, \cat{N}}(F(\phi))$.
Recall that 
$\d_{B, \cat{N}}(\th) =  \d_{Q} \circ \th - \th \circ \d_{F(P)}$.
We have to prove that this is the zero element in 
$\opn{Hom}_{B, \cat{N}} \bigl( F(P), Q \bigr)^1$.
 
Writing the cones as column modules:
\[ P = \bmat{ M_1 \\[0.2em] \opn{T}_{A, \cat{M}}(M_0) }  
\quad \tup{and}  \quad 
Q = \bmat{ F(M_1) \\[0.2em] \opn{T}_{B, \cat{N}}(F(M_0)) } \, , \]
the matrices representing the morphisms in question are
\[ \th = 
\bmat{\opn{id}_{F(M_1)} & 0 \\[0.2em] 0 & \tau_{F, M_0} } 
\, , \quad 
\d_P = \bmat{\d_{M_1} & \phi \circ \opn{t}_{M_0}^{-1} 
\\[0.2em] 0 & \d_{\opn{T}_{A, \cat{M}}(M_0)} } \]
and \[  
\d_Q = \bmat{\d_{F(M_1)} & F(\phi) \circ \opn{t}_{F(M_0)}^{-1} 
\\[0.2em] 0 & \d_{\opn{T}_{B, \cat{N}}(F(M_0))} } \ . \] 

Let us write $\ga := \ga_F$ for simplicity. 
According to Theorem \ref{thm:1225}, the gauge 
$\ga : F \to F$ is a degree $1$ morphism of DG functors
$\dcat{C}(A, \cat{M}) \to \dcat{C}(B, \cat{N})$.
Because the decomposition (\ref{eqn:1197}) is in the 
category $\dcat{G}_{\mrm{str}}(A, \cat{M})$, Lemma \ref{lem:1200} tells us that 
$\ga_P$ decomposes too, i.e.\ 
\[ \ga_P =  \bmat{ \ga_{M_1} & 0 \\[0.2em] 0 & \ga_{\opn{T}_{A, \cat{M}}(M_0)} 
} . \]

By definition of $\ga_P$ we have 
$\d_{F(P)} = F(\d_P) + \ga_P$ in 
$\opn{Hom}_{B, \cat{N}} \bigl( F(P), F(P) \bigr)^1$.
It follows that 
\[ \begin{aligned}
& \d_{F(P)} = F(\d_P) + \ga_P  
\\[0.2em]
& \qquad 
= \bmat{ F(\d_{M_1})  & 
F(\phi \circ \opn{t}_{M_0}^{-1}) 
\\[0.2em] 0 & F(\d_{\opn{T}_{A, \cat{M}}(M_0)})  } +
\bmat{ \ga_{M_1} & 0 \\[0.2em] 0 & \ga_{\opn{T}_{A, \cat{M}}(M_0)} } 
\\[0.2em]
& \qquad 
= \bmat{ F(\d_{M_1}) +  \ga_{M_1} & 
F(\phi \circ \opn{t}_{M_0}^{-1}) 
\\[0.2em] 0 & 
F(\d_{\opn{T}_{A, \cat{M}}(M_0)}) + \ga_{\opn{T}_{A, \cat{M}}(M_0)} }
\\[0.2em]
& \qquad 
= \bmat{ \d_{F(M_1)} & 
F(\phi \circ \opn{t}_{M_0}^{-1}) 
\\[0.2em] 0 & \d_{F(\opn{T}_{A, \cat{M}}(M_0))} } .
\end{aligned}   \]

Finally we will check that 
$\th \circ \d_{F(P)}$ and $ \d_{Q} \circ \th$ are equal as matrices of 
morphisms. We do that in each matrix entry separately. The two left 
entries in the matrices $\th \circ \d_{F(P)}$ and $ \d_{Q} \circ \th$
agree trivially. The bottom right entries in these matrices
are $\tau_{F, M_0} \circ \d_{F(\opn{T}_{A, \cat{M}}(M_0))}$
and 
$\d_{\opn{T}_{B, \cat{N}}(F(M_0))} \circ \tau_{F, M_0}$
respectively; they are equal by Proposition \ref{prop:1179}. 
And in the top right  entries we have 
$F(\phi \, \circ \opn{t}_{M_0}^{-1})$
and 
$F(\phi) \, \circ \opn{t}_{F(M_0)}^{-1} \circ \, \tau_{F, M_0}$
respectively. Now
$F(\phi \, \circ \opn{t}_{M_0}^{-1}) = F(\phi) \circ F(\opn{t}_{M_0}^{-1})$;
so it suffices to prove that 
$F(\opn{t}_{M_0}^{-1}) = \opn{t}_{F(M_0)}^{-1} \circ \, \tau_{F, M_0}$.
This is immediate from the definition of $\tau_{F, M_0}$.

\medskip \noindent 
(2) By definition of $\th = \opn{cone}(F, \phi)$, the diagram is commutative in 
$\dcat{G}_{\mrm{str}}(A, \cat{M})$. But by part (1) we know that all morphisms 
in it commute with the differentials, so they lie in 
$\dcat{C}_{\mrm{str}}(B, \cat{N})$.
And the functor $\opn{Und}$ from Proposition \ref{prop:2365}
is faithful. 
\end{proof}

\begin{cor} \label{cor:1915}
In the situation of Theorem \tup{\ref{thm:1185}}, the diagram 
\[ \UseTips  \xymatrix @C=6.2ex @R=6ex {
F(M_0) 
\ar[r]^{F(\phi)}
\ar[d]_{\opn{id}}
& 
F(M_1) 
\ar[r]^(0.4){F(e_{\phi})}
\ar[d]_{\opn{id}}
& 
F(\opn{Cone}_{A, \cat{M}}(\phi))
\ar[rr]^{\tau_{F, M_0} \, \circ \, F(p_{\phi})}
\ar[d]_{\opn{cone}(F, \phi)}
& &
\opn{T}_{B, \cat{N}}(F(M_0))
\ar[d]_{\opn{id}}
\\
F(M_0)
\ar[r]^{F(\phi)}
& 
F(M_1)
\ar[r]^(0.4){e_{F(\phi)}}
& 
\opn{Cone}_{B, \cat{N}}(F(\phi))
\ar[rr]^{p_{F(\phi)}}
& &
\opn{T}_{B, \cat{N}}(F(M_0))
} \]
is an isomorphism of triangles in $\dcat{C}_{\mrm{str}}(B, \cat{N})$. 
\end{cor}

\begin{proof}
Just rearrange the diagram in item (2) of the theorem. 
\end{proof}




\mysubsection{Examples of DG Functors} \label{subsec:exa-DGfun}

Recall that $\cat{M}$ and $\cat{N}$ are abelian categories, and $A$ 
and $B$ are DG rings.
Here are four examples of DG functors, of various types. 
We work out in detail the translation isomorphism, the cone isomorphism and 
the gauge in each example. These examples 
should serve as templates for constructing other DG functors. 

\begin{exa} \label{exa:1117}
Here $A = B = \K$, so 
$\dcat{C}(A, \cat{M}) = \dcat{C}(\cat{M})$
and $\dcat{C}(B, \cat{N}) = \dcat{C}(\cat{N})$. 
Let $F : \cat{M} \to \cat{N}$ be a $\K$-linear functor. It extends to a functor 
$\dcat{C}(F) : \dcat{C}(\cat{M}) \to \dcat{C}(\cat{N})$
as follows: on objects, a complex 
\[ M = \bigl( \{ M^i \}_{i \in \Z}, \{ \d_M^i \}_{i \in \Z} \bigr) 
\in \dcat{C}(\cat{M}) \]
goes to the complex 
\[ \dcat{C}(F)(M) := \bigl( \{ F(M^i) \}, \{ F(\d_M^i) \} \bigr) \in 
\dcat{C}(\cat{N}) . \]
A morphism $\phi = \{ \phi^j \}_{j \in \Z}$ in $\dcat{C}(\cat{M})$ goes 
to the morphism 
$\dcat{C}(F)(\phi) : = \lb \{ F(\phi^j) \}_{j \in \Z}$
in $\dcat{C}(\cat{N})$. A slightly tedious calculation shows that $\dcat{C}(F)$ 
is a graded functor.

Given a complex $M \in \dcat{C}(\cat{M})$, let 
$N := \dcat{C}(F)(M) \in \dcat{C}(\cat{N})$.
Then the translation of $N$ is 
$\opn{T}_{\cat{N}}(N) = \dcat{C}(F) \bigl( \opn{T}_{\cat{M}}(M) \bigr)$;
and the little t operator of $N$ is 
$\opn{t}_{N} = \dcat{C}(F)(\opn{t}_{M})$. 
So the translation isomorphism 
\[ \tau_{\dcat{C}(F)} :  \dcat{C}(F) \circ \opn{T}_{\cat{M}} \iso 
\opn{T}_{\cat{N}} \circ \, \dcat{C}(F) \]
of functors 
$\dcat{C}_{\mrm{str}}(\cat{M}) \to \dcat{C}_{\mrm{str}}(\cat{N})$
is the identity automorphism of this functor. 

Let $\phi : M_0 \to M_1$ be a morphism in 
$\dcat{C}_{\mrm{str}}(\cat{M})$, 
whose image under $\dcat{C}(F)$ is the morphism 
$\psi : N_0 \to N_1$  in $\dcat{C}_{\mrm{str}}(\cat{N})$.
Then 
\[ \opn{Cone}(\psi) = N_1 \oplus \opn{T}_{\cat{N}}(N_0) =
\dcat{C}(F) \bigl( \opn{Cone}(\phi) \bigr)  \]
as graded objects in $\cat{N}$, with differential 
\[ \d_{\opn{Cone}(\psi)} = 
\bmat{\d_{N_1} & \psi \circ \opn{t}_{N_0}^{-1} \\[0.2em] 0 & \d_{\opn{T}(N_0)}} 
= \dcat{C}(F) \Bigl( 
\bmat{\d_{M_1} & \phi \circ \opn{t}_{M_0}^{-1} \\[0.2em] 0 & \d_{\opn{T}(M_0)}}
\Bigr) = \dcat{C}(F) \bigl( \d_{\opn{Cone}(\phi)} \bigr) . \]
We see that the cone isomorphism $\opn{cone}(F, \phi) $ is the identity 
automorphism of the DG module $\opn{Cone}(\psi)$. The gauge $\ga_{\dcat{C}(F)}$ 
of the graded functor $\dcat{C}(F)$ is zero. Therefore, by Theorem 
\ref{thm:1225}, $\dcat{C}(F)$ is a DG functor. 
\end{exa}

The next example is much more complicated, and we work out the full details 
(only once -- later on, such details will be left to the reader). 

\begin{exa} \label{exa:1118}
Let $A$ and $B$  be DG rings, and fix some 
$N \in \cat{DGMod} B \ot A^{\mrm{op}} = \dcat{C}(B \ot A^{\mrm{op}})$.
In other words, $N$ is a DG $B$-$A$-bimodule. 
For every $M \in \dcat{C}(A)$ we have a DG $\K$-module
$F(M) := N \ot_{A} M$,
as in Definition \ref{dfn:1100}. 
The differential of $F(M)$ is
\begin{equation} \label{eqn:1220}
\d_{F(M)} = \d_N \ot \opn{id}_M + \opn{id}_N \ot \, \d_M . 
\end{equation}
See formula (\ref{eqn:4105}) 
regarding the Koszul sign rule that's involved.
But $F(M)$ has the structure of a DG $B$-module: for every $b \in B$, 
$n \in N$ and $m \in M$, the action is 
$b \cd (n \ot m) := (b \cd n) \ot m$.
Clearly 
\[ F : \dcat{C}(A) = \cat{DGMod} A \to \dcat{C}(B) = \cat{DGMod} B \]
is a linear functor. We will show that it is actually a DG functor. 

Let $M_0, M_1 \in \dcat{C}(A)$, and consider the $\K$-linear homomorphism 
\begin{equation} \label{eqn:1120}
F : \opn{Hom}_A(M_0, M_1) \to 
\opn{Hom}_B(N \ot_{A} M_0, N \ot_{A} M_1) .
\end{equation}
Take any $\phi \in \opn{Hom}_A(M_0, M_1)^i$. Then 
\[  F(\phi) \in \opn{Hom}_B(N \ot_{A} M_0, N \ot_{A} M_1) \]
is the homomorphism that on a homogeneous tensor 
$n \ot m \in (N \ot_{A} M_0)^{k + j}$, 
with  $n \in N^k$ and $m \in M_0^j$, has the value
\[ F(\phi) (n \ot m) = (-1)^{i \cd k} \cd n \ot \phi(m) \in 
(N \ot_{A} M_1)^{k + j + i} .  \]
In other words, 
\begin{equation}  \label{eqn:1222}
F(\phi) = \opn{id}_N \ot \, \phi . 
\end{equation}
We see that the homomorphism $F(\phi)$ has degree $i$. So
$F$ is a graded functor. 

Let us calculate $\ga_F$, the gauge of $F$. 
From (\ref{eqn:1222}) and (\ref{eqn:1220}) we get 
$\ga_{F, M} = \d_N \ot \opn{id}_M$,
which is often a nonzero endomorphism of $F(M)$. Still, take any degree $i$ 
morphism $\phi : M_0 \to M_1$ 
in $\dcat{C}(A)$. Then 
\[ \begin{aligned}
& \ga_{M_1} \circ F(\phi) = 
(\d_N \ot \opn{id}_{M_1}) \circ (\opn{id}_N \ot \, \phi)
\\ & \qquad
= \d_N \ot \phi = 
(-1)^i \cd (\opn{id}_N  \ot \, \phi) \circ (\d_N \ot \opn{id}_{M_0})  
= (-1)^i \cd F(\phi) \circ \ga_{M_0} . 
\end{aligned} \]
We see that $\ga_F$ satisfies the condition of Definition \ref{dfn:1195}(1), 
which is really Definition \ref{dfn:1695}.
By Theorem \ref{thm:1225}, $F$ is a DG functor. (It is also possible to 
calculate directly that $F$ is a DG functor.) 

Finally let us figure out what is the translation isomorphism $\tau_F$ of the 
functor $F$. Take $M \in \dcat{C}(A)$. 
Then 
$\tau_{F, M} : F(\opn{T}_{A}(M)) \to \opn{T}_{B}(F(M))$
is an isomorphism in $\dcat{C}_{\mrm{str}}(B)$. 
By Definition \ref{dfn:1150} we have 
$\tau_{F, M} = \opn{t}_{F(M)} \circ \, F(\opn{t}_M)^{-1}$.
Take any $n \in N^k$ and $m \in M^{j + 1}$, so that 
$n \ot \opn{t}_M(m) \in (N \ot_A \opn{T}_{A}(M))^{k + j} 
= F(\opn{T}_{A}(M))^{k + j}$,
a degree $k + j$ element of $F(\opn{T}_{A}(M))$. But 
\[ n \ot \opn{t}_M(m) = (-1)^k \cd (\opn{id}_N \ot \opn{t}_M)(n \ot m) = 
(-1)^k \cd F(\opn{t}_M)(n \ot m) . \]
Therefore 
\[ \tau_{F, M}(n \ot \opn{t}_M(m)) = 
(-1)^k \cd \opn{t}_{F(M)}(n \ot m) \in \opn{T}_{B}(F(M))^{k + j} .  \]
Observe that when $N$ is concentrated in degree $0$, we are back in the 
situation of Example \ref{exa:1117}, in which there are no sign twists, and 
$\tau_{F, M}$ is the identity automorphism. 
\end{exa}

\begin{exa} \label{exa:1120}
Let $A$ and $B$  be DG rings, and fix some 
$N \in \cat{DGMod} A \ot B^{\mrm{op}} = \dcat{C}(B \ot A^{\mrm{op}})$.
For any $M \in \cat{DGMod} A$ we define 
$F(M) := \opn{Hom}_A(N, M)$.
This is a DG $B$-module: for every $b \in B^i$ and
$\phi \in \opn{Hom}_A(N, M)^j$, the homomorphism 
$b \cd \phi \in \opn{Hom}_A(N, M)^{i + j}$ has value
$(b \cd \phi)(n) := (-1)^{i \cd (j + k)} \cd \phi(n \cd b) \in M^{i + j + k}$
on $n \in N^k$. As in the previous example, 
\[ F : \dcat{C}(A) = \cat{DGMod} A \to \dcat{C}(B) = \cat{DGMod} B \]
is a $\K$-linear graded functor. 

The value of the gauge $\ga_F$ at $M \in \dcat{C}(A)$  is 
$\ga_{F, M} = \opn{Hom}(\d_N, \opn{id}_M)$.
See formula (\ref{eqn:4106})  
regarding this notation. Namely for 
$\psi \in F(M)^j = \opn{Hom}_A(N, M)^j$
we have 
$\ga_{F, M}(\psi) = (-1)^{\lms j} \cd \psi \circ \d_N$.
It is not too hard to check that $\ga_F$ is a degree $1$ morphism of graded 
functors. Hence, by Theorem \ref{thm:1225}, $F$ is a DG functor. 

The formula for the translation isomorphism $\tau_F$ is as follows. 
Take $M \in \dcat{C}(A)$. Then 
\[ \tau_{F, M} : F(\opn{T}_{A}(M)) = \opn{Hom}_A(N, \opn{T}_{A}(M))
\to \opn{T}_{B}(F(M)) = \opn{T}_{B}(\opn{Hom}_A(N, M)) \]
is, by definition, 
$\tau_{F, M} = \opn{t}_{F(M)} \circ \, F(\opn{t}_{M})^{-1}$.
Now
$F(\opn{t}_{M})^{-1} = \opn{Hom}(\opn{id}_N, \opn{t}_{M}^{-1})$.
So given any 
$\psi \in F(\opn{T}_{A}(M))^k$, we have 
$\tau_{F, M}(\psi) = \opn{t}_{F(M)}(\opn{t}_{M}^{-1} \circ \, \psi) \in 
\opn{T}_{B}(F(M))^k$.
\end{exa}

We end with a contravariant example. 

\begin{exa} \label{exa:2300}
Let $A$ be a commutative ring. Fix some complex $N \in \dcat{C}(A)$. 
For any $M \in \dcat{C}(A)$ let 
$F(M) := \opn{Hom}_A(M, N) \in \dcat{C}(A)$.
For every degree $i$ morphism 
$\phi : M_0 \to M_1$ in $\dcat{C}(A)$ let
$F(\phi) : F(M_1) \to F(M_0)$ be the degree $i$ morphism
$F(\phi) := \opn{Hom}_A(\phi, \opn{id}_N)$.

A direct calculation, that we leave to the reader, shows that 
\[ F : \opn{Hom}_{\dcat{C}(A)}(M_0, M_1) \to 
\opn{Hom}_{\dcat{C}(A)} \bigl( F(M_1), F(M_0) \bigr) \]
is a strict homomorphism of DG $\K$-modules. It also satisfies conditions (a) 
and (b) of Definition \ref{dfn:2495}. Thus we have a contravariant DG 
functor 
$F : \dcat{C}(A) \to \dcat{C}(A)$.
For contravariant DG functors we do not talk about translation isomorphisms or 
gauges. 
\end{exa}

We will return to the {\em DG bifunctors} $\opn{Hom}(-, -)$ and 
$(- \ot -)$ later in the book, in Subsection \ref{subsec:dg-bifun}.

